\chapter{Data Structures}
\label{ch:data_structures}

This chapter provides an exhaustive, examination-oriented foundation in data structures, covering linear and non-linear abstractions, algorithmic complexity, implementation nuances in C and C++, and high-frequency problem sets required for Bangladesh Government IT Recruitment Examinations (including Assistant Programmer, Assistant Director (ICT), Assistant Network Engineer, Assistant Database Administrator, and Maintenance Engineer posts).

The chapter is systematically organized into six interconnected core areas:
\begin{enumerate}[label=\bfseries Part \Roman*:]
    \item \textbf{Stack and Queue Architectures:} LIFO and FIFO operations, queue implementation using two stacks, amortized analysis, real-world systems applications, and common MCQ traps.
    \item \textbf{Linked Lists and Binary Heaps:} Singly, doubly, and circular linked lists, array vs. pointer trade-offs, max-heap and min-heap construction (heapify), insertion, deletion, and Heap Sort.
    \item \textbf{Circular Queues and Priority Queues:} Eliminating linear memory exhaustion via modulo arithmetic indexing, priority queue heap mappings, and scheduling algorithms.
    \item \textbf{Trees, Binary Trees, and B+ Trees:} Tree hierarchies, binary search trees (BST), balanced AVL trees, and exact mathematical formulations for minimum/maximum keys in B+ trees.
    \item \textbf{Graphs, Adjacency Representations, and Hash Tables:} Graph theory, adjacency list vs. matrix space/time trade-offs, hash functions, and collision resolution (chaining vs. open addressing).
    \item \textbf{Data Structure Recruitment Examination Problem Bank:} Essential worked algorithm-tracing problems, programming solutions, and analytical viva questions with complete step-by-step solutions.
\end{enumerate}


\begin{center}
\vspace{1em}
{\color{primary}\rule{0.5\textwidth}{1.5pt}}\\[6pt]
{\Large\bfseries\color{primary} Part I: Stack and Queue Architectures}\\[3pt]
{\color{primary}\rule{0.5\textwidth}{1.5pt}}
\vspace{1em}
\end{center}

This chapter covers the fundamental differences between stacks and queues, followed by a detailed explanation and implementation of a queue using two stacks. The emphasis is on conceptual understanding, algorithms, complexity analysis, practical applications, and common examination questions.

\section{Stack and Queue: Fundamental Concepts}

\subsection{What is a Data Structure?}

A \textbf{data structure} is a way of organizing and storing data so that it can be accessed, modified, and processed efficiently. Examples include arrays, linked lists, stacks, queues, trees, graphs, heaps, and hash tables.

A stack and a queue are both \textbf{linear data structures} because their elements are logically arranged in a sequence. However, they differ in the order in which elements are inserted and removed.

\subsection{What is a Stack?}

A \textbf{stack} is a linear data structure that follows the \textbf{LIFO (Last In, First Out)} principle. The last element inserted into the stack is the first element removed.

A real-life example is a stack of plates in a cafeteria. When a new plate is placed on top, it becomes the first plate available to be removed.

Suppose we insert the following elements in order:
\[
10, \; 20, \; 30
\]
The top element is $30$. Removing elements produces this sequence:
\[
30 \longrightarrow 20 \longrightarrow 10
\]
The most recently inserted element is removed first.

\subsubsection{Basic Stack Operations}

\begin{table}[htbp]
\centering
\small
\begin{tabularx}{\textwidth}{l p{7.5cm} X}
\toprule
\textbf{Operation} & \textbf{Meaning} & \textbf{Typical Time Complexity} \\
\midrule
\texttt{push(x)}           & Insert element $x$ at the top                & $O(1)$ \\
\texttt{pop()}             & Remove the top element                       & $O(1)$ \\
\texttt{top()} or \texttt{peek()} & Read the top element without removing it     & $O(1)$ \\
\texttt{isEmpty()}         & Check whether the stack contains no elements & $O(1)$ \\
\texttt{size()}            & Return the number of elements                & $O(1)$ when maintained \\
\bottomrule
\end{tabularx}
\caption{Basic stack operations and standard complexities.}
\end{table}

These complexities assume a standard stack implementation and exclude dynamic-array reallocation during an individual insertion.

\subsubsection{Important Conditions}
\begin{itemize}
    \item \textbf{Stack overflow:} An insertion is attempted when a fixed-capacity stack is completely full.
    \item \textbf{Stack underflow:} A removal or top-element access is attempted when the stack is empty.
    \item \textbf{Top:} The pointer or index indicating the position at which stack insertion and removal take place.
\end{itemize}

In a dynamically allocated stack (such as a linked list implementation or \texttt{std::vector}), capacity can grow until system memory is exhausted, so fixed-capacity overflow does not apply.

\subsubsection{Applications of Stacks}
\begin{enumerate}
    \item \textbf{Function Calls and Recursion:} Programming language runtimes use a call stack to manage active function calls. Each call creates a stack frame storing local variables, parameters, and the return address.
    \item \textbf{Undo Operations:} Software applications (such as text editors and graphic tools) store user action history in an undo stack. The most recent action is always the first one to be undone.
    \item \textbf{Expression Evaluation and Conversion:} Stacks are indispensable for evaluating postfix (Reverse Polish) expressions and converting infix expressions to postfix or prefix notations (e.g., Dijkstra's Shunting-yard algorithm). For example, the postfix expression \texttt{2 3 +} evaluates to $5$.
    \item \textbf{Parentheses and Delimiter Matching:} Compilers and syntax checkers use stacks to verify that opening and closing brackets (\texttt{()}, \texttt{[]}, \texttt{\{\}}) are properly matched and nested.
    \item \textbf{Depth-First Search (DFS) and Backtracking:} Stacks facilitate graph traversal, maze exploration, game-tree evaluation, and constraint-satisfaction problems.
\end{enumerate}

\subsection{What is a Queue?}

A \textbf{queue} is a linear data structure that follows the \textbf{FIFO (First In, First Out)} principle. The first element inserted into the queue is the first element removed.

A real-life example is a line of customers waiting at a service counter. The customer who arrived first is attended to first.

Suppose we insert in order:
\[
10, \; 20, \; 30
\]
The front element is $10$, and the rear element is $30$. Removing elements produces this sequence:
\[
10 \longrightarrow 20 \longrightarrow 30
\]
The earliest inserted element is removed first.

\subsubsection{Basic Queue Operations}

\begin{table}[htbp]
\centering
\small
\begin{tabularx}{\textwidth}{l p{7.5cm} X}
\toprule
\textbf{Operation} & \textbf{Meaning} & \textbf{Typical Time Complexity} \\
\midrule
\texttt{enqueue(x)}        & Insert element $x$ at the rear             & $O(1)$ \\
\texttt{dequeue()}         & Remove the front element                   & $O(1)$ \\
\texttt{front()} or \texttt{peek()} & Read the front element without removing it & $O(1)$ \\
\texttt{isEmpty()}         & Check whether the queue is empty           & $O(1)$ \\
\texttt{size()}            & Return the number of elements              & $O(1)$ when maintained \\
\bottomrule
\end{tabularx}
\caption{Basic queue operations and standard complexities.}
\end{table}

These operation complexities assume a suitable queue implementation, such as a circular array queue or a linked list queue maintaining both head (front) and tail (rear) pointers. A naive linear array implementation that shifts elements upon removal incurs an undesirable $O(n)$ time for \texttt{dequeue()}.

\subsubsection{Important Queue Terminology}
\begin{itemize}
    \item \textbf{Front (Head):} The end from which elements are inspected and removed.
    \item \textbf{Rear (Tail):} The end at which new elements are inserted.
    \item \textbf{Enqueue:} The operation of inserting an element at the rear.
    \item \textbf{Dequeue:} The operation of removing an element from the front.
    \item \textbf{Queue underflow:} Attempting to dequeue or peek when the queue contains zero elements.
    \item \textbf{Queue overflow:} Attempting to enqueue into a fixed-capacity queue that is already full.
\end{itemize}

\subsubsection{Applications of Queues}
\begin{enumerate}
    \item \textbf{CPU and Process Scheduling:} Operating systems maintain queues (ready queues, I/O queues) to schedule processes awaiting execution (e.g., Round-Robin scheduling).
    \item \textbf{Printer Spooling:} Print jobs sent from multiple applications wait in a FIFO print buffer queue until the printer finishes previous jobs.
    \item \textbf{Network Packet Buffering:} Switches, routers, and network interfaces buffer incoming and outgoing packets in queues.
    \item \textbf{Breadth-First Search (BFS):} Graph and tree BFS algorithms explore vertices level-by-level using a FIFO queue.
    \item \textbf{Asynchronous Request Processing:} Web servers, message brokers (Kafka, RabbitMQ), and worker pools buffer tasks in queues to decouple production from consumption.
    \item \textbf{Customer Service and Ticketing:} Ticket reservation systems, call centers, and banking counters order requests strictly by arrival time.
\end{enumerate}

\section{Differences Between Stack and Queue}

The comparison between stacks and queues is a frequent subject in technical interviews, written recruitment exams, and viva voce.

\begin{table}[htbp]
\centering
\footnotesize
\begin{tabularx}{\textwidth}{l X X}
\toprule
\textbf{Basis} & \textbf{Stack} & \textbf{Queue} \\
\midrule
Operating Principle             & LIFO (Last In, First Out)               & FIFO (First In, First Out) \\
Full Form                       & Last In, First Out                      & First In, First Out \\
Insertion Operation             & Push                                    & Enqueue \\
Removal Operation               & Pop                                     & Dequeue \\
Insertion Position              & Top                                     & Rear \\
Removal Position                & Top                                     & Front \\
Main Access Points              & Top only                                & Front and Rear \\
Element Removed First           & Most recently inserted element          & Earliest inserted element \\
Typical Implementation          & Dynamic array or linked list            & Circular array or linked list \\
Common Applications             & Recursion, undo, parsing, backtracking  & Scheduling, buffering, BFS, spooling \\
Empty-State Error               & Stack underflow                         & Queue underflow \\
Fixed-Capacity Full Error       & Stack overflow                          & Queue overflow \\
Insertion Complexity            & $O(1)$                                  & $O(1)$ \\
Removal Complexity              & $O(1)$                                  & $O(1)$ \\
\bottomrule
\end{tabularx}
\caption{Comprehensive comparison between Stack and Queue data structures.}
\end{table}

Both data structures represent abstract data types (ADTs) and can be realized using contiguous arrays or node-based linked lists. Neither is restricted to a single underlying memory layout.

\subsection{The Key Difference: Order of Removal}

Consider an identical sequence of insertions into both structures:
\begin{enumerate}
    \item Insert $A$
    \item Insert $B$
    \item Insert $C$
    \item Remove one element
\end{enumerate}

\begin{table}[htbp]
\centering
\small
\begin{tabularx}{0.85\textwidth}{l X X}
\toprule
\textbf{Operation} & \textbf{Stack (bottom to top)} & \textbf{Queue (front to rear)} \\
\midrule
Insert $A$ & $A$       & $A$ \\
Insert $B$ & $A, B$    & $A, B$ \\
Insert $C$ & $A, B, C$ & $A, B, C$ \\
Remove one & Removes $C$ (leaving $A, B$) & Removes $A$ (leaving $B, C$) \\
\bottomrule
\end{tabularx}
\caption{Visualizing the distinct removal sequences for identical insertion orders.}
\end{table}

\begin{quote}
\textbf{Core Exam Takeaway:} A stack and a queue may store identical elements, yet their processing semantics diverge completely: a stack inverts arrival chronology, whereas a queue strictly honors it.
\end{quote}

\subsection{Can a Stack or Queue Access an Arbitrary Element Directly?}

Neither standard stack nor standard queue supports direct random access ($O(1)$ indexed lookup). 
\begin{itemize}
    \item A stack restricts visibility and modification exclusively to its \texttt{top}.
    \item A queue restricts insertion to its \texttt{rear} and deletion/reading to its \texttt{front}.
\end{itemize}
To retrieve the $k$-th element, all preceding elements must be sequentially popped or dequeued, necessitating auxiliary storage if their original values must be preserved. Although an underlying array possesses indexing capability, exposing arbitrary access breaks the ADT abstraction boundary.

\subsection{Stack versus Queue: Common Exam Traps}
\begin{enumerate}
    \item \textbf{Principle Mismatch:} Remember: LIFO belongs to stacks; FIFO belongs to queues.
    \item \textbf{End Symmetry:} In a stack, insertion and removal happen at the \emph{same end} (the top). In a queue, insertion and removal happen at \emph{opposite ends} (rear and front).
    \item \textbf{Priority Queues are Not Pure FIFO:} In a priority queue, elements are extracted according to priority keys rather than arrival chronology.
    \item \textbf{Deques (Double-Ended Queues):} A deque generalizes both stacks and queues by permitting insertions and deletions at both ends in $O(1)$ time.
    \item \textbf{Array Queue Naivety:} A queue built on a standard array without circular wrapping either wastes memory or requires $O(n)$ shifting on dequeue.
    \item \textbf{Mutual Emulation:} A stack can be implemented using two queues, and a queue can be implemented using two stacks.
\end{enumerate}

\section{Implementing a Queue Using Two Stacks}

This classic data structure problem appears frequently in recruitment tests (e.g., BPDB, Biman, BCS, BAPEX, Bangladesh Bank IT).

\subsection{The Problem}
A stack is strictly LIFO; a queue must behave as FIFO. The task is to construct a fully functioning queue supporting:
\begin{itemize}
    \item \texttt{enqueue(x)}: Insert an item into the queue.
    \item \texttt{dequeue()}: Remove the oldest item from the queue.
    \item \texttt{front() / peek()}: Inspect the oldest item without removal.
    \item \texttt{isEmpty()}: Determine whether the queue is empty.
\end{itemize}
using \textbf{only two standard stack instances}, conventionally denoted $S_1$ and $S_2$, with standard stack primitives (\texttt{push}, \texttt{pop}, \texttt{top}, \texttt{empty}).

\subsection{Why Are Two Stacks Needed?}

Suppose we push the values $10, 20, 30$ onto a single stack $S_1$:
\[
\text{Stack } S_1 \text{ (bottom to top): } [10, 20, 30] \quad \Longrightarrow \quad \text{Top} = 30
\]
Popping from $S_1$ yields $30, 20, 10$, which violates FIFO order.

Now, pop each element from $S_1$ and push it onto a second stack $S_2$:
\begin{enumerate}
    \item Pop $30$ from $S_1$, push to $S_2$.
    \item Pop $20$ from $S_1$, push to $S_2$.
    \item Pop $10$ from $S_1$, push to $S_2$.
\end{enumerate}
Stack $S_2$ now contains:
\[
\text{Stack } S_2 \text{ (bottom to top): } [30, 20, 10] \quad \Longrightarrow \quad \text{Top} = 10
\]
When we pop from $S_2$, the first element extracted is $10$---the very first element that was originally entered!

\begin{quote}
\textbf{The Inversion Principle:} Transferring elements from one LIFO stack to another reverses their ordering:
\[
\text{LIFO} \times \text{LIFO} = \text{FIFO}
\]
Two reversals restore the original arrival sequence.
\end{quote}

\section{Method 1: Costly Enqueue, Efficient Dequeue}

In this design, every \texttt{enqueue} proactively rearranges elements so that the oldest element perpetually resides at the top of $S_1$.

\subsection{Algorithm}

\subsubsection{Enqueue Operation (\texttt{ENQUEUE(x)})}
To insert element $x$:
\begin{enumerate}
    \item While $S_1$ is not empty, pop all elements from $S_1$ and push them into $S_2$.
    \item Push the new element $x$ into $S_1$.
    \item While $S_2$ is not empty, pop all elements from $S_2$ and push them back into $S_1$.
\end{enumerate}
After this two-way transfer, $x$ is placed at the very bottom of $S_1$, and the oldest element is restored to the top.

\subsubsection{Dequeue Operation (\texttt{DEQUEUE()})}
\begin{enumerate}
    \item If $S_1$ is empty, report queue underflow.
    \item Otherwise, pop and return the top element of $S_1$.
\end{enumerate}

\subsubsection{Front Operation (\texttt{FRONT()})}
\begin{enumerate}
    \item If $S_1$ is empty, report queue underflow.
    \item Otherwise, return the top element of $S_1$ without popping.
\end{enumerate}

\subsection{Pseudocode}

\begin{lstlisting}[style=pseudocode]
ENQUEUE(x):
    while S1 is not empty:
        S2.push(S1.pop())
    
    S1.push(x)
    
    while S2 is not empty:
        S1.push(S2.pop())

DEQUEUE():
    if S1.isEmpty():
        error "Queue Underflow"
    return S1.pop()

FRONT():
    if S1.isEmpty():
        error "Queue Underflow"
    return S1.top()
\end{lstlisting}

\subsection{Step-by-Step Example}
Insert $10$, then $20$, then $30$:
\begin{itemize}
    \item \textbf{Enqueue 10:} $S_1$ is empty. Push $10$ into $S_1$. $S_1 = [10]$ (top is $10$).
    \item \textbf{Enqueue 20:} Move $10$ to $S_2$. Push $20$ to $S_1$. Move $10$ back to $S_1$. $S_1 = [20, 10]$ (top is $10$).
    \item \textbf{Enqueue 30:} Move $10, 20$ to $S_2$. Push $30$ to $S_1$. Move $20, 10$ back to $S_1$. $S_1 = [30, 20, 10]$ (top is $10$).
\end{itemize}
Calling \texttt{dequeue()} immediately pops $10$, then $20$, then $30$, exactly preserving FIFO order.

\subsection{Complexity Analysis}
Let $n$ denote the number of elements in the queue:
\begin{table}[htbp]
\centering
\small
\begin{tabularx}{0.7\textwidth}{l X}
\toprule
\textbf{Operation} & \textbf{Time Complexity} \\
\midrule
Enqueue      & $O(n)$ \\
Dequeue      & $O(1)$ \\
Front / Peek & $O(1)$ \\
isEmpty      & $O(1)$ \\
Space        & $O(n)$ total storage \\
\bottomrule
\end{tabularx}
\caption{Complexity profile for Method 1.}
\end{table}

\textbf{Trade-off:} Dequeue is instantaneous ($O(1)$), but every single enqueue costs $2n$ stack operations ($O(n)$). This approach is only advantageous when reads/removals vastly outnumber insertions.

\section{Method 2: Efficient Enqueue, Costly Dequeue (Recommended)}

This is the standard and widely preferred method because it achieves \textbf{$O(1)$ amortized time complexity} across both insertions and removals.

Here, the two stacks maintain distinct functional responsibilities:
\begin{itemize}
    \item $S_1$ (\textbf{Input Stack}): Receives all newly enqueued items.
    \item $S_2$ (\textbf{Output Stack}): Dispenses items for dequeue and peek operations.
\end{itemize}

\begin{quote}
\textbf{The Golden Transfer Rule:} Elements are transferred from $S_1$ to $S_2$ \textbf{if and only if} $S_2$ is completely empty. If $S_2$ already contains elements, never transfer from $S_1$, as doing so would violate FIFO chronology.
\end{quote}

\subsection{Algorithm}

\subsubsection{Enqueue Operation (\texttt{ENQUEUE(x)})}
\begin{enumerate}
    \item Push $x$ directly onto stack $S_1$.
    \item Do not move any other elements.
\end{enumerate}
Time complexity: $O(1)$.

\subsubsection{Dequeue Operation (\texttt{DEQUEUE()})}
\begin{enumerate}
    \item If $S_2$ is not empty, pop and return the top element of $S_2$.
    \item If $S_2$ is empty:
    \begin{itemize}
        \item While $S_1$ is not empty, pop the top of $S_1$ and push it onto $S_2$.
    \end{itemize}
    \item If $S_2$ is still empty (meaning both $S_1$ and $S_2$ were empty), report queue underflow.
    \item Otherwise, pop and return the top of $S_2$.
\end{enumerate}

\subsubsection{Front Operation (\texttt{FRONT()})}
Same logic as \texttt{dequeue()}, but inspect $S_2.\texttt{top()}$ without popping.

\subsubsection{Pseudocode}

\begin{lstlisting}[style=pseudocode]
ENQUEUE(x):
    S1.push(x)

DEQUEUE():
    if S2.isEmpty():
        while not S1.isEmpty():
            S2.push(S1.pop())
    
    if S2.isEmpty():
        error "Queue Underflow"
    
    return S2.pop()

FRONT():
    if S2.isEmpty():
        while not S1.isEmpty():
            S2.push(S1.pop())
            
    if S2.isEmpty():
        error "Queue Underflow"
        
    return S2.top()

IS_EMPTY():
    return S1.isEmpty() and S2.isEmpty()
\end{lstlisting}

\subsection{Step-by-Step Execution Trace}

Consider the execution sequence:
\texttt{enqueue(10)}, \texttt{enqueue(20)}, \texttt{enqueue(30)}, \texttt{dequeue()}, \texttt{enqueue(40)}, \texttt{dequeue()}, \texttt{dequeue()}, \texttt{dequeue()}, \texttt{dequeue()}.

\begin{table}[htbp]
\centering
\small
\begin{tabularx}{\textwidth}{c l X X l}
\toprule
\textbf{Step} & \textbf{Operation} & \textbf{Stack $S_1$ (bottom $\to$ top)} & \textbf{Stack $S_2$ (bottom $\to$ top)} & \textbf{Result} \\
\midrule
1 & \texttt{enqueue(10)} & $[10]$         & Empty       & --- \\
2 & \texttt{enqueue(20)} & $[10, 20]$     & Empty       & --- \\
3 & \texttt{enqueue(30)} & $[10, 20, 30]$ & Empty       & --- \\
4 & \texttt{dequeue()}   & Empty          & $[30, 20]$  & Returns $10$ \\
5 & \texttt{enqueue(40)} & $[40]$         & $[30, 20]$  & --- \\
6 & \texttt{dequeue()}   & $[40]$         & $[30]$      & Returns $20$ \\
7 & \texttt{dequeue()}   & $[40]$         & Empty       & Returns $30$ \\
8 & \texttt{dequeue()}   & Empty          & Empty       & Returns $40$ (transferred) \\
9 & \texttt{dequeue()}   & Empty          & Empty       & Underflow (Error) \\
\bottomrule
\end{tabularx}
\caption{Detailed step-by-step state trace of the two stacks.}
\end{table}

\textbf{Key Observation at Step 5:} When $40$ is enqueued, it is pushed onto $S_1$. It does not displace the elements already waiting in $S_2$ ($20$ and $30$). Those older elements are correctly dequeued first. Only after $S_2$ is exhausted does $40$ transfer to $S_2$. The resulting departure order is $10 \to 20 \to 30 \to 40$.

\subsection{Complete C++ Implementation}

The following standard C++ implementation encapsulates the queue, properly handles underflow via return booleans and output references, and exposes complete queue functionality.

\begin{lstlisting}[style=cppstyle]
#include <iostream>
#include <stack>
using namespace std;

class QueueUsingTwoStacks {
private:
    stack<int> S1;  // Input stack for enqueued elements
    stack<int> S2;  // Output stack for dequeued elements

    // Helper: transfer elements from S1 to S2 only when S2 is empty
    void moveIfNeeded() {
        if (S2.empty()) {
            while (!S1.empty()) {
                S2.push(S1.top());
                S1.pop();
            }
        }
    }

public:
    // Insert an element at the rear of the queue: O(1)
    void enqueue(int value) {
        S1.push(value);
    }

    // Remove the front element: O(1) amortized
    // Returns false if queue underflow occurs
    bool dequeue(int& value) {
        moveIfNeeded();

        if (S2.empty()) {
            return false; // Queue is empty
        }

        value = S2.top();
        S2.pop();
        return true;
    }

    // Read the front element without removing it: O(1) amortized
    bool front(int& value) {
        moveIfNeeded();

        if (S2.empty()) {
            return false;
        }

        value = S2.top();
        return true;
    }

    // Check whether the queue is empty: O(1)
    bool isEmpty() const {
        return S1.empty() && S2.empty();
    }

    // Return the total count of queued elements: O(1)
    size_t size() const {
        return S1.size() + S2.size();
    }
};

int main() {
    QueueUsingTwoStacks q;
    int value;

    q.enqueue(10);
    q.enqueue(20);
    q.enqueue(30);

    if (q.front(value)) {
        cout << "Front: " << value << '\n';
    }

    if (q.dequeue(value)) {
        cout << "Removed: " << value << '\n';
    }

    q.enqueue(40);

    while (q.dequeue(value)) {
        cout << "Removed: " << value << '\n';
    }

    cout << "Queue empty: " << (q.isEmpty() ? "Yes" : "No") << '\n';

    return 0;
}
\end{lstlisting}

\subsubsection{Expected Output}
\begin{lstlisting}[style=pseudocode]
Front: 10
Removed: 10
Removed: 20
Removed: 30
Removed: 40
Queue empty: Yes
\end{lstlisting}

\subsubsection{Explanation of Implementation Details}
\begin{itemize}
    \item \texttt{moveIfNeeded()}: Centralizes the transfer logic. By enforcing \texttt{if (S2.empty())}, elements are only transferred when the output pipeline is exhausted, preserving strict FIFO ordering.
    \item \texttt{dequeue(int\& value)}: Uses pass-by-reference to safely return the value while providing a boolean return code to distinguish successful extraction from queue underflow.
    \item \texttt{isEmpty()}: Returns \texttt{true} only when \emph{both} $S_1$ and $S_2$ are empty. A queue is not empty simply because $S_1$ is empty (as unserved elements may still exist in $S_2$).
\end{itemize}

\section{Time and Space Complexity of the Two-Stack Queue}

Recruitment exams frequently test the distinction between \textbf{worst-case} and \textbf{amortized} complexity for Method 2.

\subsection{Complexity Comparison}

\begin{table}[htbp]
\centering
\small
\begin{tabularx}{\textwidth}{l X X}
\toprule
\textbf{Operation} & \textbf{Method 1 (Costly Enqueue)} & \textbf{Method 2 (Efficient Enqueue)} \\
\midrule
Enqueue         & $O(n)$     & $O(1)$ worst-case \\
Dequeue         & $O(1)$     & $O(1)$ amortized; $O(n)$ worst-case \\
Front / Peek    & $O(1)$     & $O(1)$ amortized; $O(n)$ worst-case \\
isEmpty         & $O(1)$     & $O(1)$ \\
Auxiliary Space & $O(n)$     & $O(n)$ \\
\bottomrule
\end{tabularx}
\caption{Detailed performance comparison between the two approaches.}
\end{table}

\subsection{Why is the Amortized Complexity $O(1)$?}

To establish that Method 2 requires $O(1)$ amortized time per operation, track the full lifecycle of an arbitrary element $e$:
\begin{enumerate}
    \item \textbf{Push onto $S_1$:} Executed once when $e$ is enqueued ($1$ stack push).
    \item \textbf{Pop from $S_1$:} Executed at most once during a batch transfer to $S_2$ ($1$ stack pop).
    \item \textbf{Push onto $S_2$:} Executed at most once during the same batch transfer ($1$ stack push).
    \item \textbf{Pop from $S_2$:} Executed once when $e$ is dequeued ($1$ stack pop).
\end{enumerate}

Notice that element $e$ is pushed onto $S_1$ exactly once, transferred from $S_1$ to $S_2$ at most once, and popped from $S_2$ exactly once. It is \emph{never} transferred back and forth repeatedly.

Thus, over the entire lifetime of any element, it incurs at most \textbf{4 primitive stack operations}.

For an arbitrary sequence of $m$ queue operations involving $n$ elements:
\[
\text{Total Operations} \leq 4n \in O(m)
\]
Averaging the total cost over the $m$ operations yields:
\[
\text{Amortized Cost per Operation} = \frac{O(m)}{m} = O(1)
\]

\begin{quote}
\textbf{Essential Distinction for Viva:}
\begin{itemize}
    \item \textbf{Worst-case complexity:} The cost of the single most expensive operation. For Method 2 \texttt{dequeue()}, if $S_2$ is empty and $S_1$ holds $n$ elements, the single call takes $O(n)$ time.
    \item \textbf{Amortized complexity:} The guaranteed upper bound on average time per operation over any sequence of operations. For Method 2, \texttt{dequeue()} is guaranteed $O(1)$ amortized.
\end{itemize}
\end{quote}

\section{Comparing Both Two-Stack Approaches}

\begin{table}[htbp]
\centering
\small
\begin{tabularx}{\textwidth}{l X X}
\toprule
\textbf{Feature} & \textbf{Costly Enqueue (Method 1)} & \textbf{Costly Dequeue (Method 2)} \\
\midrule
Front Element Location & Top of $S_1$ at all times & Top of $S_2$ (when $S_2$ nonempty) \\
Enqueue Strategy       & Rebalance both stacks ($O(n)$) & Direct push to $S_1$ ($O(1)$) \\
Dequeue Strategy       & Direct pop from $S_1$ ($O(1)$) & Transfer when $S_2$ empty, then pop \\
Primary Advantage      & Predictable $O(1)$ removal & Fast insertions and $O(1)$ amortized dequeues \\
Primary Drawback       & Every enqueue is slow & Occasional batch transfer on dequeue \\
Recommended Use-Case   & Real-time systems prioritizing predictable latency on read & General-purpose workloads (standard solution) \\
\bottomrule
\end{tabularx}
\caption{Strategic architectural comparison between Method 1 and Method 2.}
\end{table}

\section{Common Mistakes and Edge Cases}

These edge cases are frequent traps in coding tests and interview discussions:

\subsection{Mistake 1: Transferring Elements on Every Dequeue}
In Method 2, transferring from $S_1$ to $S_2$ on \emph{every} dequeue destroys FIFO order if $S_2$ already holds elements. If $S_2$ holds older elements $[10, 20]$ and $S_1$ holds $30$, pouring $S_1$ into $S_2$ would place the newer element $30$ on top of $10$, which violates queue ordering. 
\begin{center}
\textbf{Rule: Only transfer from $S_1$ to $S_2$ when $S_2$ is completely empty.}
\end{center}

\subsection{Mistake 2: Checking Only One Stack for Emptiness}
A common bug is writing \texttt{isEmpty() \{ return S1.empty(); \}}. After several enqueues and one dequeue, $S_1$ may be empty while $S_2$ contains multiple queued items. Conversely, after several enqueues without a dequeue, $S_2$ is empty while $S_1$ is full. 
\begin{center}
\textbf{Rule: The queue is empty if and only if BOTH $S_1$ and $S_2$ are empty.}
\end{center}

\subsection{Mistake 3: Popping an Empty Stack}
In C++, invoking \texttt{S.top()} or \texttt{S.pop()} on an empty \texttt{std::stack} causes undefined behavior (often a segmentation fault). Always guard access with an emptiness check.

\subsection{Mistake 4: Equating Stack Top to Queue Front}
In Method 2, the top of $S_1$ contains the \emph{newest} element, whereas the front of the queue is the \emph{oldest} element. The front is at the top of $S_2$ (or at the bottom of $S_1$ if no transfer has occurred).

\subsection{Mistake 5: Stating Dequeue is $O(1)$ Worst-Case}
Never claim that Method 2 \texttt{dequeue()} is $O(1)$ worst-case. In written answers, state: \emph{Enqueue is $O(1)$ worst-case; Dequeue is $O(1)$ amortized and $O(n)$ worst-case.}

\section{Likely Recruitment Examination Questions: Stack and Queue}

\begin{quote}
\textbf{Q1: What is the primary difference between a stack and a queue?}\\
\textbf{Answer:} A stack operates on the LIFO principle (Last In, First Out), where insertions and deletions occur at the same end (the top). A queue operates on the FIFO principle (First In, First Out), where insertions occur at the rear and deletions occur at the front.
\end{quote}

\begin{quote}
\textbf{Q2: How can a FIFO queue be implemented using two LIFO stacks?}\\
\textbf{Answer:} Use an input stack ($S_1$) and an output stack ($S_2$). Push new elements directly onto $S_1$. For dequeue or peek operations, if $S_2$ is empty, transfer all elements from $S_1$ to $S_2$ by repeated popping and pushing, which reverses their order. Then pop from $S_2$. If $S_2$ is not empty, pop directly from $S_2$.
\end{quote}

\begin{quote}
\textbf{Q3: What are the time complexities of enqueue and dequeue in the efficient-enqueue method?}\\
\textbf{Answer:} Enqueue takes $O(1)$ worst-case time. Dequeue takes $O(1)$ amortized time, although an individual dequeue takes $O(n)$ worst-case time when a batch transfer from $S_1$ to $S_2$ is triggered.
\end{quote}

\begin{quote}
\textbf{Q4: Why must elements be transferred from $S_1$ to $S_2$ only when $S_2$ is empty?}\\
\textbf{Answer:} Elements already inside $S_2$ arrived earlier than any elements currently waiting in $S_1$. Pouring elements from $S_1$ while $S_2$ is nonempty would place newer elements on top of older elements, violating FIFO order.
\end{quote}

\begin{quote}
\textbf{Q5: What is the auxiliary space complexity of implementing a queue using two stacks?}\\
\textbf{Answer:} $O(n)$, where $n$ is the total number of elements stored in the queue, because $S_1$ and $S_2$ collectively hold exactly $n$ items.
\end{quote}

\begin{quote}
\textbf{Q6: What happens when dequeue is invoked on an empty queue?}\\
\textbf{Answer:} The operation triggers queue underflow. In a robust implementation, it should be intercepted and handled gracefully (e.g., throwing an underflow exception or returning a boolean status code) without attempting to read from empty internal stacks.
\end{quote}

\begin{quote}
\textbf{Q7: What is amortized analysis, and how does it justify the $O(1)$ dequeue cost?}\\
\textbf{Answer:} Amortized analysis evaluates the average running time per operation over a worst-case sequence of operations. Because every enqueued element is transferred from $S_1$ to $S_2$ at most once in its entire lifetime, $n$ insertions and removals incur at most $4n$ primitive stack operations, guaranteeing an average $O(1)$ cost per operation.
\end{quote}

\begin{quote}
\textbf{Q8: Can a stack be implemented using two queues? If so, what is the complexity?}\\
\textbf{Answer:} Yes. In a costly-push approach, each push transfers elements between the two queues so that the newest element remains at the front ($O(n)$ push, $O(1)$ pop). Alternatively, in a costly-pop approach, push is $O(1)$ and pop shifts $n-1$ elements to the second queue to extract the last inserted element ($O(n)$ pop).
\end{quote}

\section{Quick Revision Summary}

\begin{table}[htbp]
\centering
\small
\begin{tabularx}{\textwidth}{l X}
\toprule
\textbf{Topic} & \textbf{Key Fact} \\
\midrule
Stack Principle                & LIFO (Last In, First Out) \\
Queue Principle                & FIFO (First In, First Out) \\
Stack Operations               & Push and Pop both take place at the Top ($O(1)$) \\
Queue Operations               & Enqueue at Rear, Dequeue at Front ($O(1)$) \\
Boundary Errors                & Underflow (empty access) and Overflow (full insertion) \\
Queue via Two Stacks Core Idea & Transferring elements between two stacks reverses their order \\
Efficient Enqueue Architecture & Enqueue pushes to $S_1$; Dequeue pops from $S_2$ \\
Transfer Condition             & Transfer $S_1 \to S_2$ \textbf{only} when $S_2$ is empty \\
Enqueue Complexity             & $O(1)$ worst-case \\
Dequeue Complexity             & $O(1)$ amortized; $O(n)$ worst-case \\
Total Space Complexity         & $O(n)$ \\
\bottomrule
\end{tabularx}
\caption{Quick revision reference for Stack and Queue fundamentals.}
\end{table}

\begin{quote}
\textbf{Final Golden Rule:} Two LIFO stacks simulate a FIFO queue because reversing a reversed sequence restores original arrival order. Method 2 guarantees $O(1)$ amortized operation times because every element is transferred at most once during its lifetime.
\end{quote}


\begin{center}
\vspace{1em}
{\color{primary}\rule{0.5\textwidth}{1.5pt}}\\[6pt]
{\Large\bfseries\color{primary} Part II: Linked Lists and Binary Heaps}\\[3pt]
{\color{primary}\rule{0.5\textwidth}{1.5pt}}
\vspace{1em}
\end{center}

This chapter covers three major linked-list structures---singly linked lists, doubly linked lists, and circular linked lists---followed by heaps, with particular emphasis on step-by-step max-heap construction, insertion, deletion, and complexity analysis.

The explanations and implementations are aimed at Bangladesh government IT recruitment examinations, including Assistant Programmer, AD (ICT), Assistant Network Engineer, Assistant Maintenance Engineer, and Assistant Database Administrator positions.

\begin{center}
\vspace{0.5em}
{\color{primary}\rule{0.45\textwidth}{1.2pt}}\\[4pt]
{\Large\bfseries\color{primary} Section I: Linked Lists}\\[2pt]
{\color{primary}\rule{0.45\textwidth}{1.2pt}}
\vspace{0.5em}
\end{center}

\section{What is a Linked List?}

A \textbf{linked list} is a linear data structure consisting of a sequence of nodes. Each node stores data and one or more links (pointers) to other nodes.

Unlike an array, whose elements are conventionally stored in contiguous memory locations, the nodes of a linked list can occupy different memory locations across the heap. Pointers connect them into a coherent logical sequence.

A node in a singly linked list contains two fields:
\begin{itemize}
    \item \textbf{Data:} The value stored in the node.
    \item \textbf{Next pointer:} The memory address of the next node in the sequence.
\end{itemize}

For example, consider a list containing $10$, $20$, and $30$. The first node stores $10$ and points to the node containing $20$; the second points to the node containing $30$; the final node's next pointer is \texttt{NULL} (or \texttt{nullptr}) in a conventional, non-circular singly linked list:
\[
\text{Head} \longrightarrow [10 \mid \bullet] \longrightarrow [20 \mid \bullet] \longrightarrow [30 \mid \text{nullptr}]
\]

\subsection{Important Terminology}

\begin{table}[htbp]
\centering
\small
\begin{tabularx}{\textwidth}{l X}
\toprule
\textbf{Term} & \textbf{Meaning} \\
\midrule
Node                      & An individual unit containing data and pointer/link fields \\
Head                      & A pointer or reference addressing the first node in the list \\
Tail                      & A pointer or reference addressing the last node, if maintained \\
Link                      & A pointer or reference connecting one node to another \\
\texttt{NULL} / \texttt{nullptr} & A sentinel value indicating the absence of a successor node \\
Traversal                 & Visiting each node sequentially from head to end \\
Insertion                 & Allocating and linking a new node into the list \\
Deletion                  & Unlinking and deallocating an existing node \\
Dynamic Memory Allocation & Allocating node memory at runtime as needed \\
\bottomrule
\end{tabularx}
\caption{Fundamental linked list terminology.}
\end{table}

In C++, nodes are commonly created dynamically using \texttt{new} and deallocated using \texttt{delete}. In C, \texttt{malloc()} and \texttt{free()} are standard.

\section{Why Use Linked Lists Instead of Arrays?}

An array is usually preferable when fast indexed access ($O(1)$ random access) is important. A linked list is advantageous when the application performs frequent insertions or deletions at positions where the target node or link is already known.

Consider an array containing:
\[
[10, \; 20, \; 30, \; 40]
\]
Inserting $25$ between $20$ and $30$ requires shifting all subsequent elements ($30$ and $40$) to the right, incurring $O(n)$ time.

In a linked list, if a pointer to the node containing $20$ is already available, a new node containing $25$ can be inserted in $O(1)$ time by adjusting two pointers. Existing nodes remain stationary in memory. However, locating the node containing $20$ from the head may itself require an $O(n)$ linear traversal.

\subsection{Array versus Linked List Comparison}

\begin{table}[htbp]
\centering
\footnotesize
\begin{tabularx}{\textwidth}{l X X}
\toprule
\textbf{Basis} & \textbf{Array} & \textbf{Linked List} \\
\midrule
Memory Layout                   & Contiguous memory block                           & Scattered nodes linked via pointers \\
Size Flexibility                & Fixed (static) or reallocation cost (dynamic)      & Fully dynamic; grows/shrinks per node \\
Random Access                   & $O(1)$ via index arithmetic                       & $O(n)$ sequential link traversal \\
Access to First Element         & $O(1)$                                            & $O(1)$ via \texttt{head} pointer \\
Insert at Beginning             & $O(n)$ due to element shifting                    & $O(1)$ pointer reassignment \\
Delete at Beginning             & $O(n)$ due to element shifting                    & $O(1)$ pointer reassignment \\
Insert After Known Node         & $O(n)$ due to element shifting                    & $O(1)$ pointer updates \\
Search in Unsorted Data         & $O(n)$                                            & $O(n)$ \\
Extra Memory Overhead           & Minimal structural overhead                       & Extra memory per node for pointers \\
Cache Locality                  & Excellent (spatial locality)                      & Generally poor (cache misses common) \\
Binary Search Capability        & $O(\log n)$ on sorted array                       & Infeasible for $O(\log n)$ (middle node is not $O(1)$) \\
\bottomrule
\end{tabularx}
\caption{Architectural comparison between Arrays and Linked Lists.}
\end{table}

\begin{quote}
\textbf{Important Exam Distinction:} Inserting a node into a linked list takes $O(1)$ time \emph{if the predecessor pointer is already held}. Finding that position from scratch still takes $O(n)$ traversal time.
\end{quote}

\section{Singly Linked List}

\subsection{Definition and Node Structure}

A \textbf{singly linked list (SLL)} is a linked list in which each node contains data and exactly one pointer to the next node. Each node can directly identify its successor, but cannot directly identify its predecessor.

The first node is accessed through \texttt{head}. The last node is identified by \texttt{next == nullptr}.

\begin{lstlisting}[style=cppstyle]
struct Node {
    int data;
    Node* next;

    Node(int value) : data(value), next(nullptr) {}
};
\end{lstlisting}

\subsubsection{Main Characteristics}
\begin{enumerate}
    \item Each node has exactly one link field.
    \item Traversal proceeds strictly forward.
    \item Nodes are inserted or removed without shifting remaining elements in memory.
    \item Reverse traversal is not directly supported without auxiliary structures (e.g., stack) or pointer reversal.
    \item Accessing the $k$-th element requires $O(k)$ pointer dereferences from \texttt{head}.
\end{enumerate}

\subsection{Traversing a Singly Linked List}

Traversal visits every node in order from \texttt{head} until \texttt{nullptr} is reached.

\begin{lstlisting}[style=cppstyle]
void display(Node* head) {
    Node* current = head;
    while (current != nullptr) {
        cout << current->data << " ";
        current = current->next;
    }
    cout << '\n';
}
\end{lstlisting}

\textbf{Complexity:} Time is $O(n)$ as all $n$ nodes are visited once; auxiliary space is $O(1)$.

\subsection{Insertion Operations}

\subsubsection{A. Insert at the Beginning (Front)}
To insert value $x$ before the current head:
\begin{enumerate}
    \item Allocate a new node containing $x$.
    \item Point \texttt{newNode->next} to the current \texttt{head}.
    \item Reassign \texttt{head} to \texttt{newNode}.
\end{enumerate}

\begin{lstlisting}[style=cppstyle]
void insertFront(Node*& head, int value) {
    Node* newNode = new Node(value);
    newNode->next = head;
    head = newNode;
}
\end{lstlisting}

\textbf{Time Complexity:} $O(1)$. If the list is initially empty (\texttt{head == nullptr}), \texttt{newNode->next} becomes \texttt{nullptr}, and \texttt{newNode} correctly becomes the head.

\subsubsection{B. Insert at the End (Tail)}
To append a node at the tail:
\begin{enumerate}
    \item Allocate the new node with \texttt{next = nullptr}.
    \item If \texttt{head == nullptr}, set \texttt{head = newNode}.
    \item Otherwise, traverse from \texttt{head} to the node where \texttt{current->next == nullptr}.
    \item Assign \texttt{current->next = newNode}.
\end{enumerate}

\begin{lstlisting}[style=cppstyle]
void insertEnd(Node*& head, int value) {
    Node* newNode = new Node(value);
    if (head == nullptr) {
        head = newNode;
        return;
    }
    Node* current = head;
    while (current->next != nullptr) {
        current = current->next;
    }
    current->next = newNode;
}
\end{lstlisting}

\textbf{Time Complexity:} $O(n)$ when only \texttt{head} is tracked. If an explicit \texttt{tail} pointer is maintained, insertion at the end is $O(1)$.

\subsubsection{C. Insert After a Specified Node}
Suppose \texttt{current} points to the node after which the new node is to be inserted:
\begin{enumerate}
    \item Allocate \texttt{newNode}.
    \item Set \texttt{newNode->next = current->next}.
    \item Set \texttt{current->next = newNode}.
\end{enumerate}

\begin{lstlisting}[style=cppstyle]
void insertAfter(Node* current, int value) {
    if (current == nullptr) return;
    Node* newNode = new Node(value);
    newNode->next = current->next;
    current->next = newNode;
}
\end{lstlisting}

\begin{quote}
\textbf{Common Exam Trap:} Always set \texttt{newNode->next} \emph{before} overwriting \texttt{current->next}. Overwriting \texttt{current->next} first breaks the link to the remainder of the list, causing memory leaks and orphaned nodes.
\end{quote}

\subsection{Deletion Operations}

\subsubsection{A. Delete the First Node (Front)}
\begin{lstlisting}[style=cppstyle]
void deleteFront(Node*& head) {
    if (head == nullptr) return;
    Node* temp = head;
    head = head->next;
    delete temp;
}
\end{lstlisting}
\textbf{Time Complexity:} $O(1)$.

\subsubsection{B. Delete a Node Containing a Target Value}
Deleting a value requires maintaining two pointers: \texttt{current} (the node being examined) and \texttt{previous} (the node preceding \texttt{current}).

\begin{lstlisting}[style=cppstyle]
void deleteValue(Node*& head, int value) {
    Node* current = head;
    Node* previous = nullptr;

    while (current != nullptr && current->data != value) {
        previous = current;
        current = current->next;
    }

    if (current == nullptr) return; // Value not found

    if (previous == nullptr) {
        head = current->next; // Target is head
    } else {
        previous->next = current->next; // Bypass current
    }
    delete current;
}
\end{lstlisting}

\textbf{Time Complexity:} $O(n)$ worst-case to locate the element.

\subsection{Advantages and Disadvantages of Singly Linked Lists}

\textbf{Advantages:}
\begin{itemize}
    \item Dynamic sizing without pre-allocation overhead.
    \item $O(1)$ insertion and deletion at the front.
    \item No need to shift memory elements upon modification.
    \item Simpler pointer management and lower memory overhead than doubly linked lists.
\end{itemize}

\textbf{Disadvantages:}
\begin{itemize}
    \item Unidirectional traversal only.
    \item No $O(1)$ random access.
    \item Deleting an arbitrary known node requires locating its predecessor ($O(n)$ search).
    \item Extra pointer memory per node ($4$ or $8$ bytes).
    \item Unfriendly to CPU cache lines due to non-contiguous heap allocations.
\end{itemize}

\subsection{Applications}
\begin{itemize}
    \item Implementing dynamic stacks and queues.
    \item Representing adjacency lists in graph algorithms.
    \item Separate chaining collision resolution in hash tables.
    \item Polynomial manipulation and arbitrary-precision arithmetic.
\end{itemize}

\section{Doubly Linked List}

\subsection{Definition and Node Structure}

A \textbf{doubly linked list (DLL)} is a linked list in which each node contains:
\begin{enumerate}
    \item Data payload.
    \item A pointer to the previous node (\texttt{prev}).
    \item A pointer to the next node (\texttt{next}).
\end{enumerate}

The first node has \texttt{prev == nullptr}, and the last node has \texttt{next == nullptr}.

\begin{lstlisting}[style=cppstyle]
struct DNode {
    int data;
    DNode* prev;
    DNode* next;

    DNode(int value) : data(value), prev(nullptr), next(nullptr) {}
};
\end{lstlisting}

\subsection{Singly versus Doubly Linked Lists}

\textbf{Advantages of DLL over SLL:}
\begin{itemize}
    \item Supports bidirectional traversal (forward and backward).
    \item Deletion of a known node is $O(1)$ because the predecessor is accessible directly via \texttt{node->prev}.
    \item Inserting before a known node is $O(1)$.
\end{itemize}

\textbf{Disadvantages of DLL:}
\begin{itemize}
    \item Requires two pointer fields per node, increasing memory overhead.
    \item Every insertion and deletion requires updating up to four pointers, increasing code complexity and bug potential.
\end{itemize}

\subsection{Insertion Operations}

\subsubsection{A. Insert at the Beginning}
\begin{lstlisting}[style=cppstyle]
void insertFront(DNode*& head, int value) {
    DNode* newNode = new DNode(value);
    newNode->next = head;
    if (head != nullptr) {
        head->prev = newNode;
    }
    head = newNode;
}
\end{lstlisting}
\textbf{Time Complexity:} $O(1)$.

\subsubsection{B. Insert After a Known Node}
\begin{lstlisting}[style=cppstyle]
void insertAfter(DNode* current, int value) {
    if (current == nullptr) return;
    DNode* newNode = new DNode(value);

    newNode->next = current->next;
    newNode->prev = current;

    if (current->next != nullptr) {
        current->next->prev = newNode;
    }
    current->next = newNode;
}
\end{lstlisting}
\textbf{Time Complexity:} $O(1)$ when \texttt{current} is given.

\subsection{Deletion of a Known Node}

A major advantage of a DLL is unlinking a known node in $O(1)$ without needing to traverse from the head:

\begin{lstlisting}[style=cppstyle]
void deleteNode(DNode*& head, DNode* node) {
    if (node == nullptr) return;

    if (node->prev != nullptr) {
        node->prev->next = node->next;
    } else {
        head = node->next; // node was head
    }

    if (node->next != nullptr) {
        node->next->prev = node->prev;
    }
    delete node;
}
\end{lstlisting}
\textbf{Time Complexity:} $O(1)$.

\subsection{Applications of Doubly Linked Lists}
\begin{itemize}
    \item Web browser navigation history (Back and Forward buttons).
    \item Undo and Redo functionality in editors.
    \item Least Recently Used (LRU) Cache implementation (combined with a hash map).
    \item Media player playlists with Next and Previous track navigation.
    \item Operating system thread and process schedulers.
\end{itemize}

\section{Circular Linked List}

\subsection{Definition and Principal Variants}

A \textbf{circular linked list} is a linked list in which the last node links back to the first node instead of holding \texttt{nullptr}. There are two variants:
\begin{enumerate}
    \item \textbf{Circular Singly Linked List (CSLL):} The last node's \texttt{next} points to \texttt{head}.
    \item \textbf{Circular Doubly Linked List (CDLL):} The last node's \texttt{next} points to \texttt{head}, and \texttt{head->prev} points to the last node.
\end{enumerate}

\subsection{Traversing a Circular Singly Linked List}

Because a circular list contains no \texttt{nullptr}, a standard \texttt{while (current != nullptr)} loop will run indefinitely. Instead, use a \texttt{do-while} loop:

\begin{lstlisting}[style=cppstyle]
void displayCircular(Node* head) {
    if (head == nullptr) return;
    Node* current = head;
    do {
        cout << current->data << " ";
        current = current->next;
    } while (current != head);
    cout << '\n';
}
\end{lstlisting}

\subsection{Insertion Using a Tail Pointer}

Maintaining a pointer to the \textbf{last node (\texttt{tail})} provides $O(1)$ insertion at both the head and the tail because:
\[
\text{Head} = \texttt{tail->next}
\]

\begin{lstlisting}[style=cppstyle]
class CircularLinkedList {
private:
    struct Node {
        int data;
        Node* next;
        Node(int val) : data(val), next(nullptr) {}
    };
    Node* tail = nullptr;

public:
    // Insert at front: O(1)
    void insertFront(int value) {
        Node* newNode = new Node(value);
        if (tail == nullptr) {
            newNode->next = newNode;
            tail = newNode;
            return;
        }
        newNode->next = tail->next;
        tail->next = newNode;
    }

    // Insert at end: O(1)
    void insertEnd(int value) {
        Node* newNode = new Node(value);
        if (tail == nullptr) {
            newNode->next = newNode;
            tail = newNode;
            return;
        }
        newNode->next = tail->next;
        tail->next = newNode;
        tail = newNode; // Only difference: tail moves to newNode
    }

    // Delete front: O(1)
    void deleteFront() {
        if (tail == nullptr) return;
        Node* head = tail->next;
        if (head == tail) {
            delete head;
            tail = nullptr;
            return;
        }
        tail->next = head->next;
        delete head;
    }

    void display() const {
        if (tail == nullptr) return;
        Node* current = tail->next;
        do {
            cout << current->data << " ";
            current = current->next;
        } while (current != tail->next);
        cout << '\n';
    }

    ~CircularLinkedList() {
        while (tail != nullptr) deleteFront();
    }
};
\end{lstlisting}

\subsection{Applications of Circular Linked Lists}
\begin{itemize}
    \item \textbf{Round-Robin CPU Scheduling:} Processes are allocated fixed time-slices in a cyclic rotation.
    \item \textbf{Continuous Media Streaming:} Repeating audio playlists or slideshow loops.
    \item \textbf{Multiplayer Turn-Based Games:} Cycling through players continuously.
\end{itemize}

\section{Comparison of Linked-List Types}

\begin{table}[htbp]
\centering
\small
\begin{tabularx}{\textwidth}{l X X X}
\toprule
\textbf{Feature} & \textbf{Singly Linked List} & \textbf{Doubly Linked List} & \textbf{Circular Singly Linked List} \\
\midrule
Links per Node       & One (\texttt{next})         & Two (\texttt{prev}, \texttt{next}) & One (\texttt{next}) \\
Traversal Direction  & Forward only                & Bidirectional               & Forward (cyclic) \\
End of List          & \texttt{next == nullptr}    & \texttt{next == nullptr}    & \texttt{next == head} \\
Insert at Head       & $O(1)$                      & $O(1)$                      & $O(1)$ (with tail pointer) \\
Insert at End        & $O(n)$ ($O(1)$ with tail)   & $O(n)$ ($O(1)$ with tail)   & $O(1)$ (with tail pointer) \\
Delete First Node    & $O(1)$                      & $O(1)$                      & $O(1)$ (with tail pointer) \\
Delete Known Node    & Requires predecessor ($O(n)$) & $O(1)$ unlinking           & Requires predecessor ($O(n)$) \\
Stopping Condition   & Reach \texttt{nullptr}      & Reach \texttt{nullptr}      & Return to starting node \\
Pointer Overhead     & 1 pointer per node          & 2 pointers per node         & 1 pointer per node \\
\bottomrule
\end{tabularx}
\caption{Comprehensive comparison of linked-list variants.}
\end{table}

\begin{center}
\vspace{1em}
{\color{primary}\rule{0.45\textwidth}{1.2pt}}\\[4pt]
{\Large\bfseries\color{primary} Section II: Heap Data Structures}\\[2pt]
{\color{primary}\rule{0.45\textwidth}{1.2pt}}
\vspace{0.5em}
\end{center}

\section{What is a Heap?}

A \textbf{heap} is a specialized tree-based data structure that satisfies two fundamental properties:
\begin{enumerate}
    \item \textbf{Shape Property:} The tree is a \textbf{complete binary tree}.
    \item \textbf{Heap Property:} The keys satisfy either the \textbf{max-heap} or \textbf{min-heap} ordering rule.
\end{enumerate}

\subsection{What is a Complete Binary Tree?}

A binary tree is \textbf{complete} if:
\begin{itemize}
    \item Every level, except possibly the last, is completely filled with nodes.
    \item All nodes at the last level reside as far left as possible, without gaps.
\end{itemize}

This property enables a binary heap to be represented compactly inside a contiguous 1D array without storing explicit parent or child pointers.

\begin{table}[htbp]
\centering
\small
\begin{tabularx}{\textwidth}{l X}
\toprule
\textbf{Binary Tree Type} & \textbf{Formal Definition} \\
\midrule
Complete Binary Tree & All levels full except possibly the last, which is filled strictly left-to-right \\
Full Binary Tree     & Every internal node has exactly two children (no node has 1 child) \\
Perfect Binary Tree  & All internal nodes have two children and all leaves reside at identical depth \\
\bottomrule
\end{tabularx}
\caption{Distinguishing Complete, Full, and Perfect binary trees.}
\end{table}

\subsection{Max-Heap vs Min-Heap}

\begin{itemize}
    \item \textbf{Max-Heap:} For every node $i$ other than the root:
    \[
    A[\text{parent}(i)] \geq A[i]
    \]
    The maximum element resides perpetually at the root ($A[0]$).
    \item \textbf{Min-Heap:} For every node $i$ other than the root:
    \[
    A[\text{parent}(i)] \leq A[i]
    \]
    The minimum element resides perpetually at the root ($A[0]$).
\end{itemize}

\subsection{Is a Heap the Same as a Binary Search Tree (BST)?}

No. This is a common viva question:
\begin{table}[htbp]
\centering
\small
\begin{tabularx}{\textwidth}{l X X}
\toprule
\textbf{Basis} & \textbf{Binary Heap} & \textbf{Binary Search Tree (BST)} \\
\midrule
Ordering Rule     & Vertical ordering: Parent $\geq$ children (no left vs right rule) & Horizontal ordering: Left $<$ Root $<$ Right \\
Root Node         & Always the global Maximum (max-heap) or Minimum (min-heap)        & Arbitrary key depending on insertion \\
Access to Extremum & $O(1)$                                                            & $O(h)$ ($O(\log n)$ balanced, $O(n)$ skewed) \\
Search for Arbitrary Key & $O(n)$ linear scan (no directional guarantee)                & $O(\log n)$ average / balanced \\
Underlying Storage & Implicit 1D Array (zero pointer overhead)                         & Explicit linked nodes with left/right pointers \\
\bottomrule
\end{tabularx}
\caption{Key distinctions between Binary Heap and Binary Search Tree.}
\end{table}

\section{Representing a Binary Heap Using an Array}

Because a heap is a complete binary tree, parent-child relationships map directly to arithmetic indices.

\subsection{Zero-Based Indexing Formulas}
For a node located at index $i$ in a zero-indexed array ($i = 0, 1, \dots, n-1$):
\begin{align*}
\text{Parent}(i) &= \left\lfloor \frac{i - 1}{2} \right\rfloor \quad (\text{for } i > 0) \\
\text{Left Child}(i) &= 2i + 1 \\
\text{Right Child}(i) &= 2i + 2
\end{align*}

\subsection{One-Based Indexing Formulas}
For a one-indexed array ($i = 1, 2, \dots, n$):
\begin{align*}
\text{Parent}(i) &= \left\lfloor \frac{i}{2} \right\rfloor \\
\text{Left Child}(i) &= 2i \\
\text{Right Child}(i) &= 2i + 1
\end{align*}

\subsubsection{Example Mapping}
Consider the array $[50, 30, 40, 10, 20]$ of size $n=5$:
\begin{table}[htbp]
\centering
\small
\begin{tabularx}{0.85\textwidth}{ccccc}
\toprule
\textbf{Index $i$} & \textbf{Value $A[i]$} & \textbf{Parent Index} & \textbf{Left Child ($2i+1$)} & \textbf{Right Child ($2i+2$)} \\
\midrule
0 & 50 & None & 1 (value 30) & 2 (value 40) \\
1 & 30 & 0    & 3 (value 10) & 4 (value 20) \\
2 & 40 & 0    & 5 (out of range) & 6 (out of range) \\
3 & 10 & 1    & --- & --- \\
4 & 20 & 1    & --- & --- \\
\bottomrule
\end{tabularx}
\caption{Zero-based array index mapping of a binary heap.}
\end{table}

\section{Max-Heap Construction: Step by Step}

Given an arbitrary array, two strategies exist to convert it into a valid max-heap:
\begin{enumerate}
    \item \textbf{Repeated Insertion:} Insert elements one by one, sifting up ($O(n \log n)$ time).
    \item \textbf{Bottom-Up Heap Construction (Floyd's Algorithm):} Place elements into array, then call \texttt{maxHeapify} downward starting from the last non-leaf node up to the root ($O(n)$ time).
\end{enumerate}

\subsection{The Heapify (Sift-Down) Operation}

\texttt{maxHeapify(A, i, n)} assumes that the subtrees rooted at $\text{left}(i)$ and $\text{right}(i)$ are valid max-heaps, but $A[i]$ may violate the heap property. It floats $A[i]$ downward by repeatedly swapping with its largest child until the property is restored.

\subsection{Step-by-Step Bottom-Up Construction Trace}

Given the input array:
\[
A = [4, \; 10, \; 3, \; 5, \; 1], \quad n = 5
\]
The last non-leaf node index in a zero-indexed array is:
\[
\text{Last non-leaf} = \left\lfloor \frac{n}{2} \right\rfloor - 1 = \left\lfloor \frac{5}{2} \right\rfloor - 1 = 1
\]
Leaves reside at indices $2, 3, 4$ and are trivial heaps. We only process indices $1$ down to $0$.

\begin{enumerate}
    \item \textbf{Step 1: Heapify Node at Index 1 ($A[1] = 10$)}
    \begin{itemize}
        \item Left child at index $3$: $A[3] = 5$.
        \item Right child at index $4$: $A[4] = 1$.
        \item Since $10 \geq 5$ and $10 \geq 1$, node $1$ already satisfies max-heap ordering. No swap.
        \item Array remains: $[4, 10, 3, 5, 1]$.
    \end{itemize}

    \item \textbf{Step 2: Heapify Node at Index 0 ($A[0] = 4$)}
    \begin{itemize}
        \item Left child at index $1$: $A[1] = 10$.
        \item Right child at index $2$: $A[2] = 3$.
        \item Largest among $\{4, 10, 3\}$ is $10$ at index $1$.
        \item Swap $A[0]$ and $A[1]$.
        \item Array becomes: $[10, 4, 3, 5, 1]$.
    \end{itemize}

    \item \textbf{Step 3: Continue Sifting Down Displaced Node ($4$ at Index 1)}
    \begin{itemize}
        \item Children of index $1$ are $A[3] = 5$ and $A[4] = 1$.
        \item Largest among $\{4, 5, 1\}$ is $5$ at index $3$.
        \item Swap $A[1]$ and $A[3]$.
        \item Array becomes: $[10, 5, 3, 4, 1]$.
        \item Node $4$ is now at index $3$ (a leaf), so sift-down terminates.
    \end{itemize}
\end{enumerate}

\textbf{Final Valid Max-Heap:} $[10, 5, 3, 4, 1]$.

\begin{table}[htbp]
\centering
\small
\begin{tabularx}{0.85\textwidth}{cc X c}
\toprule
\textbf{Step} & \textbf{Index Processed} & \textbf{Action Taken} & \textbf{Array State} \\
\midrule
Initial & --- & Original sequence & $[4, 10, 3, 5, 1]$ \\
1       & 1   & Checked children $\{5, 1\}$; no swap & $[4, 10, 3, 5, 1]$ \\
2       & 0   & Swapped $4 \leftrightarrow 10$ with child & $[10, 4, 3, 5, 1]$ \\
3       & 1   & Swapped $4 \leftrightarrow 5$ to restore leaf & $[10, 5, 3, 4, 1]$ \\
\bottomrule
\end{tabularx}
\caption{Trace of bottom-up max-heap construction.}
\end{table}

\subsection{Why Does Bottom-Up Heap Construction Take $O(n)$?}

A single \texttt{heapify} takes $O(\log n)$ worst-case time. However, bottom-up construction does not run for full height on all nodes:
\begin{itemize}
    \item Half the nodes are leaves ($h=0$), requiring $0$ operations.
    \item A quarter of the nodes have height $1$, moving down at most $1$ step.
    \item Only $1$ node (the root) can travel the full tree height $\lfloor \log n \rfloor$.
\end{itemize}

Summing work over all tree heights:
\[
T(n) = \sum_{h=0}^{\lfloor \log n \rfloor} \left\lceil \frac{n}{2^{h+1}} \right\rceil O(h) = O\left( n \sum_{h=0}^{\infty} \frac{h}{2^h} \right)
\]
Since the infinite series $\sum_{h=0}^{\infty} \frac{h}{2^h} = 2$, we obtain:
\[
T(n) = O(2n) = O(n)
\]
This proves that bottom-up construction runs in \textbf{linear time $O(n)$}.

\section{Max-Heap Insertion}

Insertion maintains both the complete binary tree shape and the max-heap ordering via \textbf{sift-up} (percolate-up).

\subsection{Algorithm}
\begin{enumerate}
    \item Append the new value $x$ at the end of the array ($A[n]$), expanding the heap size.
    \item Compare $x$ with its parent: $p = \lfloor (i - 1)/2 \rfloor$.
    \item If $x > A[p]$, swap $A[i]$ and $A[p]$, and set $i = p$.
    \item Repeat until $x \leq A[p]$ or $x$ reaches the root ($i = 0$).
\end{enumerate}

\subsection{Step-by-Step Insertion Example}
Insert $15$ into existing max-heap $[10, 5, 3, 4, 1]$:
\begin{enumerate}
    \item Append $15$ at index $5$: $[10, 5, 3, 4, 1, \mathbf{15}]$.
    \item Parent of index $5$ is $\lfloor(5-1)/2\rfloor = 2$ ($A[2] = 3$). Since $15 > 3$, swap:
    \[
    [10, 5, \mathbf{15}, 4, 1, \mathbf{3}]
    \]
    \item New parent of index $2$ is $\lfloor(2-1)/2\rfloor = 0$ ($A[0] = 10$). Since $15 > 10$, swap:
    \[
    [\mathbf{15}, 5, \mathbf{10}, 4, 1, 3]
    \]
    \item Index $0$ is the root. Insertion terminates.
\end{enumerate}
\textbf{Complexity:} Worst-case time is $O(\log n)$ (tree height). Auxiliary space is $O(1)$.

\section{Max-Heap Deletion (Extract Maximum)}

Deleting from a max-heap refers to extracting the root node containing the maximum key.

\subsection{Algorithm}
\begin{enumerate}
    \item Save root $A[0]$ for return.
    \item Copy the last element in the array to the root ($A[0] = A[n-1]$).
    \item Decrement heap size by $1$ (removing the last position).
    \item Call \texttt{maxHeapify(A, 0, n)} to sift the new root downward until the max-heap property is restored.
\end{enumerate}

\subsection{Step-by-Step Deletion Example}
Extract max from $[15, 5, 10, 4, 1, 3]$:
\begin{enumerate}
    \item Save maximum: $\text{max} = 15$.
    \item Replace root with last element $3$: $[3, 5, 10, 4, 1]$.
    \item Children of root ($3$) are $A[1] = 5$ and $A[2] = 10$. Largest child is $10$ at index $2$.
    \item Swap $3$ and $10$:
    \[
    [10, 5, \mathbf{3}, 4, 1]
    \]
    \item Node $3$ at index $2$ has no children ($2(2)+1 = 5 \geq 5$). Sift-down terminates.
\end{enumerate}
\textbf{Complexity:} $O(\log n)$ worst-case time.

\subsection{Deleting an Arbitrary Key}
To delete an arbitrary key $k$ (not necessarily the root):
\begin{enumerate}
    \item Search the heap array for $k$. Since a heap has no directional search property, finding $k$ takes $O(n)$ time.
    \item Overwrite $k$ with the last element in the array and decrement the size.
    \item Repair the heap: if the replacement is greater than its parent, sift up; if smaller than either child, sift down. Repair takes $O(\log n)$.
    \item Total time: $O(n)$ due to linear search.
\end{enumerate}

\section{Complete C++ Implementation of a Max-Heap}

\begin{lstlisting}[style=cppstyle]
#include <iostream>
#include <vector>
#include <stdexcept>
#include <utility>
using namespace std;

class MaxHeap {
private:
    vector<int> heap;

    void siftUp(size_t i) {
        while (i > 0) {
            size_t parent = (i - 1) / 2;
            if (heap[parent] >= heap[i]) break;
            swap(heap[parent], heap[i]);
            i = parent;
        }
    }

    void siftDown(size_t i) {
        size_t n = heap.size();
        while (true) {
            size_t largest = i;
            size_t left = 2 * i + 1;
            size_t right = 2 * i + 2;

            if (left < n && heap[left] > heap[largest]) {
                largest = left;
            }
            if (right < n && heap[right] > heap[largest]) {
                largest = right;
            }
            if (largest == i) break;

            swap(heap[i], heap[largest]);
            i = largest;
        }
    }

public:
    // O(n) Bottom-Up constructor
    explicit MaxHeap(vector<int> values) : heap(std::move(values)) {
        for (size_t i = heap.size() / 2; i > 0; --i) {
            siftDown(i - 1);
        }
    }

    bool empty() const { return heap.empty(); }
    size_t size() const { return heap.size(); }

    // O(1) Max element inspection
    int getMax() const {
        if (heap.empty()) throw underflow_error("Heap is empty");
        return heap[0];
    }

    // O(log n) Insertion
    void insert(int value) {
        heap.push_back(value);
        siftUp(heap.size() - 1);
    }

    // O(log n) Extract Maximum
    int extractMax() {
        if (heap.empty()) throw underflow_error("Heap is empty");
        int maximum = heap[0];
        if (heap.size() == 1) {
            heap.pop_back();
            return maximum;
        }
        heap[0] = heap.back();
        heap.pop_back();
        siftDown(0);
        return maximum;
    }

    void display() const {
        for (int val : heap) cout << val << " ";
        cout << '\n';
    }
};

int main() {
    MaxHeap h(vector<int>{4, 10, 3, 5, 1});

    cout << "After construction: ";
    h.display();

    h.insert(15);
    cout << "After inserting 15: ";
    h.display();

    cout << "Maximum element: " << h.getMax() << '\n';
    cout << "Extracted maximum: " << h.extractMax() << '\n';

    cout << "After deleting maximum: ";
    h.display();
    cout << "Heap size: " << h.size() << '\n';

    return 0;
}
\end{lstlisting}

\subsubsection{Expected Output}
\begin{lstlisting}[style=pseudocode]
After construction: 10 5 3 4 1
After inserting 15: 15 5 10 4 1 3
Maximum element: 15
Extracted maximum: 15
After deleting maximum: 10 5 3 4 1
Heap size: 5
\end{lstlisting}

\section{Heap Complexity: Complete Comparison}

\begin{table}[htbp]
\centering
\small
\begin{tabularx}{\textwidth}{l X X}
\toprule
\textbf{Operation} & \textbf{Time Complexity} & \textbf{Algorithmic Reason} \\
\midrule
Find Maximum / Minimum    & $O(1)$                    & Root is always the extremum \\
Insert Element            & $O(\log n)$ worst-case    & Sifts up tree height $\lfloor \log_2 n \rfloor$ \\
Delete Max / Min (Root)   & $O(\log n)$ worst-case    & Sifts down tree height \\
Heapify Single Node       & $O(\log n)$ worst-case    & Floats down subtree height \\
Build Heap (Bottom-Up)    & $O(n)$                    & Leaf-heavy distribution of complete tree \\
Build Heap (Insertion)    & $O(n \log n)$             & Repeated insertions \\
Search for Arbitrary Key  & $O(n)$                    & No left/right directional ordering \\
Heap Sort                 & $O(n \log n)$             & $n$ extractions of $O(\log n)$ cost \\
Storage for $n$ elements  & $O(n)$                    & Compact 1D array \\
\bottomrule
\end{tabularx}
\caption{Comprehensive operational complexity reference for Binary Heaps.}
\end{table}

\subsection{Binary Heap versus Priority Queue}
A \textbf{priority queue} is an Abstract Data Type (ADT) defining operations (\texttt{push}, \texttt{top}, \texttt{pop}). A \textbf{binary heap} is the standard concrete data structure used to implement a priority queue efficiently.

In C++ STL:
\begin{lstlisting}[style=cppstyle]
#include <queue>
#include <vector>
#include <functional>

// Default: Max-Priority Queue
priority_queue<int> maxPQ;

// Min-Priority Queue
priority_queue<int, vector<int>, greater<int>> minPQ;
\end{lstlisting}

\section{Heap Sort}

Heap sort is an in-place, comparison-based sorting algorithm using a binary heap.

\subsection{Algorithm (Ascending Order)}
\begin{enumerate}
    \item Build a max-heap from the input array ($O(n)$).
    \item Swap $A[0]$ (the maximum) with the last element $A[n-1]$ in the active heap.
    \item Reduce active heap size by $1$.
    \item Sift down the new root ($A[0]$) across the active heap ($O(\log n)$).
    \item Repeat steps 2--4 until active heap size is $1$.
\end{enumerate}

\subsection{Properties of Heap Sort}
\begin{itemize}
    \item \textbf{Time Complexity:} $O(n \log n)$ in Best, Average, and Worst cases.
    \item \textbf{Auxiliary Space:} $O(1)$ in-place.
    \item \textbf{Stability:} \textbf{Not stable}. Swapping elements across the heap destroys the relative order of identical keys.
\end{itemize}

\section{Common Heap Errors and Examination Traps}

\begin{enumerate}
    \item \textbf{Assuming Heap Array is Sorted:} A heap is \emph{not} a sorted array. Only parent-child vertical ordering holds.
    \item \textbf{Confusing Heap with BST:} Left children are not necessarily smaller than right children in a heap.
    \item \textbf{Sifting in Wrong Direction:} Insertion requires \emph{sift-up}; root deletion requires \emph{sift-down}; bottom-up construction runs \emph{sift-down} starting from the last internal node.
    \item \textbf{Selecting Wrong Child in Sift-Down:} In a max-heap, you must swap with the \textbf{larger} of the two children. Swapping with the smaller child leaves the heap property violated.
    \item \textbf{Starting Bottom-Up at the Wrong Index:} Begin at $\lfloor n/2 \rfloor - 1$, not at $n-1$. Leaves have no children and are already valid heaps.
    \item \textbf{Confusing Memory Heap with Data Structure Heap:} The runtime memory heap used by \texttt{malloc} or \texttt{new} is an operating system memory management pool, entirely distinct from the binary heap data structure.
\end{enumerate}

\section{Practical Applications of Heaps}
\begin{itemize}
    \item \textbf{Operating System Schedulers:} Priority scheduling for CPU jobs.
    \item \textbf{Graph Algorithms:} Dijkstra's shortest path and Prim's Minimum Spanning Tree (using a min-priority queue).
    \item \textbf{Top-$K$ Selection Problems:} Maintaining the $k$ largest or smallest elements from a massive streaming dataset in $O(n \log k)$ time using a bounded heap.
    \item \textbf{Discrete Event Simulation:} Processing events chronologically based on timestamps.
    \item \textbf{Huffman Coding:} Constructing optimal prefix trees for data compression.
\end{itemize}

\section{Likely Recruitment Examination Questions: Linked Lists and Heaps}

\begin{table}[htbp]
\centering
\small
\begin{tabularx}{\textwidth}{p{7cm} X}
\toprule
\textbf{Question} & \textbf{Key Answer} \\
\midrule
What is a linked list? & A linear sequence of nodes connected via pointers. \\
Difference between singly and doubly linked list? & Singly has 1 pointer (\texttt{next}); doubly has 2 pointers (\texttt{prev}, \texttt{next}). \\
How does circular linked list differ from singly? & Last node points back to the head node instead of \texttt{nullptr}. \\
Time complexity of insertion at head? & $O(1)$ pointer update. \\
Why does linked list not support $O(1)$ random access? & Nodes are non-contiguous; reaching index $k$ requires sequential traversal. \\
What is a binary heap? & A complete binary tree satisfying a heap-ordering rule. \\
What is at the root of a max-heap? & The maximum element in the heap. \\
Formula for last non-leaf in 0-indexed heap? & $\lfloor n/2 \rfloor - 1$. \\
Formulas for left and right children (0-based)? & Left: $2i + 1$; Right: $2i + 2$. \\
Formula for parent index (0-based)? & $\lfloor (i - 1) / 2 \rfloor$. \\
Complexity of bottom-up heap construction? & $O(n)$ linear time. \\
Complexity of insertion into binary heap? & $O(\log n)$ worst-case. \\
Complexity of deleting the root from heap? & $O(\log n)$ worst-case. \\
Complexity of inspecting root maximum? & $O(1)$. \\
Can arbitrary key be searched in $O(\log n)$? & No; worst-case arbitrary search is $O(n)$. \\
Time complexity of Heap Sort? & $O(n \log n)$ in all cases (best, average, worst). \\
Is Heap Sort stable? & No, it does not preserve relative order of duplicates. \\
Is a max-heap a Binary Search Tree? & No; heap enforces vertical order, BST enforces horizontal order. \\
Why does bottom-up build take $O(n)$? & Most nodes reside near the leaves and travel very small distances. \\
\bottomrule
\end{tabularx}
\caption{Rapid-fire recruitment examination questions and answers.}
\end{table}

\section{Final Quick-Revision Sheet: Linked Lists and Heaps}

\subsection{Linked Lists Quick Revision}
\begin{itemize}
    \item \textbf{Singly Linked List:} Unidirectional (\texttt{next}); $O(1)$ insert at head; $O(n)$ arbitrary search/delete.
    \item \textbf{Doubly Linked List:} Bidirectional (\texttt{prev}, \texttt{next}); $O(1)$ delete of known node.
    \item \textbf{Circular Linked List:} Last node links to first; $O(1)$ front and end insertion with \texttt{tail} pointer.
    \item Traversal loop for circular list requires \texttt{do \{ ... \} while (curr != head);}.
\end{itemize}

\subsection{Max-Heaps Quick Revision}
\begin{itemize}
    \item Complete binary tree stored in array without pointers.
    \item Parent $\geq$ Children; Maximum at $A[0]$.
    \item 0-based formulas: $\text{left} = 2i + 1$, $\text{right} = 2i + 2$, $\text{parent} = \lfloor(i-1)/2\rfloor$.
    \item Build heap bottom-up: $O(n)$.
    \item Insert: append and sift up ($O(\log n)$).
    \item Delete Max: replace root with last element and sift down ($O(\log n)$).
    \item Get Max: $O(1)$.
    \item Heap Sort: $O(n \log n)$ time, $O(1)$ auxiliary space, unstable.
\end{itemize}

\begin{quote}
\textbf{Final Architectural Takeaway:} Linked lists optimize dynamic node insertion and deletion without shifting contiguous memory. Binary heaps leverage the structural completeness of binary trees to deliver $O(1)$ extremum access and $O(\log n)$ priority updates with zero pointer memory overhead.
\end{quote}


\begin{center}
\vspace{1em}
{\color{primary}\rule{0.5\textwidth}{1.5pt}}\\[6pt]
{\Large\bfseries\color{primary} Part III: Circular Queues and Priority Queues}\\[3pt]
{\color{primary}\rule{0.5\textwidth}{1.5pt}}
\vspace{1em}
\end{center}

This chapter covers two fundamental queue variants: the \textbf{circular queue}, which eliminates memory wastage by wrapping array indices, and the \textbf{priority queue}, which processes elements according to importance or urgency rather than strictly following their arrival order.

\begin{center}
\vspace{0.5em}
{\color{primary}\rule{0.45\textwidth}{1.2pt}}\\[4pt]
{\Large\bfseries\color{primary} Section I: Circular Queue}\\[2pt]
{\color{primary}\rule{0.45\textwidth}{1.2pt}}
\vspace{0.5em}
\end{center}

\section{What is a Circular Queue?}

A \textbf{circular queue} is an array-based queue in which the end of the underlying memory buffer is logically wrapped around to connect with the beginning. When the rear pointer reaches the final array index $N-1$, the next insertion wraps back to index $0$, provided that earlier positions have been vacated by dequeue operations.

Like a standard queue, a circular queue strictly honors the \textbf{FIFO (First In, First Out)} principle. The distinction lies entirely in how physical memory slots are recycled.

A circular queue is also called a \textbf{ring buffer} when utilized as a fixed-capacity streaming or buffering data structure.

\subsection{Practical Applications}
\begin{itemize}
    \item \textbf{Network Packet Buffers:} Network routers and switches buffer incoming packets in ring buffers awaiting processing by network processors.
    \item \textbf{Keyboard Typeahead Buffers:} Operating systems capture hardware keystrokes in circular buffers until the active process reads them.
    \item \textbf{Producer-Consumer Architectures:} Thread-safe circular queues decouple producer threads from consumer threads in concurrent systems.
    \item \textbf{Audio and Video Streaming:} Multimedia players buffer sound/video frames in circular audio rings to prevent playback jitter.
    \item \textbf{Embedded Systems:} Embedded microcontrollers with strict, bounded RAM use fixed ring buffers for serial (UART/SPI) communication.
    \item \textbf{Operating System Schedulers:} Round-robin process ready queues cycle through processes efficiently.
\end{itemize}

\section{Why is a Circular Queue Needed?}

To appreciate the necessity of circular wrapping, examine the structural flaws of a naive linear array queue.

Suppose an array has five slots ($N=5$):
\begin{lstlisting}[style=pseudocode]
Index:  0    1    2    3    4
Data:  [10,  20,  30,  40,  50]
        ^                    ^
      front                 rear
\end{lstlisting}
Here, \texttt{front} points to index $0$, and \texttt{rear} points to index $4$. The queue is full.

Now, perform two dequeue operations, removing $10$ and $20$. The queue logically contains:
\[
30, \; 40, \; 50
\]
The state of the array is now:
\begin{lstlisting}[style=pseudocode]
Index:  0    1    2    3    4
Data:  [--,  --,  30,  40,  50]
                  ^          ^
                front       rear
\end{lstlisting}
Indices $0$ and $1$ are now vacant and completely unused. However, in a simple linear implementation, \texttt{rear} is already located at index $N-1$ ($4$). If another element is to be enqueued, the algorithm inspects \texttt{rear == N - 1} and reports an overflow condition!

This failure is termed \textbf{false overflow} (or space wastage in a linear queue): the queue rejects new elements despite having unoccupied slots available.

\begin{quote}
\textbf{The Circular Solution:} A circular queue resolves false overflow by wrapping \texttt{rear} around to index $0$, reusing vacated memory without shifting existing elements.
\end{quote}

\textbf{Important Qualification:} Circular wrapping does not confer infinite capacity. A fixed-capacity circular queue still experiences genuine overflow when all $N$ slots are occupied simultaneously.

\section{Front and Rear Indexing in a Circular Queue}

A circular queue manages two primary index pointers:
\begin{itemize}
    \item \textbf{Front:} The index of the next element awaiting extraction (read/dequeue).
    \item \textbf{Rear:} The index of the most recently inserted element.
\end{itemize}

For an array of capacity $N$, indexed $0$ through $N-1$, pointer advancement uses the \textbf{modulo operator}:
\[
\boxed{\text{Next Index} = (\text{Current Index} + 1) \bmod N}
\]

For example, when $N = 5$:
\begin{align*}
(4 + 1) \bmod 5 &= 0 \quad (\text{wraps from last index to start}) \\
(0 + 1) \bmod 5 &= 1 \\
(2 + 1) \bmod 5 &= 3
\end{align*}
Modulo arithmetic guarantees that indices remain strictly within $[0, N-1]$ across indefinite cycles.

\section{Circular Queue Operations}

Two primary implementation models exist:
\begin{enumerate}
    \item \textbf{Convention A (Pointer-based with $-1$ for empty):} Tracks \texttt{front} and \texttt{rear}, initializing both to $-1$.
    \item \textbf{Convention B (Count-based):} Explicitly maintains a \texttt{count} variable alongside \texttt{front} and \texttt{rear}.
\end{enumerate}

\subsection{Method A: Front and Rear with -1 Initialization}

Let capacity be $N$. Initially:
\[
\texttt{front} = -1, \quad \texttt{rear} = -1
\]

\subsubsection{A. Enqueue Operation}
To insert element $x$:
\begin{enumerate}
    \item \textbf{Check for Queue Overflow:} The queue is full if the next position after \texttt{rear} equals \texttt{front}:
    \[
    (\texttt{rear} + 1) \bmod N == \texttt{front}
    \]
    If true, report overflow and abort.
    \item \textbf{Handle Empty State:} If $\texttt{front} == -1$ (queue is empty), initialize both:
    \[
    \texttt{front} = 0, \quad \texttt{rear} = 0
    \]
    \item \textbf{Advance Rear Circularly:} If not empty:
    \[
    \texttt{rear} = (\texttt{rear} + 1) \bmod N
    \]
    \item \textbf{Store Data:} Set $\texttt{array}[\texttt{rear}] = x$.
\end{enumerate}

\subsubsection{B. Dequeue Operation}
To remove the front element:
\begin{enumerate}
    \item \textbf{Check for Queue Underflow:} If $\texttt{front} == -1$, report underflow and abort.
    \item \textbf{Retrieve Data:} Read $x = \texttt{array}[\texttt{front}]$.
    \item \textbf{Check for Single-Element Queue:} If $\texttt{front} == \texttt{rear}$, the queue contained exactly one element. Reset both pointers to empty:
    \[
    \texttt{front} = -1, \quad \texttt{rear} = -1
    \]
    \item \textbf{Advance Front Circularly:} Otherwise:
    \[
    \texttt{front} = (\texttt{front} + 1) \bmod N
    \]
    \item Return $x$.
\end{enumerate}

\subsubsection{C. Peek / Front}
If $\texttt{front} \neq -1$, return $\texttt{array}[\texttt{front}]$ without modifying indices.

\subsection{Step-by-Step Execution Trace}

Consider a circular queue with $N = 5$, initially empty ($\texttt{front} = -1, \texttt{rear} = -1$):

\begin{table}[htbp]
\centering
\small
\begin{tabularx}{\textwidth}{l c c X l}
\toprule
\textbf{Operation} & \textbf{Front} & \textbf{Rear} & \textbf{Array Layout ($0 \dots 4$)} & \textbf{Status / Result} \\
\midrule
Initial        & $-1$ & $-1$ & $[\;,\;,\;,\;,\;]$               & Empty \\
\texttt{enqueue(10)} & $0$  & $0$  & $[10, \;,\;,\;,\;]$              & First element \\
\texttt{enqueue(20)} & $0$  & $1$  & $[10, 20, \;,\;,\;]$             & Enqueued \\
\texttt{enqueue(30)} & $0$  & $2$  & $[10, 20, 30, \;,\;]$            & Enqueued \\
\texttt{enqueue(40)} & $0$  & $3$  & $[10, 20, 30, 40, \;]$           & Enqueued \\
\texttt{dequeue()}   & $1$  & $3$  & $[--, 20, 30, 40, \;]$           & Returns $10$ \\
\texttt{dequeue()}   & $2$  & $3$  & $[--, --, 30, 40, \;]$           & Returns $20$ \\
\texttt{enqueue(50)} & $2$  & $4$  & $[--, --, 30, 40, 50]$           & Enqueued at index 4 \\
\texttt{enqueue(60)} & $2$  & $0$  & $[60, --, 30, 40, 50]$           & \textbf{Wraps to index 0!} \\
\texttt{enqueue(70)} & $2$  & $1$  & $[60, 70, 30, 40, 50]$           & \textbf{Wraps to index 1; Full!} \\
\bottomrule
\end{tabularx}
\caption{Detailed step-by-step state trace of circular queue pointer wrapping.}
\end{table}

\textbf{Analysis:} At the final step, the elements in FIFO order are:
\[
30 \longrightarrow 40 \longrightarrow 50 \longrightarrow 60 \longrightarrow 70
\]
Although physically stored at indices $\{2, 3, 4, 0, 1\}$, reading circularly from $\texttt{front} = 2$ to $\texttt{rear} = 1$ extracts them in precise chronological order. The condition $(1 + 1) \bmod 5 = 2 == \texttt{front}$ confirms that the queue is full.

\section{Circular Queue Implementation in C++}

The production-grade implementation below maintains an explicit \texttt{count} variable. This elegantly eliminates ambiguities between full and empty states when indices match.

\begin{lstlisting}[style=cppstyle]
#include <iostream>
#include <vector>
using namespace std;

class CircularQueue {
private:
    vector<int> data;
    size_t frontIndex = 0;
    size_t rearIndex = 0;   // Points to the NEXT available slot
    size_t count = 0;

public:
    explicit CircularQueue(size_t capacity)
        : data(capacity) {}

    bool isEmpty() const {
        return count == 0;
    }

    bool isFull() const {
        return count == data.size();
    }

    size_t size() const {
        return count;
    }

    // Insert an element at the rear: O(1)
    bool enqueue(int value) {
        if (isFull()) {
            return false; // Overflow
        }
        data[rearIndex] = value;
        rearIndex = (rearIndex + 1) % data.size();
        ++count;
        return true;
    }

    // Remove the front element: O(1)
    bool dequeue(int& value) {
        if (isEmpty()) {
            return false; // Underflow
        }
        value = data[frontIndex];
        frontIndex = (frontIndex + 1) % data.size();
        --count;
        return true;
    }

    // Read the front element without removing it: O(1)
    bool front(int& value) const {
        if (isEmpty()) {
            return false;
        }
        value = data[frontIndex];
        return true;
    }
};

int main() {
    CircularQueue q(5);
    int value;

    q.enqueue(10);
    q.enqueue(20);
    q.enqueue(30);

    if (q.dequeue(value)) cout << "Removed: " << value << '\n';

    q.enqueue(40);
    q.enqueue(50);
    q.enqueue(60);

    cout << "Queue size: " << q.size() << '\n';

    if (q.front(value)) cout << "Front: " << value << '\n';

    if (!q.enqueue(70)) cout << "Queue is full\n";

    while (q.dequeue(value)) {
        cout << "Removed: " << value << '\n';
    }

    if (!q.dequeue(value)) cout << "Queue is empty\n";

    return 0;
}
\end{lstlisting}

\subsubsection{Expected Output}
\begin{lstlisting}[style=pseudocode]
Removed: 10
Queue size: 5
Front: 20
Queue is full
Removed: 20
Removed: 30
Removed: 40
Removed: 50
Removed: 60
Queue is empty
\end{lstlisting}

\subsection{Complexity and Advantages}

\begin{table}[htbp]
\centering
\small
\begin{tabularx}{0.75\textwidth}{l X}
\toprule
\textbf{Operation} & \textbf{Time / Space Complexity} \\
\midrule
Enqueue      & $O(1)$ worst-case \\
Dequeue      & $O(1)$ worst-case \\
Front / Peek & $O(1)$ \\
isEmpty      & $O(1)$ \\
isFull       & $O(1)$ \\
Auxiliary Space & $O(N)$ fixed pre-allocated buffer \\
\bottomrule
\end{tabularx}
\caption{Operational complexity profile of Circular Queue.}
\end{table}

\textbf{Key Advantages:}
\begin{itemize}
    \item Reuses vacated memory automatically without element shifting ($O(1)$ dequeue).
    \item Fixed memory footprint suitable for mission-critical and embedded real-time systems.
    \item Eliminates memory fragmentation caused by dynamic node allocations.
\end{itemize}

\begin{center}
\vspace{1em}
{\color{primary}\rule{0.45\textwidth}{1.2pt}}\\[4pt]
{\Large\bfseries\color{primary} Section II: Priority Queue}\\[2pt]
{\color{primary}\rule{0.45\textwidth}{1.2pt}}
\vspace{0.5em}
\end{center}

\section{What is a Priority Queue?}

A \textbf{priority queue} is an Abstract Data Type (ADT) in which each stored element is associated with a priority value. Elements are extracted according to priority order rather than strictly following their arrival sequence.

Unlike a FIFO queue, the earliest arriving element is not necessarily served first. A critical task arriving late can preempt older, lower-priority tasks.

\subsection{Real-Life Analogy: Hospital Emergency Triage}
Consider hospital emergency patients arriving at a clinic:
\begin{table}[htbp]
\centering
\small
\begin{tabularx}{0.85\textwidth}{l c X}
\toprule
\textbf{Patient} & \textbf{Priority Level} & \textbf{Condition} \\
\midrule
Patient A & 2 & Mild fever (arrived 10:00 AM) \\
Patient B & 5 & Cardiac arrest (arrived 10:15 AM) \\
Patient C & 3 & Bone fracture (arrived 10:05 AM) \\
\bottomrule
\end{tabularx}
\caption{Emergency triage priority comparison.}
\end{table}
Although Patient A arrived first, Patient B has the highest medical urgency (priority 5) and is attended to immediately. This encapsulates priority queue semantics.

\subsection{Key Terminology}
\begin{itemize}
    \item \textbf{Max-Priority Queue:} Extracts the element with the largest priority key first.
    \item \textbf{Min-Priority Queue:} Extracts the element with the smallest priority key first.
    \item \textbf{Priority Key:} The metric (e.g., urgency score, shortest distance, timestamp) determining extraction precedence.
    \item \textbf{Stable Priority Queue:} An implementation that guarantees FIFO ordering among items with identical priorities.
\end{itemize}

\section{Types of Priority Queues}

\subsection{Max-Priority Queue vs Min-Priority Queue}

\begin{itemize}
    \item \textbf{Max-Priority Queue:} Given inputs $\{10, 40, 20, 30\}$, extractions proceed:
    \[
    40 \longrightarrow 30 \longrightarrow 20 \longrightarrow 10
    \]
    This is the default configuration of C++ STL's \texttt{std::priority\_queue<int>}.
    \item \textbf{Min-Priority Queue:} For the same inputs, extractions proceed:
    \[
    10 \longrightarrow 20 \longrightarrow 30 \longrightarrow 40
    \]
    Indispensable for shortest-path algorithms (Dijkstra) and Minimum Spanning Trees (Prim).
\end{itemize}

\subsection{Does a Priority Queue Preserve FIFO Order for Equal Priorities?}

\textbf{Not by default.} Standard priority queue specifications do not mandate stability. If Task 1, Task 2, and Task 3 all share priority 5, an ordinary heap may extract them in any arbitrary order. If strict FIFO tie-breaking is required, a sequence number or secondary timestamp must be incorporated into the comparator.

\section{Priority Queue Implementation Strategies}

A priority queue can be realized using multiple underlying data structures, each presenting distinct trade-offs:

\begin{table}[htbp]
\centering
\small
\begin{tabularx}{\textwidth}{l X X X}
\toprule
\textbf{Implementation} & \textbf{Insert} & \textbf{Peek Highest} & \textbf{Extract Highest} \\
\midrule
Unsorted Array          & $O(1)$ amortized & $O(n)$ linear scan    & $O(n)$ search and shift \\
Unsorted Linked List    & $O(1)$           & $O(n)$                & $O(n)$ search and delete \\
Sorted Array            & $O(n)$ insertion & $O(1)$                & $O(1)$ pop back \\
Sorted Linked List      & $O(n)$ traversal & $O(1)$                & $O(1)$ pop head \\
\textbf{Binary Heap}    & $\mathbf{O(\log n)}$ & $\mathbf{O(1)}$   & $\mathbf{O(\log n)}$ \\
Balanced BST (AVL/RB)   & $O(\log n)$      & $O(\log n)$ (or $O(1)$) & $O(\log n)$ \\
\bottomrule
\end{tabularx}
\caption{Comparison of underlying priority queue data structures.}
\end{table}

\textbf{Why the Binary Heap is Universally Preferred:}
A binary heap strikes the optimal practical balance: both insertions and extractions run in guaranteed $O(\log n)$ time, inspection is instantaneous ($O(1)$), and memory storage requires zero explicit pointer overhead because it maps into a flat 1D array.

\section{Implementing a Priority Queue Using a Binary Heap}

A binary heap is a complete binary tree obeying the heap-order property:
\begin{itemize}
    \item \textbf{Insert:} Append element at end of heap array, then \textbf{sift-up} ($O(\log n)$).
    \item \textbf{Extract (Max/Min):} Swap root with last array element, reduce size, and \textbf{sift-down} ($O(\log n)$).
    \item \textbf{Peek:} Read root element $A[0]$ ($O(1)$).
\end{itemize}

\section{C++ Priority Queue Using the Standard Library}

C++ provides \texttt{std::priority\_queue} in the \texttt{<queue>} header.

\subsection{Standard Max-Priority Queue}
\begin{lstlisting}[style=cppstyle]
#include <iostream>
#include <queue>
using namespace std;

int main() {
    priority_queue<int> pq; // Max-heap by default

    pq.push(10);
    pq.push(40);
    pq.push(20);
    pq.push(30);

    cout << "Highest priority: " << pq.top() << '\n'; // 40

    while (!pq.empty()) {
        cout << pq.top() << " ";
        pq.pop();
    }
    cout << '\n'; // Prints: 40 30 20 10
    return 0;
}
\end{lstlisting}

\subsection{Standard Min-Priority Queue}
To construct a min-heap, pass \texttt{greater<T>} as the comparator:
\begin{lstlisting}[style=cppstyle]
#include <iostream>
#include <queue>
#include <vector>
#include <functional>
using namespace std;

int main() {
    // Min-Priority Queue declaration
    priority_queue<int, vector<int>, greater<int>> minPQ;

    minPQ.push(10);
    minPQ.push(40);
    minPQ.push(20);
    minPQ.push(30);

    while (!minPQ.empty()) {
        cout << minPQ.top() << " ";
        minPQ.pop();
    }
    cout << '\n'; // Prints: 10 20 30 40
    return 0;
}
\end{lstlisting}

\begin{table}[htbp]
\centering
\small
\begin{tabularx}{\textwidth}{l p{7cm} X}
\toprule
\textbf{Method} & \textbf{Description} & \textbf{Complexity} \\
\midrule
\texttt{push(x)} & Inserts element $x$ into the heap & $O(\log n)$ \\
\texttt{emplace(args)} & In-place constructs and inserts element & $O(\log n)$ \\
\texttt{top()} & Inspects highest-priority element & $O(1)$ \\
\texttt{pop()} & Removes highest-priority element & $O(\log n)$ \\
\texttt{empty()} & Checks if container is empty & $O(1)$ \\
\texttt{size()} & Returns count of active elements & $O(1)$ \\
\bottomrule
\end{tabularx}
\caption{C++ STL priority\_queue member functions and complexities.}
\end{table}

\begin{quote}
\textbf{Crucial API Detail:} In C++, \texttt{pop()} returns \texttt{void}. To read the extracted value, you must inspect \texttt{top()} \emph{before} calling \texttt{pop()}. Invoking \texttt{top()} or \texttt{pop()} on an empty container causes undefined behavior.
\end{quote}

\section{Priority Queue with Real-World Tasks and Custom Comparators}

In production software, queues manage complex objects (jobs, network requests, database alerts) with custom priorities.

\subsection{Stable Priority Queue Implementation}

The following C++ program demonstrates a custom struct \texttt{Task} prioritized by urgency. A secondary \texttt{sequence} counter enforces strict FIFO arrival order whenever priorities collide.

\begin{lstlisting}[style=cppstyle]
#include <iostream>
#include <queue>
#include <string>
#include <vector>
using namespace std;

struct Task {
    string name;
    int priority;
    unsigned long long sequence; // Arrival counter for stability
};

struct CompareTask {
    bool operator()(const Task& a, const Task& b) const {
        // Higher numeric priority value goes first
        if (a.priority != b.priority) {
            return a.priority < b.priority;
        }
        // For identical priorities, smaller sequence number (earlier) goes first
        return a.sequence > b.sequence;
    }
};

int main() {
    priority_queue<Task, vector<Task>, CompareTask> tasks;
    unsigned long long nextSeq = 0;

    tasks.push({"Database backup", 2, nextSeq++});
    tasks.push({"Critical server alert", 5, nextSeq++});
    tasks.push({"Routine report", 1, nextSeq++});
    tasks.push({"Transaction recovery", 4, nextSeq++});
    tasks.push({"Second critical alert", 5, nextSeq++});

    while (!tasks.empty()) {
        Task current = tasks.top();
        tasks.pop();
        cout << current.name << " | Priority: " << current.priority << '\n';
    }
    return 0;
}
\end{lstlisting}

\subsubsection{Expected Output}
\begin{lstlisting}[style=pseudocode]
Critical server alert | Priority: 5
Second critical alert | Priority: 5
Transaction recovery | Priority: 4
Database backup | Priority: 2
Routine report | Priority: 1
\end{lstlisting}
Notice that the two priority-5 tasks depart in exact insertion order.

\section{Priority Queue Applications and Important Concepts}

\subsection{Operating System Scheduling and Starvation}
Priority-based schedulers select highest-priority threads first. However, continuous arrival of high-priority jobs can cause \textbf{starvation} (indefinite postponement) for low-priority tasks.
\begin{itemize}
    \item \textbf{Aging Technique:} Operating systems counter starvation by periodically incrementing the priority of tasks waiting in the queue. Over time, an aging low-priority task eventually becomes the highest-priority task.
\end{itemize}

\subsection{Dijkstra's Shortest Path Algorithm}
Dijkstra's algorithm repeatedly extracts the unvisited vertex with the smallest tentative distance. Using a \textbf{min-priority queue} reduces the running time from $O(V^2)$ to $O((V + E) \log V)$.

\subsection{Huffman Coding Tree Construction}
Huffman compression uses a min-priority queue to repeatedly extract the two smallest frequency symbols, merge them into an internal tree node, and push the combined sum back into the queue.

\subsection{Top-$K$ Problem in Streaming Big Data}
To extract the $k$ largest elements from $n$ continuous numbers ($n \gg k$):
\begin{enumerate}
    \item Maintain a \textbf{min-heap of size $k$}.
    \item For every incoming element $x$: if $x > \text{heap.top()}$, pop the root and insert $x$.
    \item The min-heap perpetually retains the $k$ largest elements in $O(n \log k)$ time and $O(k)$ memory.
\end{enumerate}

\section{Circular Queue versus Priority Queue}

\begin{table}[htbp]
\centering
\small
\begin{tabularx}{\textwidth}{l X X}
\toprule
\textbf{Basis} & \textbf{Circular Queue} & \textbf{Priority Queue} \\
\midrule
Primary Purpose       & Recycles array storage while preserving arrival order & Selects elements based on priority metrics \\
Processing Order      & Strict FIFO (First In, First Out)                    & Urgency / Priority (Highest or Lowest first) \\
Underlying Structure  & Fixed-size circular 1D array                         & Binary heap, balanced BST, or sorted array \\
Enqueue / Insert      & $O(1)$ worst-case                                    & $O(\log n)$ with binary heap \\
Dequeue / Extract     & $O(1)$ worst-case                                    & $O(\log n)$ with binary heap \\
Front / Peek Access   & Oldest waiting element                               & Highest-priority element \\
Buffer Capacity       & Fixed (pre-allocated array)                          & Bounded or dynamically resizable \\
Typical Application   & Network buffers, audio streams, UART rings           & CPU scheduling, Dijkstra, Huffman, event queues \\
\bottomrule
\end{tabularx}
\caption{Architectural comparison between Circular Queue and Priority Queue.}
\end{table}

\section{Common Examination Mistakes and Traps}

\begin{enumerate}
    \item \textbf{Circular Queue vs Circular Linked List:} A circular queue is a FIFO buffer (typically implemented with a circular array). A circular linked list is a pointer-based node cycle.
    \item \textbf{Forgetting Modulo Wrapping:} Pointer advancement requires $(\text{index} + 1) \bmod N$, not simple incrementation.
    \item \textbf{Mixing Front/Rear Indexing Conventions:} In some textbooks, \texttt{rear} holds the last inserted element; in others, \texttt{rear} points to the next available slot. Full/empty conditions must match the convention used.
    \item \textbf{Assuming Priority Queues are FIFO:} A priority queue is inherently \emph{not} FIFO. Duplicate priorities require explicit secondary sequence keys if stability is needed.
    \item \textbf{Confusing Priority Queue with Binary Heap:} A priority queue is an abstract data type (ADT); a binary heap is a concrete implementation.
    \item \textbf{C++ STL Default Comparator Direction:} \texttt{std::priority\_queue<int>} is a max-heap. Creating a min-heap requires \texttt{greater<int>}.
\end{enumerate}

\section{Likely Recruitment Examination Questions: Circular and Priority Queues}

\begin{table}[htbp]
\centering
\small
\begin{tabularx}{\textwidth}{p{7cm} X}
\toprule
\textbf{Question} & \textbf{Key Answer} \\
\midrule
What is a circular queue? & A FIFO queue where the last array slot wraps to the first via modulo arithmetic. \\
Why is it better than a linear array queue? & It eliminates false overflow by recycling vacated slots without shifting. \\
What is the index wrapping formula? & $(\text{index} + 1) \bmod N$. \\
Overflow condition in conventional circular queue? & $(\texttt{rear} + 1) \bmod N == \texttt{front}$. \\
What is a priority queue? & An ADT that serves elements based on priority rather than arrival order. \\
Difference between FIFO and priority processing? & FIFO follows chronological arrival; priority queue follows importance/urgency. \\
What does a max-priority queue return first? & The element possessing the maximum priority key. \\
What does a min-priority queue return first? & The element possessing the minimum priority key. \\
Standard implementation of a priority queue? & A binary heap (complete binary tree in array). \\
Insertion complexity in binary-heap priority queue? & $O(\log n)$ worst-case. \\
Extraction complexity from max-heap? & $O(\log n)$ worst-case. \\
Complexity of inspecting highest-priority key? & $O(1)$ at root index. \\
Does C++ priority\_queue preserve arrival order? & No, stability is not guaranteed for equal priorities. \\
How to prevent starvation in priority scheduling? & Use the aging technique to progressively increase waiting tasks' priority. \\
Which priority queue is used in Dijkstra's algorithm? & Min-priority queue (extracts vertex with smallest tentative distance). \\
Difference between priority queue and heap? & Priority queue is the interface (ADT); binary heap is the data structure. \\
\bottomrule
\end{tabularx}
\caption{Rapid-fire recruitment examination questions and answers.}
\end{table}

\section{Final Quick-Revision Sheet: Circular and Priority Queues}

\subsection{Circular Queue Summary}
\begin{itemize}
    \item \textbf{FIFO Semantics:} Strictly preserves chronological arrival order.
    \item \textbf{Index Advancement:} $\text{next} = (i + 1) \bmod N$.
    \item \textbf{Empty Condition:} $\texttt{front} == \texttt{rear} == -1$ (or $\texttt{count} == 0$).
    \item \textbf{Full Condition:} $(\texttt{rear} + 1) \bmod N == \texttt{front}$ (or $\texttt{count} == N$).
    \item \textbf{Complexities:} Enqueue: $O(1)$; Dequeue: $O(1)$; Peek: $O(1)$; Space: $O(N)$ static buffer.
\end{itemize}

\subsection{Priority Queue Summary}
\begin{itemize}
    \item \textbf{Priority Semantics:} Elements dequeued by priority key, not arrival chronology.
    \item \textbf{Max-PQ:} Highest key first; \textbf{Min-PQ:} Lowest key first.
    \item \textbf{Binary Heap Complexity:} Insert: $O(\log n)$; Extract: $O(\log n)$; Peek: $O(1)$; Build Heap: $O(n)$.
    \item \textbf{C++ Declarations:}
    \begin{itemize}
        \item Max-heap: \texttt{priority\_queue<int> pq;}
        \item Min-heap: \texttt{priority\_queue<int, vector<int>, greater<int>> minPQ;}
    \end{itemize}
    \item \textbf{Stability:} Requires explicit tie-breaking sequence counter.
\end{itemize}

\begin{quote}
\textbf{Final Architectural Takeaway:} A circular queue optimizes storage reuse while strictly maintaining FIFO chronological arrival order. A priority queue optimizes dynamic selection according to urgency or importance, leveraging binary heaps for logarithmic operational efficiency.
\end{quote}


\begin{center}
\vspace{1em}
{\color{primary}\rule{0.5\textwidth}{1.5pt}}\\[6pt]
{\Large\bfseries\color{primary} Part IV: Trees, Binary Trees, and B+ Trees}\\[3pt]
{\color{primary}\rule{0.5\textwidth}{1.5pt}}
\vspace{1em}
\end{center}

This chapter covers tree terminology, fundamental tree properties, binary trees, binary search trees, balanced trees, heaps, B-trees, B+ trees, and the mathematical calculations commonly encountered in technical recruitment examinations.

Special attention is given to \textbf{calculating the maximum and minimum number of keys in a B+ tree from its order and height}, as the formulation strictly depends on the architectural definitions employed in the question.

\begin{center}
\vspace{0.5em}
{\color{primary}\rule{0.45\textwidth}{1.2pt}}\\[4pt]
{\Large\bfseries\color{primary} Section I: Fundamentals of Trees}\\[2pt]
{\color{primary}\rule{0.45\textwidth}{1.2pt}}
\vspace{0.5em}
\end{center}

\section{What is a Tree?}

A \textbf{tree} is a hierarchical, non-linear data structure consisting of nodes connected by directed or undirected edges. In a rooted tree, exactly one node is designated as the \textbf{root}, and every other node has exactly one parent.

From a graph-theoretic perspective, a tree is a \textbf{connected, acyclic undirected graph}. This implies that between any two nodes there exists exactly one simple path, and no closed cycles exist.

Trees are indispensable when modeling hierarchical or branching relationships, such as file system folder structures, organizational hierarchies, database indexing (B/B+ trees), Document Object Models (DOM), and compiler abstract syntax trees (AST).

\subsection{Basic Terminology}

\begin{table}[htbp]
\centering
\small
\begin{tabularx}{\textwidth}{l p{7.5cm} X}
\toprule
\textbf{Term} & \textbf{Formal Definition} & \textbf{Example} \\
\midrule
Node             & An individual data element in the tree                      & A directory or file \\
Root             & The topmost node, possessing no parent                      & The root directory (\texttt{/}) \\
Edge             & A directed/undirected link connecting two nodes             & Parent-child relationship \\
Parent           & The immediate predecessor of a node                         & Directory containing a subdirectory \\
Child            & A node directly succeeding another node                     & Subdirectory inside a folder \\
Siblings         & Nodes sharing the identical parent                          & Files in the same directory \\
Leaf Node        & An external node having zero children                       & An empty file / leaf directory \\
Internal Node    & A non-leaf node possessing at least one child               & A folder containing items \\
Ancestor         & Any node along the upward path from a node to the root      & Parent, grandparent, etc. \\
Descendant       & Any node situated in the subtree below a given node         & Child, grandchild, etc. \\
Degree of a Node & Total number of children belonging to the node              & Node with 3 children has degree 3 \\
Degree of a Tree & The maximum degree of any node within the tree              & A binary tree has degree $\leq 2$ \\
Path             & A continuous sequence of connected nodes                    & Root-to-leaf traversal \\
Path Length      & Number of edges traversed along the path                    & 4 edges = path length 4 \\
Depth of a Node  & Number of edges from the root to that node                  & Root depth is $0$ \\
Height of a Node & Number of edges on the longest downward path to a leaf      & A leaf node has height $0$ \\
Height of a Tree & The height of its root node (maximum depth among all nodes) & Edge count on longest path \\
\bottomrule
\end{tabularx}
\caption{Fundamental tree terminology and definitions.}
\end{table}

\begin{quote}
\textbf{Standard Height Convention:} Throughout this text, \textbf{height is measured in edges}. Consequently, a single-node tree has height $0$, and a tree of height $h$ contains $h + 1$ levels (levels $0, 1, \dots, h$).
\end{quote}

\subsection{Fundamental Mathematical Properties of Trees}

\subsubsection{Property 1: Number of Edges}
A tree containing $n$ nodes has exactly:
\[
\boxed{E = n - 1}
\]
where $E$ is the total edge count.
\begin{itemize}
    \item $1$ node $\implies 0$ edges.
    \item $10$ nodes $\implies 9$ edges.
    \item $100$ nodes $\implies 99$ edges.
\end{itemize}
\textbf{Proof:} The root node has zero incoming edges. Each of the remaining $n - 1$ non-root nodes possesses exactly one incoming edge from its unique parent. Hence, total edges $= n - 1$.

\subsubsection{Property 2: Unique Simple Path}
Between any two arbitrary nodes in a tree, there exists \textbf{exactly one simple path}. If two distinct paths connected the same pair, their union would form a cycle, violating the acyclic definition.

\subsubsection{Property 3: Invariance of Edge Count to Topology}
A tree with $n$ nodes contains $n - 1$ edges \emph{regardless} of whether it is perfectly balanced or degenerate (skewed into a linked list). Structural topology governs tree height and query complexity, but never the edge count.

\subsubsection{Property 4: Edge Addition}
Adding an edge between any two existing nodes in a tree creates exactly one cycle.

\subsubsection{Property 5: Edge Removal}
Removing any single edge from a tree partitions it into exactly two disconnected components (subtrees).

\section{General Tree Mathematics}

For general trees where internal nodes can have up to $m$ children ($m$-ary trees):

\subsection{Maximum Nodes at Level $i$}
In a rooted tree of maximum degree $m$, the maximum node capacity at level $i$ (with root at level $0$) is:
\[
\boxed{N_i \leq m^i}
\]
Level $0$ has $m^0 = 1$ node; Level $1$ has $m^1 = m$; Level $2$ has $m^2$; Level $i$ has $m^i$.

\subsection{Maximum Total Nodes for Height $h$}
Summing the geometric progression from level $0$ to level $h$:
\[
N_{\max} = 1 + m + m^2 + \dots + m^h = \boxed{\frac{m^{h+1} - 1}{m - 1}} \quad (m > 1)
\]
For example, for a ternary tree ($m = 3$) of height $h = 3$:
\[
N_{\max} = \frac{3^{3+1} - 1}{3 - 1} = \frac{81 - 1}{2} = 40 \text{ nodes}
\]

\subsection{Minimum Possible Height for $n$ Nodes}
To minimize height, all levels must be packed maximally. Solving $n \leq \frac{m^{h+1}-1}{m-1}$ for $h$:
\[
\boxed{h_{\min} = \left\lceil \log_m\big(n(m - 1) + 1\big) \right\rceil - 1}
\]
For a binary tree ($m = 2$) with $n = 100$:
\[
h_{\min} = \lceil \log_2(101) \rceil - 1 = 7 - 1 = 6 \text{ edges}
\]

\subsection{Maximum Possible Height for $n$ Nodes}
Without structural balancing, a tree can degenerate into a linear chain:
\[
\boxed{h_{\max} = n - 1}
\]

\section{Traversing a Tree}

Tree traversal systematically visits every node in a defined sequence. The four standard traversal methods are:

\subsection{Preorder Traversal (Root $\to$ Left $\to$ Right)}
\begin{lstlisting}[style=pseudocode]
PREORDER(node):
    if node is not null:
        visit(node)
        PREORDER(node.left)
        PREORDER(node.right)
\end{lstlisting}
Used for tree cloning, serialization, and prefix expression evaluation.

\subsection{Inorder Traversal (Left $\to$ Root $\to$ Right)}
\begin{lstlisting}[style=pseudocode]
INORDER(node):
    if node is not null:
        INORDER(node.left)
        visit(node)
        INORDER(node.right)
\end{lstlisting}
\textbf{Key Rule:} In a Binary Search Tree (BST), an inorder traversal visits keys in \textbf{strictly ascending sorted order}.

\subsection{Postorder Traversal (Left $\to$ Right $\to$ Root)}
\begin{lstlisting}[style=pseudocode]
POSTORDER(node):
    if node is not null:
        POSTORDER(node.left)
        POSTORDER(node.right)
        visit(node)
\end{lstlisting}
Used for bottom-up node deallocation (deleting a tree) and postfix (Reverse Polish) expression evaluation.

\subsection{Level-Order Traversal (Breadth-First)}
Visits nodes level-by-level from left to right using a FIFO queue. Used in BFS, tree serialization, and finding shortest root-to-leaf paths.

\textbf{Traversal Complexity:} All four traversals visit every node exactly once:
\[
\text{Time Complexity} = O(n)
\]
Auxiliary space is $O(h)$ for recursive call stacks ($O(n)$ worst-case) and $O(n)$ for level-order queue storage.

\begin{center}
\vspace{1em}
{\color{primary}\rule{0.45\textwidth}{1.2pt}}\\[4pt]
{\Large\bfseries\color{primary} Section II: Binary Trees}\\[2pt]
{\color{primary}\rule{0.45\textwidth}{1.2pt}}
\vspace{0.5em}
\end{center}

\section{What Is a Binary Tree?}

A \textbf{binary tree} is a rooted tree where every node has \textbf{at most two children}, distinguished as the \textbf{left child} and \textbf{right child}. A left child is distinct from a right child even if a parent has only one child.

\subsection{Types of Binary Trees}

\begin{enumerate}
    \item \textbf{Full Binary Tree (Strict / Proper):} Every node has either $0$ or exactly $2$ children. No node has $1$ child.
    \[
    L = I + 1, \quad n = 2I + 1, \quad I = \frac{n - 1}{2}, \quad L = \frac{n + 1}{2}
    \]
    A full binary tree \textbf{must always possess an odd number of nodes}.
    \item \textbf{Complete Binary Tree:} All levels are completely full except possibly the last level, which is filled strictly from left to right without gaps. Used in binary heaps.
    \item \textbf{Perfect Binary Tree:} All internal nodes have $2$ children and all leaves reside at the identical depth. For height $h$:
    \[
    n = 2^{h+1} - 1, \quad L = 2^h, \quad I = 2^h - 1
    \]
    \item \textbf{Balanced Binary Tree:} A binary tree whose height is bounded to $O(\log n)$ to guarantee efficient operations (e.g., AVL, Red-Black).
    \item \textbf{Skewed (Degenerate) Binary Tree:} Every internal node has exactly $1$ child, forming a linear list of height $h = n - 1$.
\end{enumerate}

\subsection{Full vs Complete vs Perfect Comparison}

\begin{table}[htbp]
\centering
\small
\begin{tabularx}{\textwidth}{l X X X}
\toprule
\textbf{Property} & \textbf{Full Binary Tree} & \textbf{Complete Binary Tree} & \textbf{Perfect Binary Tree} \\
\midrule
Internal node degree    & Exactly 2 children        & May have 1 child (at most one) & Exactly 2 children \\
All levels filled       & Not necessarily           & All except possibly last      & Completely filled \\
Last level layout       & No restriction            & Left-to-right without gaps    & Full left-to-right \\
All leaves equal depth  & Not necessarily           & Not necessarily               & Yes \\
Total node parity       & Always odd ($2I+1$)       & Odd or Even                   & Always odd ($2^{h+1}-1$) \\
Node count at height $h$& Varies                    & $2^h \dots 2^{h+1}-1$         & Exactly $2^{h+1}-1$ \\
\bottomrule
\end{tabularx}
\caption{Detailed comparison between Full, Complete, and Perfect binary trees.}
\end{table}

\section{Important Binary Tree Formulas}

Assuming height $h$ is measured in edges and root is at level $0$:

\subsection{Fundamental Extremum Bounds}
\begin{itemize}
    \item \textbf{Max nodes at level $i$:} $N_i^{\max} = 2^i$.
    \item \textbf{Max total nodes for height $h$:} $N_{\max} = 2^{h+1} - 1$.
    \item \textbf{Max leaves for height $h$:} $L_{\max} = 2^h$.
    \item \textbf{Min nodes for height $h$:} $N_{\min} = h + 1$ (skewed tree).
    \item \textbf{Min height for $n$ nodes:} $h_{\min} = \lceil \log_2(n + 1) \rceil - 1$.
\end{itemize}

\subsection{Theorem: Leaves vs Two-Child Nodes ($n_0 = n_2 + 1$)}
In \textbf{any} nonempty binary tree (not just full binary trees), let $n_0$ be leaves, $n_1$ be single-child nodes, and $n_2$ be two-child nodes:
\[
\boxed{n_0 = n_2 + 1}
\]
\textbf{Derivation:}
The total number of nodes is $n = n_0 + n_1 + n_2$.
The total number of edges is $E = n - 1 = n_1 + 2n_2$.
Substituting $n$:
\[
(n_0 + n_1 + n_2) - 1 = n_1 + 2n_2 \implies n_0 - 1 = n_2 \implies \boxed{n_0 = n_2 + 1}
\]

\subsection{Number of Null Child Pointers}
A binary tree with $n$ nodes has $2n$ total pointer fields. Exactly $n - 1$ fields point to child nodes (edges). Hence:
\[
\boxed{\text{Null Pointers} = 2n - (n - 1) = n + 1}
\]

\section{Array Representation of a Binary Tree}

\subsection{Zero-Based Indexing Formulas}
Root at index $0$:
\[
\text{Left}(i) = 2i + 1, \quad \text{Right}(i) = 2i + 2, \quad \text{Parent}(i) = \left\lfloor \frac{i - 1}{2} \right\rfloor
\]

\subsection{One-Based Indexing Formulas}
Root at position $1$:
\[
\text{Left}(i) = 2i, \quad \text{Right}(i) = 2i + 1, \quad \text{Parent}(i) = \left\lfloor \frac{i}{2} \right\rfloor
\]

\section{Number of Structurally Different Binary Trees}

For $n$ unlabeled nodes, the number of distinct binary tree shapes is given by the \textbf{Catalan Number} $C_n$:
\[
\boxed{C_n = \frac{1}{n+1} \binom{2n}{n} = \frac{(2n)!}{(n+1)! \, n!}}
\]
Values: $C_0=1, \; C_1=1, \; C_2=2, \; C_3=5, \; C_4=14, \; C_5=42, \; C_6=132$.
This also equals the number of distinct Binary Search Trees (BSTs) that can be formed from $n$ distinct keys.

\begin{center}
\vspace{1em}
{\color{primary}\rule{0.45\textwidth}{1.2pt}}\\[4pt]
{\Large\bfseries\color{primary} Section III: Binary Search Trees \& Balanced Trees}\\[2pt]
{\color{primary}\rule{0.45\textwidth}{1.2pt}}
\vspace{0.5em}
\end{center}

\section{Binary Search Tree (BST)}

A \textbf{Binary Search Tree} enforces the search ordering property:
\[
\text{All keys in Left Subtree} < \text{Parent Key} < \text{All keys in Right Subtree}
\]

\subsection{Time Complexity}
Search, insertion, and deletion operate in $O(h)$ time:
\begin{itemize}
    \item Balanced BST: $h = O(\log n) \implies O(\log n)$ operations.
    \item Skewed BST: $h = O(n) \implies O(n)$ worst-case operations.
\end{itemize}

\subsection{BST Node Deletion: Three Cases}
\begin{enumerate}
    \item \textbf{Node is a Leaf:} Directly unlink and delete ($O(1)$).
    \item \textbf{Node has One Child:} Bypass node, linking parent directly to child ($O(1)$).
    \item \textbf{Node has Two Children:} Replace value with its \textbf{Inorder Successor} (smallest node in right subtree) or \textbf{Inorder Predecessor} (largest in left subtree), then recursively delete that successor/predecessor node ($O(h)$).
\end{enumerate}

\section{AVL Tree (Adelson-Velsky and Landis)}

An \textbf{AVL tree} is a strictly height-balanced BST where every node satisfies:
\[
\boxed{\text{Balance Factor (BF)} = h(\text{Left Subtree}) - h(\text{Right Subtree}) \in \{-1, \; 0, \; +1\}}
\]
If $|BF| \geq 2$, the tree is rebalanced using rotations:
\begin{table}[htbp]
\centering
\small
\begin{tabularx}{\textwidth}{l X X}
\toprule
\textbf{Imbalance Case} & \textbf{Triggering Insertion Location} & \textbf{Rebalancing Rotation} \\
\midrule
LL (Left-Left)   & Inserted in left subtree of left child   & Single Right Rotation \\
RR (Right-Right) & Inserted in right subtree of right child & Single Left Rotation \\
LR (Left-Right)  & Inserted in right subtree of left child  & Left Rotation on child, then Right on parent \\
RL (Right-Left)  & Inserted in left subtree of right child  & Right Rotation on child, then Left on parent \\
\bottomrule
\end{tabularx}
\caption{The four AVL tree imbalance scenarios and rotation solutions.}
\end{table}

\subsection{Minimum Nodes in an AVL Tree of Height $h$}
Let $N(h)$ be the minimum node count for an AVL tree of height $h$:
\[
\boxed{N(h) = 1 + N(h - 1) + N(h - 2)}, \quad N(-1) = 0, \; N(0) = 1
\]
This recurrence connects to the Fibonacci sequence:
\[
\boxed{N(h) = F_{h+3} - 1} \quad (F_0=0, F_1=1, F_2=1, F_3=2, \dots)
\]
\begin{center}
\small
\begin{tabular}{ccccccccc}
\toprule
$h$ & 0 & 1 & 2 & 3 & 4 & 5 & 6 \\
\midrule
$N(h)$ & 1 & 2 & 4 & 7 & 12 & 20 & 33 \\
\bottomrule
\end{tabular}
\end{center}

\section{Red-Black Tree}

A \textbf{Red-Black tree} is a self-balancing BST with node colors (Red/Black) ensuring:
\begin{enumerate}
    \item Every node is either Red or Black.
    \item The root is always Black.
    \item All external NIL leaves are Black.
    \item If a node is Red, both of its children must be Black (no two consecutive Reds).
    \item Every path from a node to descendant NIL leaves contains identical Black nodes (equal Black-height).
\end{enumerate}

\subsection{Height Bound}
For $n$ internal nodes, the height $h$ in edges satisfies:
\[
\boxed{h \leq 2 \log_2(n + 1)}
\]

\subsection{AVL vs Red-Black Tree Comparison}
\begin{itemize}
    \item \textbf{AVL:} More rigidly balanced. Faster lookups ($O(\log n)$ with smaller constant factors). Best for lookup-heavy workloads.
    \item \textbf{Red-Black:} More relaxed balance (height up to $2 \log_2 n$). Fewer rotations during insert/delete. Best for write-intensive workloads (used in C++ \texttt{std::map}, Java \texttt{TreeMap}, Linux CFS).
\end{itemize}

\begin{center}
\vspace{1em}
{\color{primary}\rule{0.45\textwidth}{1.2pt}}\\[4pt]
{\Large\bfseries\color{primary} Section IV: B-Trees and B+ Trees}\\[2pt]
{\color{primary}\rule{0.45\textwidth}{1.2pt}}
\vspace{0.5em}
\end{center}

\section{B-Tree}

A \textbf{B-tree} is a balanced, multiway search tree optimized for systems with block-based storage (disk pages, SSD blocks).

\subsection{Order $m$ Specification}
Let order $m$ denote the \textbf{maximum number of children} an internal node can have:
\begin{itemize}
    \item Maximum children per internal node: $m$.
    \item Maximum keys per node: $m - 1$.
    \item Minimum children in non-root internal node: $\lceil m/2 \rceil$.
    \item Minimum keys in non-root internal node: $\lceil m/2 \rceil - 1$.
    \item Root node: Has at least $2$ children (if non-leaf) and at least $1$ key.
    \item All leaf nodes reside at the exact same depth.
\end{itemize}

\section{What Is a B+ Tree?}

A \textbf{B+ tree} is a balanced multiway search tree where:
\begin{enumerate}
    \item \textbf{All actual data records or record pointers reside exclusively in leaf nodes.}
    \item Internal nodes store only \textbf{separator (routing) keys} and child pointers.
    \item All leaf nodes are linked sequentially in a \textbf{doubly linked list}, facilitating $O(1)$ contiguous range scans.
\end{enumerate}

\begin{table}[htbp]
\centering
\small
\begin{tabularx}{\textwidth}{l X X}
\toprule
\textbf{Feature} & \textbf{B-Tree} & \textbf{B+ Tree} \\
\midrule
Data Storage Location     & Internal nodes AND leaf nodes           & Leaf nodes exclusively \\
Internal Node Contents    & Keys, child pointers, AND data pointers & Separator routing keys and child pointers \\
Leaf Node Linking         & Leaves not linked together              & Leaves linked in sequential linked list \\
Range Queries             & Requires full tree inorder traversal    & Extremely fast sequential leaf scan \\
Redundant Keys            & Every key appears exactly once          & Keys in internal nodes duplicate leaf keys \\
Tree Height for same data & Taller (fewer keys fit per internal node)& Shorter (higher fanout per page) \\
Application               & File systems, search engines            & Database indexing (InnoDB, PostgreSQL, SQLite) \\
\bottomrule
\end{tabularx}
\caption{Structural comparison between B-Tree and B+ Tree architectures.}
\end{table}

\section{Critical B+ Tree Calculation: Maximum Number of Keys}

\subsection{Standard Convention Formulation}
Assuming:
\begin{itemize}
    \item Order $m$ = maximum children per internal node.
    \item Each internal node holds at most $m - 1$ separator keys.
    \item Each leaf holds at most $m - 1$ data records.
    \item Height $h$ = edge count from root to leaf (root at level $0$, leaves at level $h$).
\end{itemize}

The maximum number of leaf nodes at level $h$ is $m^h$. Therefore, the maximum number of data records is:
\[
\boxed{N_{\max} = (m - 1) m^h}
\]

\subsubsection{Example: Order $m=4$, Height $h=3$ Edges}
\begin{align*}
\text{Max leaf nodes} &= 4^3 = 64 \\
\text{Max data entries per leaf} &= 4 - 1 = 3 \\
\boxed{N_{\max}} &= 64 \times 3 = \mathbf{192 \text{ records}}
\end{align*}

\subsection{When Height Means Levels ($H$)}
If height is defined as total levels $H$ (where $h = H - 1$):
\[
\boxed{N_{\max} = (m - 1) m^{H - 1}}
\]
For $m=4, H=3$ levels: $N_{\max} = 3 \times 4^2 = 48$ records.

\section{Maximum Key Entries Across All B+ Tree Nodes}

If an exam question asks for \textbf{all physical key entries} stored across both internal and leaf nodes:
\begin{align*}
\text{Internal Nodes} &= \sum_{i=0}^{h-1} m^i = \frac{m^h - 1}{m - 1} \\
\text{Internal Separator Keys} &= (m - 1) \left( \frac{m^h - 1}{m - 1} \right) = m^h - 1 \\
\text{Leaf Data Keys} &= (m - 1) m^h \\
\text{Total Physical Key Slots} &= (m^h - 1) + (m - 1) m^h = \boxed{m^{h+1} - 1}
\end{align*}
For $m=4, h=3$: $K_{\text{total}} = 4^4 - 1 = 255$ key slots ($192$ data records $+$ $63$ internal separators).

\section{Minimum Number of Keys in a B+ Tree}

Let $a = \lceil m/2 \rceil$ be the minimum child capacity of non-root internal nodes, and $q = \lceil (m - 1)/2 \rceil$ be the minimum entry capacity of non-root leaves.

For height $h \geq 1$, the root has at least $2$ children, and subsequent internal nodes have at least $a$ children.
\[
\text{Min Leaf Nodes} = L_{\min} = 2 a^{h - 1} = 2 \left\lceil \frac{m}{2} \right\rceil^{h - 1}
\]
\[
\boxed{N_{\min} = L_{\min} \times q = 2 \left\lceil \frac{m}{2} \right\rceil^{h - 1} \left\lceil \frac{m - 1}{2} \right\rceil}
\]
For $m=4, h=3$:
\begin{itemize}
    \item $a = \lceil 4/2 \rceil = 2$, $q = \lceil 3/2 \rceil = 2$.
    \item $L_{\min} = 2 \times 2^{3-1} = 8$ leaves.
    \item $\boxed{N_{\min} = 8 \times 2 = 16 \text{ records}}$.
\end{itemize}

\section{B-Tree Key-Count Calculations}

In a B-tree of order $m$ and height $h$ edges, every node (internal and leaf) holds actual data keys:
\[
\boxed{K_{\max} = m^{h+1} - 1}
\]
\[
\boxed{K_{\min} = 2 \left\lceil \frac{m}{2} \right\rceil^h - 1}
\]
For $m=4, h=3$:
\[
K_{\max} = 4^4 - 1 = 255, \quad K_{\min} = 2(2^3) - 1 = 15 \text{ keys}
\]

\begin{center}
\vspace{1em}
{\color{primary}\rule{0.45\textwidth}{1.2pt}}\\[4pt]
{\Large\bfseries\color{primary} Section V: Other Important Tree Structures}\\[2pt]
{\color{primary}\rule{0.45\textwidth}{1.2pt}}
\vspace{0.5em}
\end{center}

\section{Trie (Prefix Tree)}
A tree structure used for fast string retrieval by sharing common prefixes.
\begin{itemize}
    \item \textbf{Height:} $h \leq L$, where $L$ is the length of the longest string.
    \item \textbf{Lookup Complexity:} $O(L)$ time, completely independent of the number of words $n$.
    \item \textbf{Max Nodes for alphabet $\sigma$ and length $L$:}
    \[
    N_{\max} = \frac{\sigma^{L+1} - 1}{\sigma - 1}
    \]
\end{itemize}

\section{Expression Trees}
Binary trees where leaves are operands and internal nodes are operators ($+,-,\times,\div$).
\begin{itemize}
    \item Preorder Traversal yields \textbf{Prefix notation}.
    \item Inorder Traversal yields \textbf{Infix notation}.
    \item Postorder Traversal yields \textbf{Postfix notation}.
    \item If all operators are binary: $L = I + 1 \implies \mathbf{n = 2I + 1}$.
\end{itemize}

\section{Huffman Coding Tree}
A full binary tree used for variable-length prefix-free data compression:
\begin{itemize}
    \item Internal Nodes: $I = L - 1$; Total Nodes: $n = 2L - 1$.
    \item Weighted Path Length (Cost): $C = \sum_{i=1}^L f_i d_i$, where $f_i$ is frequency and $d_i$ is tree depth.
\end{itemize}

\section{Segment Tree}
Stores interval information (range sums, minimums, maximums):
\begin{itemize}
    \item Full binary tree with $n$ leaves has exactly $\mathbf{2n - 1}$ logical nodes.
    \item Standard array allocation size is bounded by $\mathbf{4n}$.
    \item Build: $O(n)$; Range Query: $O(\log n)$; Point Update: $O(\log n)$.
\end{itemize}

\begin{center}
\vspace{1em}
{\color{primary}\rule{0.45\textwidth}{1.2pt}}\\[4pt]
{\Large\bfseries\color{primary} Section VI: Master Complexity Comparison}\\[2pt]
{\color{primary}\rule{0.45\textwidth}{1.2pt}}
\vspace{0.5em}
\end{center}

\begin{table}[htbp]
\centering
\small
\begin{tabularx}{\textwidth}{l X X X}
\toprule
\textbf{Tree Data Structure} & \textbf{Search Complexity} & \textbf{Insert Complexity} & \textbf{Delete Complexity} \\
\midrule
Ordinary Binary Tree & $O(n)$              & $O(n)$              & $O(n)$ \\
Ordinary BST         & $O(h)$ ($O(n)$ max)& $O(h)$              & $O(h)$ \\
Balanced BST (AVL)   & $O(\log n)$         & $O(\log n)$         & $O(\log n)$ \\
Red-Black Tree       & $O(\log n)$         & $O(\log n)$         & $O(\log n)$ \\
Binary Heap          & $O(n)$ arbitrary    & $O(\log n)$         & $O(\log n)$ extract root \\
B-Tree               & $O(\log n)$         & $O(\log n)$         & $O(\log n)$ \\
B+ Tree              & $O(\log n)$         & $O(\log n)$         & $O(\log n)$ \\
Trie (Prefix Tree)   & $O(L)$ key length   & $O(L)$              & $O(L)$ \\
Segment Tree         & $O(\log n)$ range   & Static build $O(n)$ & Point update $O(\log n)$ \\
Huffman Tree         & Static tree         & Build $O(n \log n)$ & Decode $O(L)$ \\
\bottomrule
\end{tabularx}
\caption{Master operational time complexity across tree architectures.}
\end{table}

\begin{center}
\vspace{1em}
{\color{primary}\rule{0.45\textwidth}{1.2pt}}\\[4pt]
{\Large\bfseries\color{primary} Section VII: Solved Mathematical Problems}\\[2pt]
{\color{primary}\rule{0.45\textwidth}{1.2pt}}
\vspace{0.5em}
\end{center}

\begin{quote}
\textbf{Problem 1: Nodes in Perfect Binary Tree of Height 6}\\
$n = 2^{6+1} - 1 = 2^7 - 1 = \mathbf{127 \text{ nodes}}$. Leaves $L = 2^6 = \mathbf{64}$. Internal $I = 64 - 1 = \mathbf{63}$.
\end{quote}

\begin{quote}
\textbf{Problem 2: Full Binary Tree with 25 Leaves}\\
$I = L - 1 = 25 - 1 = \mathbf{24 \text{ internal nodes}}$. Total $n = 24 + 25 = \mathbf{49 \text{ nodes}}$.
\end{quote}

\begin{quote}
\textbf{Problem 3: Minimum and Maximum Height for 40 Nodes}\\
$h_{\min} = \lceil \log_2(41) \rceil - 1 = 6 - 1 = \mathbf{5 \text{ edges}}$. $h_{\max} = n - 1 = \mathbf{39 \text{ edges}}$.
\end{quote}

\begin{quote}
\textbf{Problem 4: Leaves from Two-Child Node Count ($n=50, n_2=18$)}\\
$n_0 = n_2 + 1 = 18 + 1 = \mathbf{19 \text{ leaves}}$. Single child nodes $n_1 = 50 - 19 - 18 = \mathbf{13}$.
\end{quote}

\begin{quote}
\textbf{Problem 5: Can a Full Binary Tree Have Exactly 100 Nodes?}\\
No. A full binary tree has $n = 2I + 1$, which is strictly \textbf{odd}. Since $100$ is even, it is impossible.
\end{quote}

\begin{quote}
\textbf{Problem 6: Distinct BST Structures for 4 Keys}\\
$C_4 = \frac{1}{5} \binom{8}{4} = \frac{70}{5} = \mathbf{14 \text{ distinct BSTs}}$.
\end{quote}

\begin{quote}
\textbf{Problem 7: B+ Tree Capacity ($m=4$, Height 4 edges, 3 records/leaf)}\\
$L_{\max} = 4^4 = 256$ leaves. $\mathbf{N_{\max}} = 256 \times 3 = \mathbf{768 \text{ data records}}$.
\end{quote}

\begin{quote}
\textbf{Problem 8: B-Tree Keys ($m=6$, Height 2 edges)}\\
$K_{\max} = 6^{2+1} - 1 = 216 - 1 = \mathbf{215 \text{ keys}}$.\\
$K_{\min} = 2 \lceil 6/2 \rceil^2 - 1 = 2(3^2) - 1 = \mathbf{17 \text{ keys}}$.
\end{quote}

\begin{quote}
\textbf{Problem 9: Feasible Height for 18-Node AVL Tree}\\
$N(4) = 12, N(5) = 20$. Since $18 < 20$, height $5$ is impossible. Max feasible height is $\mathbf{h = 4}$.
\end{quote}

\begin{quote}
\textbf{Problem 10: Segment Tree Nodes for Array of Size 75}\\
$N = 2(75) - 1 = \mathbf{149 \text{ logical nodes}}$ (buffer bound $\approx 4 \times 75 = 300$).
\end{quote}

\begin{center}
\vspace{1em}
{\color{primary}\rule{0.45\textwidth}{1.2pt}}\\[4pt]
{\Large\bfseries\color{primary} Section VIII: Master Formula Cheat-Sheet}\\[2pt]
{\color{primary}\rule{0.45\textwidth}{1.2pt}}
\vspace{0.5em}
\end{center}

\begin{table}[htbp]
\centering
\small
\begin{tabularx}{\textwidth}{l X}
\toprule
\textbf{Tree Concept} & \textbf{Master Formula} \\
\midrule
Edge count for $n$ nodes & $E = n - 1$ \\
Max nodes at level $i$ & $m^i$ (general) $\;/\; 2^i$ (binary) \\
Max nodes for height $h$ & $\frac{m^{h+1}-1}{m-1}$ (general) $\;/\; 2^{h+1}-1$ (binary) \\
Min height for $n$ nodes & $\lceil \log_m(n(m-1)+1) \rceil - 1$ \\
Full binary tree leaves & $L = I + 1$ \\
Full binary tree total nodes & $n = 2I + 1$ (always odd) \\
Leaf vs 2-child relation & $n_0 = n_2 + 1$ (holds for \textbf{all} binary trees) \\
Null child pointers & $n + 1$ \\
Number of binary tree shapes & $C_n = \frac{1}{n+1} \binom{2n}{n}$ (Catalan Number) \\
AVL min nodes recurrence & $N(h) = 1 + N(h-1) + N(h-2) = F_{h+3}-1$ \\
Red-Black height upper bound & $h \leq 2 \log_2(n+1)$ \\
B+ Tree max data records & $(m - 1) m^h$ \\
B+ Tree max physical key slots & $m^{h+1} - 1$ \\
B+ Tree min data records ($h \geq 1$) & $2 \lceil m/2 \rceil^{h-1} \lceil (m-1)/2 \rceil$ \\
B-Tree max keys & $m^{h+1} - 1$ \\
B-Tree min keys ($h \geq 1$) & $2 \lceil m/2 \rceil^h - 1$ \\
Binary Heap height & $h = \lfloor \log_2 n \rfloor$ \\
Segment Tree logical nodes & $2n - 1$ \\
\bottomrule
\end{tabularx}
\caption{Comprehensive master mathematical formula sheet for recruitment examinations.}
\end{table}

\begin{center}
\vspace{1em}
{\color{primary}\rule{0.45\textwidth}{1.2pt}}\\[4pt]
{\Large\bfseries\color{primary} Section IX: Common Examination Traps}\\[2pt]
{\color{primary}\rule{0.45\textwidth}{1.2pt}}
\vspace{0.5em}
\end{center}

\begin{enumerate}
    \item \textbf{Edges vs Levels:} Height in edges ($h$) has $h + 1$ levels. If an exam defines height as levels $H$, replace $h$ with $H - 1$.
    \item \textbf{Full vs Complete:} Full means every node has $0$ or $2$ children. Complete means levels are fully packed with the last level filled left-to-right.
    \item \textbf{Even Node Count Rule:} A full binary tree \emph{cannot} have an even number of nodes.
    \item \textbf{BST Ordering vs Heap Ordering:} In a BST, left $<$ parent $<$ right (horizontal order). In a heap, parent $\geq$ children (vertical order, with no left/right constraint).
    \item \textbf{B+ Tree Separators vs Data Records:} Internal keys in a B+ tree are routing separators and duplicate leaf keys. Total physical key slots ($m^{h+1}-1$) exceeds unique data records ($(m-1)m^h$).
    \item \textbf{B+ Tree Order Ambiguity:} Always check whether order $m$ denotes maximum children ($m$ children, $m-1$ keys) or maximum keys ($2d$ keys, $2d+1$ children).
    \item \textbf{Arbitrary Key Search in Heap:} Searching for an arbitrary key in a binary heap takes $O(n)$, not $O(\log n)$.
\end{enumerate}

\begin{quote}
\textbf{Final Architectural Takeaway:} Tree structures balance operational time by exploiting geometric branching. Binary trees, AVL trees, and Red-Black trees optimize memory-resident logarithmic searches, while B-trees and B+ trees optimize disk and page block I/O by minimizing tree height through massive branching fanout.
\end{quote}


\begin{center}
\vspace{1em}
{\color{primary}\rule{0.5\textwidth}{1.5pt}}\\[6pt]
{\Large\bfseries\color{primary} Part V: Graphs, Adjacency Representations, and Hash Tables}\\[3pt]
{\color{primary}\rule{0.5\textwidth}{1.5pt}}
\vspace{1em}
\end{center}

This chapter covers three important areas of data structures for Bangladesh government IT recruitment examinations: graph classifications and mathematical properties, hash tables and collision-resolution techniques, and the comparison between adjacency lists and adjacency matrices.

The emphasis is on definitions, formulas, worked calculations, algorithms, time complexity, practical applications, and common MCQ traps.

\begin{center}
\vspace{0.5em}
{\color{primary}\rule{0.45\textwidth}{1.2pt}}\\[4pt]
{\Large\bfseries\color{primary} Section I: Graphs and Their Properties}\\[2pt]
{\color{primary}\rule{0.45\textwidth}{1.2pt}}
\vspace{0.5em}
\end{center}

\section{What is a Graph?}

A \textbf{graph} is a non-linear data structure consisting of a set of vertices and a set of edges connecting those vertices.

Mathematically, a graph is commonly represented as:
\[
G = (V, E)
\]
where:
\begin{itemize}
    \item $V$ is the set of vertices (also called nodes).
    \item $E$ is the set of edges connecting vertices.
\end{itemize}

For example, consider a computer network with five computers. Each computer can be represented as a vertex, and each network connection can be represented as an edge.

Graphs can model computer networks, road maps, social networks, communication systems, dependencies among software modules, and relationships between database entities.

\subsection{Basic Graph Terminology}

\begin{table}[htbp]
\centering
\small
\begin{tabularx}{\textwidth}{l X}
\toprule
\textbf{Term} & \textbf{Formal Definition} \\
\midrule
Vertex (node)     & An individual entity in the graph. \\
Edge              & A connection between two vertices in an undirected graph. \\
Arc               & A directed connection from one vertex to another. \\
Adjacent vertices & Vertices directly connected by an edge. \\
Incident edge     & An edge connected to a particular vertex. \\
Degree            & Number of incident edges for a vertex in an undirected graph. \\
In-degree         & Number of incoming arcs to a vertex in a directed graph. \\
Out-degree        & Number of outgoing arcs from a vertex in a directed graph. \\
Path              & A sequence of vertices connected by edges. \\
Walk              & A sequence of adjacent vertices in which vertices and edges may repeat. \\
Trail             & A walk in which no edge is repeated. \\
Simple path       & A path that does not repeat vertices. \\
Cycle             & A closed path or closed traversal that returns to its starting vertex. \\
Connected graph   & An undirected graph in which every pair of vertices has a path between them. \\
Component         & A maximal connected subgraph of an undirected graph. \\
Weight            & A numerical value assigned to an edge, such as distance, cost, or latency. \\
Subgraph          & A graph formed from a subset of vertices and edges of another graph. \\
\bottomrule
\end{tabularx}
\caption{Fundamental graph terminology and definitions.}
\end{table}

\begin{quote}
\textbf{Recruitment Exam Distinction:} The terms walk, trail, path, and cycle are sometimes defined with slightly different conventions in textbooks. For recruitment exams, remember:
\begin{itemize}
    \item \textbf{Walk:} May repeat both edges and vertices.
    \item \textbf{Trail:} Cannot repeat edges (may repeat vertices).
    \item \textbf{Simple Path:} Cannot repeat vertices (and thus cannot repeat edges).
\end{itemize}
\end{quote}

\section{Types of Graphs}

\subsection{Undirected Graph}
An \textbf{undirected graph} is a graph in which edges do not have a direction. An edge connecting vertices $u$ and $v$ is bidirectional and is commonly written as an unordered pair $\{u, v\}$ or $(u, v)$.

\textbf{Properties:}
\begin{itemize}
    \item $(u, v)$ is identical to $(v, u)$.
    \item Traversal is permitted in both directions.
    \item Used to model bidirectional communication, friendship networks, and two-way roads.
\end{itemize}

\subsection{Directed Graph (Digraph)}
A \textbf{directed graph} (or \textbf{digraph}) is a graph in which every edge has a specific direction. Each directed edge is called an \textbf{arc} and is represented as an ordered pair $(u, v)$, directed from $u$ to $v$.

\textbf{Properties:}
\begin{itemize}
    \item $(u, v)$ is distinct from $(v, u)$.
    \item $u$ is the tail or source; $v$ is the head or destination.
    \item A vertex has both an \textbf{in-degree} (incoming edges) and an \textbf{out-degree} (outgoing edges).
    \item Used to model one-way traffic, task dependencies, webpage hyperlinks, and function call graphs.
\end{itemize}

\subsection{Weighted Graph}
A \textbf{weighted graph} is a graph in which each edge is assigned a numerical value called a weight. Depending on the application, the weight may represent distance, cost, travel time, bandwidth, or network delay.

Shortest-path algorithms such as Dijkstra's algorithm and minimum spanning tree algorithms such as Kruskal's and Prim's operate on weighted graphs.

\subsection{Simple Graph}
A \textbf{simple graph} is an undirected or directed graph that contains:
\begin{itemize}
    \item \textbf{No self-loops:} An edge connecting a vertex to itself.
    \item \textbf{No parallel (multiple) edges:} Multiple edges connecting the same pair of vertices in the same direction.
\end{itemize}
Unless specified otherwise, theoretical exam questions generally assume simple graphs.

\subsection{Multigraph}
A \textbf{multigraph} is a graph that allows multiple edges between the same pair of vertices. It may or may not allow self-loops depending on the textbook definition.

Multigraphs are useful when modeling multiple flight routes between the same cities or multiple telecommunication links between two switching centers.

\subsection{Pseudograph}
A \textbf{pseudograph} is a graph that allows both parallel edges and self-loops. It is the most general graph variant encountered in discrete mathematics.

\subsection{Complete Graph}
A \textbf{complete graph} is a simple undirected graph in which an edge connects every pair of distinct vertices. A complete graph with $n$ vertices is denoted by $K_n$.

Every vertex in $K_n$ is adjacent to every other vertex.

\textbf{Example:} In $K_4$, there are 4 vertices. The total number of edges is:
\[
E = \frac{4 \times 3}{2} = 6 \text{ edges.}
\]
Each vertex has a degree of:
\[
\deg(v) = n - 1 = 4 - 1 = 3.
\]

\subsection{Regular Graph}
A graph is \textbf{$k$-regular} if every vertex has the same degree $k$.

\textbf{Example:} A cycle graph $C_n$ (for $n \geq 3$) is 2-regular because every vertex has degree 2. The complete graph $K_n$ is $(n-1)$-regular.

\subsection{Null Graph and Trivial Graph}
\begin{itemize}
    \item \textbf{Null graph (Empty graph):} A graph containing a set of vertices but \textbf{no edges} ($E = \emptyset$).
    \item \textbf{Trivial graph:} A graph consisting of exactly \textbf{one vertex and zero edges}.
\end{itemize}

\section{Connectivity, Cycles, and Special Graph Classes}

\subsection{Connected and Disconnected Graphs}
An undirected graph is \textbf{connected} if there is a path between every pair of vertices. Otherwise, it is \textbf{disconnected}.

A disconnected graph consists of two or more independent connected subgraphs called \textbf{connected components}.

\textbf{Useful Properties:}
\begin{itemize}
    \item A connected undirected graph with $V$ vertices must have at least $V - 1$ edges.
    \item If an undirected graph has fewer than $V - 1$ edges, it is guaranteed to be disconnected.
    \item An undirected graph with $V$ vertices and more than $\frac{(V-1)(V-2)}{2}$ edges is guaranteed to be connected.
\end{itemize}

\subsection{Strongly Connected and Weakly Connected Directed Graphs}
For directed graphs:
\begin{itemize}
    \item \textbf{Strongly connected:} For every pair of vertices $u$ and $v$, there is a directed path from $u$ to $v$ \textbf{and} a directed path from $v$ to $u$.
    \item \textbf{Weakly connected:} The graph is not strongly connected, but replacing all directed edges with undirected edges produces a connected undirected graph.
\end{itemize}

\subsection{Cyclic and Acyclic Graphs}
\begin{itemize}
    \item \textbf{Cyclic graph:} A graph containing at least one cycle.
    \item \textbf{Acyclic graph:} A graph containing no cycles. An undirected acyclic graph is a \textbf{forest}; if it is also connected, it is a \textbf{tree}.
\end{itemize}

\subsubsection{Directed Acyclic Graph (DAG)}
A \textbf{directed acyclic graph (DAG)} is a directed graph containing no directed cycles.
\begin{itemize}
    \item \textbf{Topological Sorting:} Vertices in a DAG can be linearly ordered such that for every directed edge $(u, v)$, vertex $u$ comes before vertex $v$.
    \item \textbf{Applications:} Task scheduling, build systems (e.g., Make/Gradle), prerequisite resolution, and spreadsheet formula evaluation.
\end{itemize}

\subsection{Bipartite Graph}
A graph is \textbf{bipartite} if its vertex set $V$ can be partitioned into two disjoint subsets $V_1$ and $V_2$ such that every edge connects a vertex in $V_1$ to a vertex in $V_2$. No edge connects two vertices in the same subset.

\begin{tcolorbox}[title=\textbf{Fundamental Theorem on Bipartite Graphs}]
An undirected graph is bipartite if and only if it contains \textbf{no odd-length cycles}.
\end{tcolorbox}

\subsubsection{Complete Bipartite Graph}
A \textbf{complete bipartite graph}, denoted by $K_{a, b}$, partitions vertices into sets of size $a$ and $b$, where every vertex in $V_1$ connects to every vertex in $V_2$.
\begin{itemize}
    \item Total vertices: $V = a + b$.
    \item Total edges: $\boxed{E = a \times b}$.
\end{itemize}

\subsection{Planar Graph}
A \textbf{planar graph} is a graph that can be drawn in a plane such that no two edges cross each other.

\begin{tcolorbox}[title=\textbf{Euler's Formula for Connected Planar Graphs}]
For any connected planar graph with $V$ vertices, $E$ edges, and $F$ faces (including the unbounded outer face):
\[
\boxed{V - E + F = 2}
\]
For a simple connected planar graph with $V \geq 3$:
\[
E \leq 3V - 6
\]
\end{tcolorbox}
The complete graph $K_5$ and the complete bipartite graph $K_{3,3}$ are the two minimal non-planar graphs (Kuratowski's Theorem).

\subsection{Eulerian Graph}
\begin{itemize}
    \item \textbf{Eulerian Trail:} A walk that visits every \textbf{edge} of the graph exactly once.
    \item \textbf{Eulerian Circuit:} An Eulerian trail that starts and ends at the same vertex.
    \item \textbf{Eulerian Graph:} A connected graph containing an Eulerian circuit.
\end{itemize}

\textbf{Euler's Condition for Undirected Graphs:}
\begin{itemize}
    \item An Eulerian circuit exists if and only if the graph is connected and \textbf{every vertex has an even degree}.
    \item An Eulerian trail (open path) exists if and only if the graph is connected and \textbf{exactly two vertices have odd degrees} (which serve as start and end vertices).
\end{itemize}

\subsection{Hamiltonian Graph}
\begin{itemize}
    \item \textbf{Hamiltonian Path:} A path that visits every \textbf{vertex} in the graph exactly once.
    \item \textbf{Hamiltonian Cycle:} A closed loop that visits every vertex exactly once and returns to the start vertex.
    \item \textbf{Hamiltonian Graph:} A graph containing a Hamiltonian cycle.
\end{itemize}

\begin{quote}
\textbf{Recruitment Comparison:} Euler = visit every \textbf{edge} once ($O(V+E)$ polynomial time check). Hamiltonian = visit every \textbf{vertex} once (NP-complete decision problem).
\end{quote}

\section{Fundamental Graph Formulas}

\subsection{Handshaking Lemma for Undirected Graphs}
In any undirected graph, the sum of the degrees of all vertices equals twice the number of edges:
\[
\boxed{\sum_{v \in V} \deg(v) = 2E}
\]
\textbf{Key Consequence:} In any undirected graph, the number of vertices with an \textbf{odd degree must be even}.

\begin{tcolorbox}[title=\textbf{Worked Example: Handshaking Lemma}]
An undirected graph has six vertices with degrees: $2, 3, 1, 4, 2, 2$. Find the number of edges.\\
\textbf{Solution:}
\[
\sum \deg(v) = 2 + 3 + 1 + 4 + 2 + 2 = 14
\]
Applying the handshaking lemma:
\[
2E = 14 \implies E = \frac{14}{2} = 7 \text{ edges.}
\]
\end{tcolorbox}

\subsection{Directed Graph Degree Formulas}
In a directed graph:
\begin{itemize}
    \item The sum of in-degrees equals the total number of arcs: $\sum \text{in-deg}(v) = E$.
    \item The sum of out-degrees equals the total number of arcs: $\sum \text{out-deg}(v) = E$.
    \item Therefore:
    \[
    \boxed{\sum_{v \in V} \text{in-deg}(v) = \sum_{v \in V} \text{out-deg}(v) = E}
    \]
\end{itemize}

\subsection{Maximum Edges in a Simple Graph}
\begin{itemize}
    \item \textbf{Simple Undirected Graph with $n$ vertices:}
    \[
    \boxed{E_{\max} = \frac{n(n-1)}{2} = \binom{n}{2}}
    \]
    \item \textbf{Simple Directed Graph without self-loops with $n$ vertices:}
    \[
    \boxed{E_{\max} = n(n-1)}
    \]
\end{itemize}

\begin{tcolorbox}[title=\textbf{Worked Example: Maximum Edges}]
Find the maximum number of edges in a simple undirected graph with 12 vertices.\\
\textbf{Solution:}
\[
E_{\max} = \frac{12(12-1)}{2} = \frac{12 \times 11}{2} = 66 \text{ edges.}
\]
\end{tcolorbox}

\subsection{Trees and Forests}
\begin{itemize}
    \item A tree with $V$ vertices has exactly:
    \[
    \boxed{E = V - 1}
    \]
    \item A forest containing $V$ vertices and $C$ connected components has:
    \[
    \boxed{E = V - C}
    \]
\end{itemize}

\subsection{Density of a Graph}
The density $D$ of a simple undirected graph measures the ratio of actual edges to maximum possible edges:
\[
D = \frac{2E}{V(V-1)}
\]
\begin{itemize}
    \item If $E \ll V^2$ (typically $E \approx O(V)$), the graph is \textbf{sparse}.
    \item If $E \approx O(V^2)$ (approaching $E_{\max}$), the graph is \textbf{dense}.
\end{itemize}

\newpage
\begin{center}
\vspace{0.5em}
{\color{primary}\rule{0.45\textwidth}{1.2pt}}\\[4pt]
{\Large\bfseries\color{primary} Section II: Hash Tables and Collision Resolution}\\[2pt]
{\color{primary}\rule{0.45\textwidth}{1.2pt}}
\vspace{0.5em}
\end{center}

\section{What Is a Hash Table?}

A \textbf{hash table} is an associative data structure that maps keys to values, achieving average-case $O(1)$ time complexity for insertions, deletions, and lookups.

\subsection{Basic Components}
\begin{itemize}
    \item \textbf{Key:} The unique identifier used to look up or store data (e.g., student ID, roll number).
    \item \textbf{Value:} The data record associated with the key.
    \item \textbf{Hash Function:} A mathematical algorithm that converts a key into a valid array index.
    \item \textbf{Hash Table Array (Bucket Array):} The underlying array where key-value pairs or bucket pointers are stored.
\end{itemize}

\subsection{What is a Hash Function?}
A \textbf{hash function} $h(k)$ maps an arbitrary key $k$ from a large universe of keys $U$ to a fixed integer slot in the range $[0, m-1]$, where $m$ is the table size:
\[
h : U \longrightarrow \{0, 1, 2, \dots, m-1\}
\]

\subsection{Characteristics of a Good Hash Function}
\begin{enumerate}
    \item \textbf{Deterministic:} The same key must always produce the identical hash index.
    \item \textbf{Uniform Distribution:} Keys should be uniformly distributed across all slots to minimize clustering and collisions.
    \item \textbf{Fast Computation:} The function should compute in $O(1)$ time.
    \item \textbf{Full Range Utilization:} Every slot from $0$ to $m-1$ should be reachable.
\end{enumerate}

\section{Types of Hash Functions}

\subsection{Division Method}
The most common hash function computes the remainder of the key divided by the table size $m$:
\[
\boxed{h(k) = k \bmod m}
\]
\textbf{Best Practice:} The table size $m$ should typically be a \textbf{prime number} not close to a power of 2 or 10, ensuring better distribution and reducing collisions.

\begin{table}[htbp]
\centering
\small
\begin{tabular}{c c c}
\toprule
\textbf{Key ($k$)} & \textbf{Calculation ($k \bmod 10$)} & \textbf{Table Index} \\
\midrule
23 & $23 \bmod 10$ & 3 \\
35 & $35 \bmod 10$ & 5 \\
42 & $42 \bmod 10$ & 2 \\
78 & $78 \bmod 10$ & 8 \\
\bottomrule
\end{tabular}
\caption{Division hash function with table size $m = 10$.}
\end{table}

\subsection{Multiplication Method}
The key $k$ is multiplied by a constant real fraction $A$ where $0 < A < 1$. The fractional part of $k A$ is extracted and multiplied by $m$, followed by taking the floor:
\[
\boxed{h(k) = \lfloor m (k A \bmod 1) \rfloor}
\]
Donald Knuth suggested using the golden ratio fractional value:
\[
A \approx \frac{\sqrt{5} - 1}{2} \approx 0.6180339887\dots
\]
An advantage of this method is that the choice of table size $m$ is less critical (powers of 2 are easily accommodated via bit shifts).

\subsection{Mid-Square Method}
The key is squared, and selected middle digits (or bits) of the result are extracted as the hash code:
\[
k \xrightarrow{\quad\text{square}\quad} k^2 \xrightarrow{\quad\text{extract middle digits}\quad} \text{index}
\]
\textbf{Example:} For key $44$:
\[
44^2 = 1936
\]
Extracting the middle two digits gives $93$. This value is then mapped modulo $m$ into the table range.

\subsection{Folding Method}
The key is partitioned into multiple pieces of equal digit width. These parts are then added together, and the sum modulo $m$ is used:
\textbf{Example:} For a 9-digit student registration number $123456789$ and $m = 1000$:
\[
\text{Parts: } 123, 456, 789 \implies \text{Sum} = 123 + 456 + 789 = 1368
\]
\[
\text{Index} = 1368 \bmod 1000 = 368.
\]

\subsection{Direct Addressing versus Hashing}
\begin{itemize}
    \item \textbf{Direct Addressing:} Uses key $k$ directly as array index $A[k]$. This requires $O(1)$ worst-case time, but requires memory proportional to the size of the entire key universe $|U|$. If the universe is large and sparse (e.g., National ID numbers), it wastes prohibitive amounts of memory.
    \item \textbf{Hashing:} Uses a hash function to compress a large universe $U$ into a much smaller table array of size $m$, where $m \ll |U|$.
\end{itemize}

\section{What Is a Collision?}

A \textbf{collision} occurs when two or more distinct keys produce the identical hash index:
\[
h(k_1) = h(k_2) \quad \text{where } k_1 \neq k_2
\]
Because the universe of keys is far larger than the table size ($|U| > m$), by the \textbf{Pigeonhole Principle}, collisions are theoretically unavoidable in any realistic general hash table.

\begin{table}[htbp]
\centering
\small
\begin{tabular}{c c}
\toprule
\textbf{Key} & \textbf{Hash Index ($h(k) = k \bmod 7$)} \\
\midrule
10 & $10 \bmod 7 = 3$ \\
17 & $17 \bmod 7 = 3$ \\
24 & $24 \bmod 7 = 3$ \\
31 & $31 \bmod 7 = 3$ \\
38 & $38 \bmod 7 = 3$ \\
\bottomrule
\end{tabular}
\caption{Collision example: Multiple keys mapping to slot 3.}
\end{table}

\subsection{Main Collision-Resolution Techniques}
\begin{enumerate}
    \item \textbf{Open Hashing (Closed Addressing / Separate Chaining):} Elements that collide at the same slot are chained together in an auxiliary data structure (typically a linked list).
    \item \textbf{Closed Hashing (Open Addressing):} All elements are stored directly inside the table array. When a collision occurs, alternative vacant slots are probed systematically:
    \begin{itemize}
        \item Linear Probing
        \item Quadratic Probing
        \item Double Hashing
    \end{itemize}
\end{enumerate}

\section{Separate Chaining}

In \textbf{separate chaining}, each table slot contains a pointer/reference to a linked list (or dynamic array/BST) of elements that hash to that slot.

\begin{table}[htbp]
\centering
\small
\begin{tabular}{c l}
\toprule
\textbf{Slot Index} & \textbf{Chained Bucket Entries} \\
\midrule
0 & Empty \\
1 & Empty \\
2 & $[44] \rightarrow \text{NULL}$ \\
3 & $[10] \rightarrow [17] \rightarrow [24] \rightarrow \text{NULL}$ \\
4 & Empty \\
5 & $[12] \rightarrow \text{NULL}$ \\
6 & $[13] \rightarrow [20] \rightarrow \text{NULL}$ \\
\bottomrule
\end{tabular}
\caption{Separate chaining collision resolution.}
\end{table}

\subsection{Operations in Separate Chaining}
\begin{itemize}
    \item \textbf{Insertion:} Compute $h(k)$. Traverse bucket list to check for duplicate key (update if found), or insert the new node at the front/back in $O(1)$ time.
    \item \textbf{Searching:} Compute $h(k)$. Linearly search the linked list at index $h(k)$ until the key is found or list ends.
    \item \textbf{Deletion:} Compute $h(k)$. Traverse bucket list, locate the node, unlink it, and free memory in $O(\text{bucket length})$ time.
\end{itemize}

\subsection{Complexity of Separate Chaining}

\begin{table}[htbp]
\centering
\small
\begin{tabularx}{\textwidth}{l X c}
\toprule
\textbf{Operation} & \textbf{Average / Expected Complexity} & \textbf{Worst-Case Complexity} \\
\midrule
Search & $O(1 + \alpha)$ & $O(n)$ \\
Insert & $O(1)$ (at head) or $O(1 + \alpha)$ (with uniqueness check) & $O(n)$ \\
Delete & $O(1 + \alpha)$ & $O(n)$ \\
\bottomrule
\end{tabularx}
\caption{Time complexity of separate chaining (where $\alpha = n / m$).}
\end{table}

\section{Open Addressing}

In \textbf{open addressing}, all elements are stored directly within the table array of size $m$. No pointers or external lists are used.
\begin{itemize}
    \item The total number of stored keys $n$ can never exceed table capacity $m$ ($n \leq m$, $\alpha \leq 1$).
    \item When a collision occurs at initial index $h(k)$, a sequence of probe slots $h(k, 0), h(k, 1), h(k, 2), \dots$ is checked until an empty slot is encountered.
\end{itemize}

\subsection{Linear Probing}
Linear probing inspects consecutive slots sequentially with step size $1$:
\[
\boxed{h(k, i) = (h'(k) + i) \bmod m, \quad i = 0, 1, 2, \dots, m-1}
\]

\begin{tcolorbox}[title=\textbf{Worked Example: Linear Probing}]
Given table size $m = 7$, hash function $h'(k) = k \bmod 7$, insert keys: $10, 17, 24, 8$ using linear probing.
\begin{itemize}
    \item Insert 10: $h'(10) = 10 \bmod 7 = 3$. Slot 3 is empty $\implies$ place \textbf{10 at index 3}.
    \item Insert 17: $h'(17) = 17 \bmod 7 = 3$ (Collision). Check $(3+1)\bmod 7 = 4$. Slot 4 empty $\implies$ place \textbf{17 at index 4}.
    \item Insert 24: $h'(24) = 24 \bmod 7 = 3$ (Collision). Check slot 4 (occupied), check $(3+2)\bmod 7 = 5$. Slot 5 empty $\implies$ place \textbf{24 at index 5}.
    \item Insert 8: $h'(8) = 8 \bmod 7 = 1$. Slot 1 is empty $\implies$ place \textbf{8 at index 1}.
\end{itemize}
Final array state: $\text{Slot 1} = 8, \text{Slot 3} = 10, \text{Slot 4} = 17, \text{Slot 5} = 24$.
\end{tcolorbox}

\subsubsection*{Primary Clustering}
The primary disadvantage of linear probing is \textbf{primary clustering}. When multiple colliding or non-colliding keys occupy adjacent consecutive slots, they form long continuous occupied blocks. Any subsequent key hashing into or near this block must traverse the entire cluster, causing search and insertion times to degrade rapidly.

\subsection{Quadratic Probing}
Quadratic probing reduces primary clustering by spacing probe attempts quadratically:
\[
\boxed{h(k, i) = (h'(k) + c_1 i + c_2 i^2) \bmod m}
\]
A common standard formulation is:
\[
h(k, i) = (h'(k) + i^2) \bmod m
\]
\textbf{Secondary Clustering:} While quadratic probing resolves primary clustering, keys that have the same initial hash index $h'(k_1) = h'(k_2)$ follow the identical sequence of probes. This milder phenomenon is called \textbf{secondary clustering}.

\subsection{Double Hashing}
Double hashing is the most effective open addressing technique. It uses two independent hash functions:
\[
\boxed{h(k, i) = (h_1(k) + i \cdot h_2(k)) \bmod m}
\]
where $h_1(k)$ computes the initial index and $h_2(k)$ computes the probe step size.

\textbf{Requirements for $h_2(k)$:}
\begin{itemize}
    \item $h_2(k)$ must never evaluate to $0$ ($h_2(k) \neq 0$), as this would produce an infinite loop at index $h_1(k)$.
    \item $h_2(k)$ must be relatively prime to $m$ ($\gcd(h_2(k), m) = 1$) so that every table slot can be probed. A common scheme is $m = \text{prime}$ and $h_2(k) = 1 + (k \bmod (m - 1))$.
\end{itemize}

\section{Comparison of Collision-Resolution Methods}

\begin{table}[htbp]
\centering
\footnotesize
\begin{tabularx}{\textwidth}{p{2.5cm} X X X X}
\toprule
\textbf{Feature} & \textbf{Separate Chaining} & \textbf{Linear Probing} & \textbf{Quadratic Probing} & \textbf{Double Hashing} \\
\midrule
Storage Method & Bucket lists outside array & Directly inside table array & Directly inside table array & Directly inside table array \\
Collision Strategy & Insert into bucket list & Sequential scan ($+i$) & Quadratic offset ($+i^2$) & Step by $h_2(k)$ \\
Clustering Issue & No array clustering & Severe primary clustering & Secondary clustering & Virtually no clustering \\
Capacity Limit & Can exceed table size $m$ & Strictly limited ($n \leq m$) & Strictly limited ($n \leq m$) & Strictly limited ($n \leq m$) \\
Load Factor ($\alpha$) & Can be $> 1$ (often $1$--$3$) & Must stay $\alpha < 0.7$ & Must stay $\alpha < 0.5$ & Can stay $\alpha < 0.8$ \\
Deletion & Simple list node removal & Requires tombstone flag & Requires tombstone flag & Requires tombstone flag \\
Cache Efficiency & Poor (pointer chasing) & Excellent (contiguous) & Moderate & Moderate \\
\bottomrule
\end{tabularx}
\caption{Comprehensive comparison of collision-resolution mechanisms.}
\end{table}

\section{Load Factor and Rehashing}

\subsection{What is Load Factor?}
The \textbf{load factor} $\alpha$ measures the ratio of the number of stored keys $n$ to the total table capacity $m$:
\[
\boxed{\alpha = \frac{n}{m}}
\]
\begin{itemize}
    \item In open addressing, $\alpha$ cannot exceed $1$ ($\alpha \leq 1$).
    \item In separate chaining, $\alpha$ represents the average length of a bucket chain and can exceed $1$.
\end{itemize}

\begin{tcolorbox}[title=\textbf{Worked Example: Load Factor}]
A hash table has 100 buckets and stores 75 entries. Calculate its load factor.\\
\textbf{Solution:}
\[
\alpha = \frac{n}{m} = \frac{75}{100} = 0.75
\]
\end{tcolorbox}

\subsection{Why Does Load Factor Matter?}
As $\alpha$ increases:
\begin{itemize}
    \item Collision probability rises sharply.
    \item Probe sequences in open addressing grow longer, degrading lookup times toward $O(n)$.
    \item Bucket chains in separate chaining lengthen.
\end{itemize}

\subsection{Rehashing}
When the load factor exceeds a predetermined threshold (commonly $\alpha > 0.7$ or $0.75$), the hash table performs \textbf{rehashing}:
\begin{enumerate}
    \item Allocate a new array of roughly double the capacity (typically the smallest prime number $> 2m$).
    \item Recompute hash indices for all existing elements using the new capacity: $h_{\text{new}}(k) = k \bmod m_{\text{new}}$.
    \item Insert every element into the new table.
    \item Release old table memory.
\end{enumerate}
Although a rehashing operation costs $O(n)$ time, because it occurs infrequently, its amortized cost per insertion remains $O(1)$.

\section{Deletion in Open Addressing: Tombstones}

In open addressing, you \textbf{cannot} simply set a deleted slot to empty (e.g., \texttt{NULL} or $0$).

\textbf{The Broken Chain Problem:} When searching for a key, the algorithm continues probing until it finds the key or hits an \textbf{empty slot} (which signals that the key was never inserted). If an intervening deleted key is simply erased to empty, subsequent searches for keys that probed past that slot will terminate prematurely and return a false "Not Found".

\subsection{The Solution: Tombstones (Deleted Marker)}
Instead of marking the slot empty, place a special sentinel value called a \textbf{tombstone} (or \texttt{DELETED} flag).
\begin{itemize}
    \item \textbf{Search Operation:} When the search encounters a tombstone, it \textbf{continues probing} past it.
    \item \textbf{Insert Operation:} An insertion can reuse and overwrite a tombstone slot once it confirms the key is not already present further down the probe sequence.
\end{itemize}

\section{Complexity of Hash Tables}

\begin{table}[htbp]
\centering
\small
\begin{tabularx}{\textwidth}{l X c}
\toprule
\textbf{Operation} & \textbf{Average / Expected Time} & \textbf{Worst-Case Time} \\
\midrule
Search & $O(1)$ (assuming uniform hashing and controlled $\alpha$) & $O(n)$ \\
Insert & $O(1)$ amortized & $O(n)$ \\
Delete & $O(1)$ expected & $O(n)$ \\
Space  & $O(n + m)$ & $O(n + m)$ \\
\bottomrule
\end{tabularx}
\caption{Hash table time and space complexity summary.}
\end{table}

\subsection{Why is Worst-Case Search $O(n)$?}
The worst-case search occurs when \textbf{all $n$ keys collide into the identical slot or bucket}. In separate chaining, this reduces to searching a singly linked list of length $n$. In open addressing, it results in probing every occupied slot in the table.

\section{C++ Implementation: Hash Table with Separate Chaining}

\begin{lstlisting}[style=cppstyle, caption={C++ Hash Table implementation using separate chaining.}]
#include <iostream>
#include <vector>
#include <list>
#include <utility>
#include <string>
using namespace std;

class HashTable {
private:
    vector<list<pair<int, string>>> table;
    int capacity;

    size_t hashFunction(int key) const {
        // Convert negative keys into a valid nonnegative index
        long long rem = key % capacity;
        if (rem < 0) rem += capacity;
        return static_cast<size_t>(rem);
    }

public:
    explicit HashTable(int size = 10) : capacity(size), table(size) {}

    void insert(int key, const string& value) {
        size_t idx = hashFunction(key);
        for (auto& entry : table[idx]) {
            if (entry.first == key) {
                entry.second = value; // Key exists: update value
                return;
            }
        }
        table[idx].emplace_back(key, value); // Key absent: append
    }

    bool search(int key, string& outValue) const {
        size_t idx = hashFunction(key);
        for (const auto& entry : table[idx]) {
            if (entry.first == key) {
                outValue = entry.second;
                return true;
            }
        }
        return false;
    }

    bool remove(int key) {
        size_t idx = hashFunction(key);
        for (auto it = table[idx].begin(); it != table[idx].end(); ++it) {
            if (it->first == key) {
                table[idx].erase(it);
                return true;
            }
        }
        return false;
    }
};

int main() {
    HashTable ht(7);
    ht.insert(10, "Dhaka");
    ht.insert(17, "Chittagong");
    ht.insert(24, "Rajshahi");

    string val;
    if (ht.search(17, val)) {
        cout << "Key 17 found: " << val << '\n';
    } else {
        cout << "Key 17 not found\n";
    }

    ht.remove(17);
    if (!ht.search(17, val)) {
        cout << "Key 17 deleted successfully\n";
    }
    return 0;
}
\end{lstlisting}

\subsection{Implementation Notes}
\begin{itemize}
    \item \textbf{Negative Key Safety:} Using modulo arithmetic directly on negative integers in C++ can yield negative remainders; adding capacity ensures valid array indexing.
    \item \textbf{Collision Handling:} Elements hashing to the same slot are appended to that slot's \texttt{std::list}.
    \item \textbf{Complexity:} Under uniform hashing, expected search, insertion, and deletion are $O(1)$.
\end{itemize}

\newpage
\begin{center}
\vspace{0.5em}
{\color{primary}\rule{0.45\textwidth}{1.2pt}}\\[4pt]
{\Large\bfseries\color{primary} Section III: Adjacency List vs Adjacency Matrix}\\[2pt]
{\color{primary}\rule{0.45\textwidth}{1.2pt}}
\vspace{0.5em}
\end{center}

\section{Representing a Graph in Memory}

Graphs can be represented in memory primarily through two classical data structures:
\begin{enumerate}
    \item \textbf{Adjacency Matrix:} A 2D array of size $V \times V$.
    \item \textbf{Adjacency List:} An array of $V$ linked lists or dynamic vectors.
\end{enumerate}

\section{Adjacency Matrix}

An \textbf{adjacency matrix} is a 2D square matrix $A$ of dimensions $V \times V$, where vertices are indexed from $0$ to $V-1$.
\[
A[u][v] = 
\begin{cases} 
1 & \text{if an edge exists between } u \text{ and } v \\
0 & \text{otherwise}
\end{cases}
\]

\subsection{Adjacency Matrix Example}
Consider an undirected graph with vertices $\{A, B, C, D\}$ and edges $(A, B)$, $(A, C)$, $(B, D)$:

\begin{table}[htbp]
\centering
\small
\begin{tabular}{c c c c c}
\toprule
\textbf{Vertex} & \textbf{A} & \textbf{B} & \textbf{C} & \textbf{D} \\
\midrule
\textbf{A} & 0 & 1 & 1 & 0 \\
\textbf{B} & 1 & 0 & 0 & 1 \\
\textbf{C} & 1 & 0 & 0 & 0 \\
\textbf{D} & 0 & 1 & 0 & 0 \\
\bottomrule
\end{tabular}
\caption{Adjacency matrix for an undirected graph (symmetric across main diagonal).}
\end{table}

\textbf{Key Observations:}
\begin{itemize}
    \item For undirected graphs, the adjacency matrix is \textbf{symmetric} ($A[u][v] = A[v][u]$).
    \item The main diagonal contains all $0$s if the graph has no self-loops.
    \item The degree of vertex $u$ is the sum of entries in row $u$.
\end{itemize}

\subsection{Adjacency Matrix for a Directed Graph}
In a directed graph, $A[u][v] = 1$ denotes a directed edge from $u \to v$. The matrix is generally asymmetric.

\begin{table}[htbp]
\centering
\small
\begin{tabular}{c c c c}
\toprule
\textbf{Vertex} & \textbf{A} & \textbf{B} & \textbf{C} \\
\midrule
\textbf{A} & 0 & 1 & 0 \\
\textbf{B} & 0 & 0 & 1 \\
\textbf{C} & 1 & 0 & 0 \\
\bottomrule
\end{tabular}
\caption{Adjacency matrix for a directed graph with arcs $A \to B, B \to C, C \to A$.}
\end{table}
\begin{itemize}
    \item \textbf{Out-degree of $u$:} Sum of entries in row $u$.
    \item \textbf{In-degree of $v$:} Sum of entries in column $v$.
\end{itemize}

\subsection{Weighted Adjacency Matrix}
For weighted graphs, store edge weight $w(u, v)$ instead of $1$. If no edge exists, store $\infty$ (or $0$ if weights are strictly positive and distinct from non-edges).

\subsection{Complexity of an Adjacency Matrix}

\begin{table}[htbp]
\centering
\small
\begin{tabularx}{\textwidth}{X c}
\toprule
\textbf{Operation} & \textbf{Complexity} \\
\midrule
Space Complexity & $O(V^2)$ \\
Check if edge $(u, v)$ exists & $O(1)$ \\
Insert an edge $(u, v)$ & $O(1)$ \\
Delete an edge $(u, v)$ & $O(1)$ \\
Find all neighbors of vertex $u$ & $O(V)$ \\
Add a new vertex & $O(V^2)$ (reallocating/copying 2D array) \\
\bottomrule
\end{tabularx}
\caption{Operational complexities of an adjacency matrix.}
\end{table}

\subsection{Advantages and Disadvantages}
\begin{itemize}
    \item \textbf{Advantages:} $O(1)$ instantaneous edge queries; easy to implement; compact representation for dense graphs.
    \item \textbf{Disadvantages:} $O(V^2)$ memory waste when graphs are sparse; iterating neighbors always takes $O(V)$ even if a vertex has only 1 edge.
\end{itemize}

\section{Adjacency List}

An \textbf{adjacency list} represents a graph using an array of $V$ lists. Entry \texttt{adj[u]} contains a list of all vertices adjacent to vertex $u$.

\begin{table}[htbp]
\centering
\small
\begin{tabular}{c l}
\toprule
\textbf{Vertex} & \textbf{Adjacent Neighbors} \\
\midrule
A & B, C \\
B & A, D \\
C & A \\
D & B \\
\bottomrule
\end{tabular}
\caption{Adjacency list for the undirected graph example.}
\end{table}

\subsection{Directed Adjacency List}
For a directed graph, each list contains only outgoing neighbors ($u \to v$):

\begin{table}[htbp]
\centering
\small
\begin{tabular}{c l}
\toprule
\textbf{Vertex} & \textbf{Outgoing Neighbors} \\
\midrule
A & B \\
B & C \\
C & A \\
\bottomrule
\end{tabular}
\caption{Adjacency list for directed cycle $A \to B \to C \to A$.}
\end{table}

\subsection{Complexity of an Adjacency List}

\begin{table}[htbp]
\centering
\small
\begin{tabularx}{\textwidth}{X p{7cm}}
\toprule
\textbf{Operation} & \textbf{Complexity} \\
\midrule
Space Complexity & $O(V + E)$ ($V$ heads $+ 2E$ list nodes for undirected) \\
Check if edge $(u, v)$ exists & $O(\deg(u))$ \\
Insert an edge $(u, v)$ & $O(1)$ (insert at head without uniqueness check) \\
Delete an edge $(u, v)$ & $O(\deg(u))$ (traverse list to find node) \\
Find all neighbors of vertex $u$ & $O(\deg(u))$ \\
Add a new vertex & $O(1)$ amortized (appending new list) \\
\bottomrule
\end{tabularx}
\caption{Operational complexities of an adjacency list.}
\end{table}

\subsection{Advantages and Disadvantages}
\begin{itemize}
    \item \textbf{Advantages:} Optimal $O(V + E)$ memory for sparse graphs; iterates neighbors in optimal $O(\deg(u))$ time; adding vertices is trivial.
    \item \textbf{Disadvantages:} Checking if edge $(u, v)$ exists takes $O(\deg(u))$ rather than $O(1)$; pointer overhead in linked lists.
\end{itemize}

\section{Adjacency List vs Adjacency Matrix: Complete Comparison}

\begin{table}[htbp]
\centering
\small
\begin{tabularx}{\textwidth}{p{3.2cm} X X}
\toprule
\textbf{Evaluation Basis} & \textbf{Adjacency Matrix} & \textbf{Adjacency List} \\
\midrule
Internal Representation & $V \times V$ 2D array & Array of $V$ lists / dynamic vectors \\
Space Complexity & $O(V^2)$ & $O(V + E)$ \\
Edge Existence Test & $O(1)$ instant lookup & $O(\deg(u))$ linear scan in list \\
Iterating Neighbors & $O(V)$ scan across entire row & $O(\deg(u))$ optimal traversal \\
Edge Insertion & $O(1)$ & $O(1)$ \\
Edge Deletion & $O(1)$ & $O(\deg(u))$ \\
Vertex Insertion & $O(V^2)$ (reallocate matrix) & $O(1)$ amortized \\
Best Suited For & \textbf{Dense graphs} ($E \approx V^2$) & \textbf{Sparse graphs} ($E \ll V^2$) \\
Memory Waste in Sparse & Extremely high ($0$s dominate) & Minimal \\
Traversal (BFS/DFS) & $O(V^2)$ & $O(V + E)$ \\
\bottomrule
\end{tabularx}
\caption{Master comparison between adjacency matrix and adjacency list.}
\end{table}

\subsection{Worked Memory Comparison}

\begin{tcolorbox}[title=\textbf{Worked Example: Memory Requirements}]
Suppose a social network graph has $V = 1,000$ vertices and $E = 2,000$ edges. Compare storage requirements:
\begin{itemize}
    \item \textbf{Adjacency Matrix:} Requires $V^2 = 1,000^2 = 1,000,000$ cells.
    \item \textbf{Adjacency List:} Structural size is proportional to $V + E = 1,000 + 2,000 = 3,000$ elements (or $1,000 + 4,000 = 5,000$ directed node instances).
    \item \textbf{Maximum possible edges:}
    \[
    E_{\max} = \frac{1000 \times 999}{2} = 499,500 \text{ edges.}
    \]
    Since $2,000 \ll 499,500$, the graph is extremely sparse ($0.4\%$ density). An adjacency list saves over $99\%$ of memory!
\end{itemize}
\end{tcolorbox}

\subsection{Worked Example: Choosing the Right Representation}
\begin{itemize}
    \item \textbf{Scenario A:} A graph has 10,000 vertices but only 15,000 edges. $\implies$ \textbf{Choose Adjacency List} (sparse, matrix would need $10^8$ cells).
    \item \textbf{Scenario B:} A 200-node graph where nearly all pairs are connected and millions of edge-existence checks occur. $\implies$ \textbf{Choose Adjacency Matrix} (dense, $O(1)$ edge test).
    \item \textbf{Scenario C:} A pathfinding algorithm frequently traverses neighbors. $\implies$ \textbf{Choose Adjacency List} ($O(\deg(u))$ vs $O(V)$ scan).
\end{itemize}

\section{Graph Traversal: BFS and DFS}

\subsection{Breadth-First Search (BFS)}
\textbf{BFS} explores vertices level by level in increasing order of their distance from the source vertex.
\begin{itemize}
    \item \textbf{Data Structure:} Uses a \textbf{Queue} (FIFO) and a \texttt{visited[]} array.
    \item \textbf{Algorithm:} Mark start visited and enqueue; while queue is not empty, dequeue vertex $u$, process it, and enqueue all unvisited neighbors.
    \item \textbf{Key Feature:} Finds the \textbf{shortest path} (minimum edges) in unweighted graphs.
\end{itemize}

\subsection{Depth-First Search (DFS)}
\textbf{DFS} explores as deeply as possible along each branch before backtracking.
\begin{itemize}
    \item \textbf{Data Structure:} Uses a \textbf{Stack} (LIFO) or \textbf{System Call Stack} (recursion) and a \texttt{visited[]} array.
    \item \textbf{Algorithm:} Mark current vertex visited, process it, and recursively visit all unvisited neighbors.
    \item \textbf{Key Features:} Cycle detection, topological sorting in DAGs, connected components, finding strongly connected components (Kosaraju/Tarjan).
\end{itemize}

\subsection{Traversal Complexity}
\begin{itemize}
    \item With \textbf{Adjacency List:} $\boxed{O(V + E)}$ time complexity.
    \item With \textbf{Adjacency Matrix:} $\boxed{O(V^2)}$ time complexity.
    \item Space complexity for both: $O(V)$ for the visited array and recursion stack / queue.
\end{itemize}

\section{C++ Implementation of Adjacency List and Matrix}

\subsection{Adjacency List Implementation}

\begin{lstlisting}[style=cppstyle, caption={C++ Graph implementation using Adjacency List.}]
#include <iostream>
#include <vector>
using namespace std;

class GraphList {
private:
    int V;
    vector<vector<int>> adj;

public:
    explicit GraphList(int vertices) : V(vertices), adj(vertices) {}

    void addEdge(int u, int v) {
        if (u < 0 || u >= V || v < 0 || v >= V) return;
        adj[u].push_back(v);
        adj[v].push_back(u); // Bidirectional for undirected graph
    }

    void display() const {
        for (int i = 0; i < V; ++i) {
            cout << i << ": ";
            for (int neighbor : adj[i]) {
                cout << neighbor << ' ';
            }
            cout << '\n';
        }
    }
};

int main() {
    GraphList g(4);
    g.addEdge(0, 1);
    g.addEdge(0, 2);
    g.addEdge(1, 3);
    g.display();
    return 0;
}
\end{lstlisting}

\begin{lstlisting}[style=pseudocode, caption={Output of Adjacency List program.}]
0: 1 2 
1: 0 3 
2: 0 
3: 1 
\end{lstlisting}

\subsection{Adjacency Matrix Implementation}

\begin{lstlisting}[style=cppstyle, caption={C++ Graph implementation using Adjacency Matrix.}]
#include <iostream>
#include <vector>
using namespace std;

class GraphMatrix {
private:
    int V;
    vector<vector<int>> matrix;

public:
    explicit GraphMatrix(int vertices)
        : V(vertices), matrix(vertices, vector<int>(vertices, 0)) {}

    void addEdge(int u, int v) {
        if (u < 0 || u >= V || v < 0 || v >= V) return;
        matrix[u][v] = 1;
        matrix[v][u] = 1; // Symmetric for undirected graph
    }

    bool hasEdge(int u, int v) const {
        if (u < 0 || u >= V || v < 0 || v >= V) return false;
        return matrix[u][v] == 1;
    }

    void display() const {
        for (int i = 0; i < V; ++i) {
            for (int j = 0; j < V; ++j) {
                cout << matrix[i][j] << ' ';
            }
            cout << '\n';
        }
    }
};

int main() {
    GraphMatrix g(4);
    g.addEdge(0, 1);
    g.addEdge(0, 2);
    g.addEdge(1, 3);

    g.display();
    cout << "Edge 0-2 exists: " << (g.hasEdge(0, 2) ? "Yes" : "No") << '\n';
    return 0;
}
\end{lstlisting}

\begin{lstlisting}[style=pseudocode, caption={Output of Adjacency Matrix program.}]
0 1 1 0 
1 0 0 1 
1 0 0 0 
0 1 0 0 
Edge 0-2 exists: Yes
\end{lstlisting}

\newpage
\begin{center}
\vspace{0.5em}
{\color{primary}\rule{0.45\textwidth}{1.2pt}}\\[4pt]
{\Large\bfseries\color{primary} Section IV: Memorization Facts, Examination Q\&A, and Quick Revision}\\[2pt]
{\color{primary}\rule{0.45\textwidth}{1.2pt}}
\vspace{0.5em}
\end{center}

\section{Additional Graph Properties Worth Memorizing}

\begin{enumerate}
    \item \textbf{A tree is connected and acyclic:} A tree with $V$ vertices has exactly $V - 1$ edges.
    \item \textbf{A forest is an acyclic undirected graph:} If it contains $C$ connected components, the total number of edges is $E = V - C$.
    \item \textbf{Every tree is bipartite:} Its vertices can be partitioned into two sets by distance parity (even/odd depth) from the root.
    \item \textbf{Bipartite Odd-Cycle Rule:} A graph is bipartite if and only if it contains \textbf{no odd cycles}.
    \item \textbf{Minimum edges for connectivity:} A connected undirected graph with $V$ vertices requires at least $V - 1$ edges.
    \item \textbf{Sum of degrees is always even:} In any undirected graph, $\sum \deg(v) = 2E$, which is an even integer.
    \item \textbf{Directed degree invariant:} For any directed graph, $\sum \text{in-deg}(v) = \sum \text{out-deg}(v) = E$.
    \item \textbf{Complete simple undirected graph edges:} $K_n$ contains exactly $\frac{n(n-1)}{2}$ edges.
    \item \textbf{DAG implies topological sorting:} Every directed acyclic graph has at least one valid topological sort.
    \item \textbf{Eulerian vs Hamiltonian:} Eulerian traversal visits every \textbf{edge} once; Hamiltonian traversal visits every \textbf{vertex} once.
    \item \textbf{Strong connectivity is for directed graphs:} Undirected graphs are simply connected or disconnected.
    \item \textbf{Shortest path guarantee:} BFS guarantees the shortest path (minimum edge count) in unweighted graphs. For weighted graphs with non-negative edges, Dijkstra's algorithm must be used.
\end{enumerate}

\section{Likely Recruitment Examination Questions: Graphs and Hashing}

\subsection{Questions on Graphs}

\begin{table}[htbp]
\centering
\small
\begin{tabularx}{\textwidth}{p{8cm} X}
\toprule
\textbf{Recruitment Examination Question} & \textbf{Direct Technical Answer} \\
\midrule
What is the maximum number of edges in a simple undirected graph with $n$ vertices? & $\frac{n(n-1)}{2}$ \\
What is the maximum number of arcs in a simple directed graph without self-loops? & $n(n-1)$ \\
What is the sum of degrees in an undirected graph? & $2E$ (Handshaking Lemma) \\
What is the sum of in-degrees in a directed graph? & $E$ (Total number of directed edges) \\
How many edges does a tree with $n$ vertices have? & $n - 1$ \\
How many edges does a complete bipartite graph $K_{a,b}$ have? & $a \times b$ \\
When does an undirected graph have an Euler circuit? & When connected and \textbf{every vertex has an even degree} \\
When does an undirected graph have an Euler trail? & When connected and \textbf{exactly two vertices have odd degree} \\
What is the key defining property of a DAG? & Directed Acyclic Graph (admits topological sort) \\
What is the time complexity of BFS using adjacency list? & $O(V + E)$ \\
What data structure is used to implement BFS? & \textbf{Queue} \\
What data structure is used to implement DFS? & \textbf{Stack} (or recursion call stack) \\
\bottomrule
\end{tabularx}
\caption{Key examination questions on graphs.}
\end{table}

\subsection{Questions on Hash Tables}

\begin{table}[htbp]
\centering
\small
\begin{tabularx}{\textwidth}{p{8cm} X}
\toprule
\textbf{Recruitment Examination Question} & \textbf{Direct Technical Answer} \\
\midrule
What is a hash collision? & Two distinct keys producing identical hash index ($h(k_1) = h(k_2)$) \\
What is the division hash function formula? & $h(k) = k \bmod m$ \\
Why should table size $m$ be a prime number? & Reduces collisions caused by common factors \\
What is the definition of load factor $\alpha$? & $\alpha = n / m$ (entries / capacity) \\
Can the load factor exceed 1 in separate chaining? & \textbf{Yes} (each bucket can hold multiple items) \\
Can the load factor exceed 1 in open addressing? & \textbf{No} ($\alpha \leq 1$, all items stored inside array) \\
What is primary clustering in linear probing? & Consecutive occupied slots grouping into large blocks \\
What formula defines quadratic probing? & $h(k, i) = (h'(k) + c_1 i + c_2 i^2) \bmod m$ \\
What formula defines double hashing? & $h(k, i) = (h_1(k) + i \cdot h_2(k)) \bmod m$ \\
Why can't $h_2(k)$ evaluate to zero in double hashing? & It would produce an infinite loop at index $h_1(k)$ \\
Why are tombstones used during deletion? & To prevent breaking the probe chain for subsequent lookups \\
What is expected search complexity in a hash table? & $O(1)$ under uniform distribution and low $\alpha$ \\
What is worst-case search complexity in a hash table? & $O(n)$ (when all keys collide into one slot) \\
What is rehashing? & Resizing the table (typically $2\times$) and re-inserting all keys \\
\bottomrule
\end{tabularx}
\caption{Key examination questions on hash tables.}
\end{table}

\subsection{Questions on Graph Representation}

\begin{table}[htbp]
\centering
\small
\begin{tabularx}{\textwidth}{p{8cm} X}
\toprule
\textbf{Recruitment Examination Question} & \textbf{Direct Technical Answer} \\
\midrule
What is the space complexity of an adjacency matrix? & $O(V^2)$ \\
What is the space complexity of an adjacency list? & $O(V + E)$ \\
What is the time complexity to check if edge $(u,v)$ exists in matrix? & $O(1)$ \\
What is the time complexity to check if edge $(u,v)$ exists in list? & $O(\deg(u))$ \\
What is the time complexity to find all neighbors in matrix? & $O(V)$ \\
What is the time complexity to find all neighbors in list? & $O(\deg(u))$ \\
Which representation is preferred for sparse graphs? & \textbf{Adjacency List} (saves massive memory) \\
Which representation is preferred for dense graphs? & \textbf{Adjacency Matrix} ($O(1)$ lookup, compact) \\
What is the time complexity of BFS using adjacency matrix? & $O(V^2)$ \\
Why is an undirected adjacency matrix symmetric? & Because $(u, v)$ and $(v, u)$ represent the identical edge \\
\bottomrule
\end{tabularx}
\caption{Key examination questions on graph representations.}
\end{table}

\section{Final Quick-Revision Sheet: Graphs and Hashing}

\subsection{Graphs Summary}
\begin{itemize}
    \item \textbf{Definition:} $G = (V, E)$.
    \item \textbf{Handshaking Lemma:} $\sum \deg(v) = 2E$. Odd degree vertex count is always even.
    \item \textbf{Directed Degrees:} $\sum \text{in-deg}(v) = \sum \text{out-deg}(v) = E$.
    \item \textbf{Max Edges (Undirected):} $\frac{V(V-1)}{2}$.
    \item \textbf{Max Edges (Directed, no loops):} $V(V-1)$.
    \item \textbf{Tree:} Connected, acyclic graph with $E = V - 1$.
    \item \textbf{Forest:} Acyclic graph with $C$ components: $E = V - C$.
    \item \textbf{Complete Bipartite $K_{a,b}$:} $E = a \times b$.
    \item \textbf{Planar Euler Formula:} $V - E + F = 2$.
    \item \textbf{DAG:} Admits topological sorting; no directed cycles.
    \item \textbf{Euler vs Hamiltonian:} Euler = edges; Hamiltonian = vertices.
    \item \textbf{Traversals:} BFS uses Queue ($O(V+E)$ list, $O(V^2)$ matrix); DFS uses Stack ($O(V+E)$ list, $O(V^2)$ matrix).
\end{itemize}

\subsection{Hash Tables Summary}
\begin{itemize}
    \item \textbf{Hash Function:} Deterministic $O(1)$ mapping from key space to table index $[0, m-1]$.
    \item \textbf{Load Factor:} $\alpha = n / m$.
    \item \textbf{Division Method:} $h(k) = k \bmod m$ (use prime $m$).
    \item \textbf{Collision Resolution:}
    \begin{itemize}
        \item \textbf{Separate Chaining:} Buckets outside array; $\alpha$ can exceed 1; simple deletion.
        \item \textbf{Linear Probing:} $h(k, i) = (h'(k) + i) \bmod m$; suffers from primary clustering.
        \item \textbf{Quadratic Probing:} $h(k, i) = (h'(k) + i^2) \bmod m$; suffers from secondary clustering.
        \item \textbf{Double Hashing:} $h(k, i) = (h_1(k) + i \cdot h_2(k)) \bmod m$; best open addressing.
    \end{itemize}
    \item \textbf{Tombstones:} Required in open addressing deletion to maintain unbroken probe chains.
    \item \textbf{Complexities:} Expected $O(1)$ search/insert/delete; worst-case $O(n)$ under total collision.
    \item \textbf{Rehashing:} Allocate $2\times$ array (prime size) and re-insert all keys when $\alpha > 0.75$.
\end{itemize}

\subsection{Adjacency List vs Matrix Summary}
\begin{itemize}
    \item \textbf{Adjacency Matrix:} $O(V^2)$ memory; $O(1)$ edge query; $O(V)$ neighbor search; ideal for dense graphs.
    \item \textbf{Adjacency List:} $O(V + E)$ memory; $O(\deg(u))$ edge query; $O(\deg(u))$ neighbor search; ideal for sparse graphs.
    \item \textbf{BFS/DFS Complexity:} $O(V + E)$ with list; $O(V^2)$ with matrix.
\end{itemize}

\begin{quote}
\textbf{Core Takeaway for Recruitment Exams:} Connect each theoretical definition to its mathematical formula and operational time/space costs. For graphs, memorize handshaking degree sums, maximum edges, and tree formulas. For hash tables, master collision-resolution mechanics, load factor constraints, and tombstones. For graph representations, select matrices for dense edge lookups and lists for sparse neighbor traversals.
\end{quote}


\begin{center}
\vspace{1em}
{\color{primary}\rule{0.5\textwidth}{1.5pt}}\\[6pt]
{\Large\bfseries\color{primary} Part VI: Data Structure Recruitment Examination Problem Bank}\\[3pt]
{\color{primary}\rule{0.5\textwidth}{1.5pt}}
\vspace{1em}
\end{center}

This chapter brings together essential practical programming questions, algorithm-tracing problems, and mathematical calculations from across the Data Structures syllabus. The problems are selected to reflect the testing patterns of Bangladesh government technical recruitment examinations, including Assistant Programmer, Assistant Director (ICT), Assistant Network Engineer, Assistant Maintenance Engineer, Assistant Database Administrator, and related IT cadres.

Each problem includes the problem statement, underlying algorithmic reasoning, complete worked solution, time and space complexity analysis, and common examination pitfalls.

\begin{center}
\vspace{0.5em}
{\color{primary}\rule{0.45\textwidth}{1.2pt}}\\[4pt]
{\Large\bfseries\color{primary} Section I: Arrays and Basic Algorithms}\\[2pt]
{\color{primary}\rule{0.45\textwidth}{1.2pt}}
\vspace{0.5em}
\end{center}

\section{Arrays and Basic Algorithms}

\subsection{Problem 1: Find the Smallest Element in an Array}

\textbf{Problem Statement:} Given an array of $n$ integers, find the smallest element without sorting the array.

\begin{lstlisting}[style=pseudocode, caption={Input}]
6
25 12 8 34 5 19
\end{lstlisting}

\begin{lstlisting}[style=pseudocode, caption={Expected Output}]
Smallest element: 5
\end{lstlisting}

\subsubsection*{Algorithmic Strategy}
Initialize the running minimum with the first element $A[0]$. Traverse the array from index $1$ to $n-1$, comparing each element against the running minimum. Whenever an element $A[i] < \text{minimum}$ is encountered, update $\text{minimum} = A[i]$.

\begin{lstlisting}[style=cppstyle, caption={C++ Solution for Finding Minimum Element.}]
#include <iostream>
#include <vector>
using namespace std;

int main() {
    int n;
    if (!(cin >> n) || n <= 0) {
        cout << "Array must not be empty\n";
        return 0;
    }

    vector<int> a(n);
    for (int i = 0; i < n; i++) {
        cin >> a[i];
    }

    int minimum = a[0];
    for (int i = 1; i < n; i++) {
        if (a[i] < minimum) {
            minimum = a[i];
        }
    }

    cout << "Smallest element: " << minimum << '\n';
    return 0;
}
\end{lstlisting}

\subsubsection*{Step-by-Step Execution Trace}
\begin{table}[htbp]
\centering
\small
\begin{tabular}{c c c}
\toprule
\textbf{Step} & \textbf{Element Examined ($A[i]$)} & \textbf{Current Minimum} \\
\midrule
$i = 0$ & 25 (Initial) & 25 \\
$i = 1$ & 12 ($12 < 25$) & 12 \\
$i = 2$ & 8 ($8 < 12$) & 8 \\
$i = 3$ & 34 ($34 \geq 8$) & 8 \\
$i = 4$ & 5 ($5 < 8$) & 5 \\
$i = 5$ & 19 ($19 \geq 5$) & 5 \\
\bottomrule
\end{tabular}
\caption{Linear scan trace for finding the minimum element.}
\end{table}

\textbf{Complexity Analysis:}
\begin{itemize}
    \item \textbf{Time Complexity:} $O(n)$, since every element is examined exactly once.
    \item \textbf{Auxiliary Space:} $O(1)$, requiring only a single scalar variable.
\end{itemize}

\begin{quote}
\textbf{Exam Note:} Sorting the array first takes $O(n \log n)$ time and is unnecessary. Linear scanning in $O(n)$ is asymptotically optimal because any adversary could place the minimum in the unexamined position.
\end{quote}

\subsection{Problem 2: Find the Second-Largest Distinct Element}

\textbf{Problem Statement:} Find the second-largest \textbf{distinct} value in an integer array without sorting it.

\begin{lstlisting}[style=pseudocode, caption={Input}]
7
10 25 8 25 17 30 30
\end{lstlisting}

\begin{lstlisting}[style=pseudocode, caption={Expected Output}]
Second largest: 25
\end{lstlisting}

\subsubsection*{Algorithmic Strategy}
Maintain two variables: \texttt{largest} and \texttt{secondLargest}, initialized to $-\infty$.
\begin{itemize}
    \item If $A[i] > \text{largest}$: update $\text{secondLargest} = \text{largest}$, then $\text{largest} = A[i]$.
    \item Else if $A[i] < \text{largest}$ and $A[i] > \text{secondLargest}$: update $\text{secondLargest} = A[i]$.
    \item If $A[i] == \text{largest}$: ignore duplicates of the largest value.
\end{itemize}

\begin{lstlisting}[style=cppstyle, caption={C++ Solution for Finding Second-Largest Distinct Element.}]
#include <iostream>
#include <vector>
#include <climits>
using namespace std;

int main() {
    int n;
    if (!(cin >> n) || n < 2) {
        cout << "Array must contain at least two elements\n";
        return 0;
    }

    vector<int> a(n);
    for (int i = 0; i < n; i++) {
        cin >> a[i];
    }

    long long largest = LLONG_MIN;
    long long secondLargest = LLONG_MIN;
    bool hasLargest = false;
    bool hasSecond = false;

    for (int x : a) {
        if (!hasLargest || x > largest) {
            if (hasLargest) {
                secondLargest = largest;
                hasSecond = true;
            }
            largest = x;
            hasLargest = true;
        } else if (x < largest && (!hasSecond || x > secondLargest)) {
            secondLargest = x;
            hasSecond = true;
        }
    }

    if (hasSecond) {
        cout << "Second largest: " << secondLargest << '\n';
    } else {
        cout << "No second-largest distinct value\n";
    }
    return 0;
}
\end{lstlisting}

\textbf{Complexity Analysis:}
\begin{itemize}
    \item \textbf{Time Complexity:} $O(n)$ in a single pass.
    \item \textbf{Auxiliary Space:} $O(1)$.
\end{itemize}

\subsection{Problem 3: Binary Search on a Sorted Array}

\textbf{Problem Statement:} Search for target value $T = 23$ in the sorted array:
\[
A = [3, 8, 12, 17, 23, 31, 42]
\]

\subsubsection*{Algorithmic Strategy}
Maintain search boundaries $[\text{left}, \text{right}]$. Compute midpoint:
\[
\text{mid} = \text{left} + \left\lfloor \frac{\text{right} - \text{left}}{2} \right\rfloor
\]
\begin{enumerate}
    \item Initially: $\text{left} = 0, \text{right} = 6 \implies \text{mid} = 3$, $A[3] = 17$. Since $23 > 17$, eliminate left half: $\text{left} = 4$.
    \item New interval: $[4, 6] \implies \text{mid} = 4 + 1 = 5$, $A[5] = 31$. Since $23 < 31$, eliminate right half: $\text{right} = 4$.
    \item New interval: $[4, 4] \implies \text{mid} = 4$, $A[4] = 23$. Target found at index $4$!
\end{enumerate}

\begin{lstlisting}[style=cppstyle, caption={C++ Iterative Binary Search.}]
#include <iostream>
#include <vector>
using namespace std;

int binarySearch(const vector<int>& a, int target) {
    int left = 0;
    int right = static_cast<int>(a.size()) - 1;

    while (left <= right) {
        int mid = left + (right - left) / 2; // Prevents integer overflow
        if (a[mid] == target) return mid;
        if (a[mid] < target)  left = mid + 1;
        else                  right = mid - 1;
    }
    return -1; // Target not found
}

int main() {
    vector<int> a = {3, 8, 12, 17, 23, 31, 42};
    int index = binarySearch(a, 23);
    cout << "Target 23 found at index: " << index << '\n';
    return 0;
}
\end{lstlisting}

\textbf{Complexity Analysis:}
\begin{itemize}
    \item \textbf{Time Complexity:} $O(\log n)$ worst and average case; $O(1)$ best case.
    \item \textbf{Auxiliary Space:} $O(1)$ for iterative; $O(\log n)$ for recursive call stack.
\end{itemize}

\begin{quote}
\textbf{Exam Trap:} Binary search \textbf{strictly requires} a sorted collection. It cannot be applied directly to an unsorted array without prior sorting.
\end{quote}

\newpage
\begin{center}
\vspace{0.5em}
{\color{primary}\rule{0.45\textwidth}{1.2pt}}\\[4pt]
{\Large\bfseries\color{primary} Section II: Stack and Queue Problems}\\[2pt]
{\color{primary}\rule{0.45\textwidth}{1.2pt}}
\vspace{0.5em}
\end{center}

\section{Stack and Queue Problems}

\subsection{Problem 4: Balanced Parentheses Checking}

\textbf{Problem Statement:} Determine whether the following bracket expressions are balanced:
\begin{itemize}
    \item Expression 1: \texttt{\{[a + (b * c)]\}} $\implies$ \textbf{Balanced}
    \item Expression 2: \texttt{\{[a + (b * c]\}} $\implies$ \textbf{Not balanced}
\end{itemize}

\subsubsection*{Algorithmic Strategy}
Use a stack to record opening brackets:
\begin{enumerate}
    \item When an opening delimiter (\texttt{(}, \texttt{\{}, \texttt{[}) is read, push it onto the stack.
    \item When a closing delimiter (\texttt{)}, \texttt{\}}, \texttt{]}) is read:
    \begin{itemize}
        \item If stack is empty, return false (unmatched closing bracket).
        \item Pop top element; if it does not match the closing bracket's type, return false.
    \end{itemize}
    \item At the end of the string, the expression is balanced if and only if the stack is completely empty.
\end{enumerate}

\begin{lstlisting}[style=cppstyle, caption={C++ Solution for Balanced Parentheses.}]
#include <iostream>
#include <stack>
#include <string>
using namespace std;

bool isMatching(char opening, char closing) {
    return (opening == '(' && closing == ')') ||
           (opening == '[' && closing == ']') ||
           (opening == '{' && closing == '}');
}

bool isBalanced(const string& expression) {
    stack<char> st;
    for (char ch : expression) {
        if (ch == '(' || ch == '[' || ch == '{') {
            st.push(ch);
        } else if (ch == ')' || ch == ']' || ch == '}') {
            if (st.empty() || !isMatching(st.top(), ch)) {
                return false;
            }
            st.pop();
        }
    }
    return st.empty();
}

int main() {
    string exp1 = "{[a + (b * c)]}";
    string exp2 = "{[a + (b * c]}";

    cout << exp1 << " -> " << (isBalanced(exp1) ? "Balanced" : "Not balanced") << '\n';
    cout << exp2 << " -> " << (isBalanced(exp2) ? "Balanced" : "Not balanced") << '\n';
    return 0;
}
\end{lstlisting}

\textbf{Complexity Analysis:}
\begin{itemize}
    \item \textbf{Time Complexity:} $O(n)$ where $n$ is string length.
    \item \textbf{Auxiliary Space:} $O(n)$ in the worst case (e.g., all opening brackets).
\end{itemize}

\subsection{Problem 5: Postfix Expression Evaluation}

\textbf{Problem Statement:} Evaluate the postfix (Reverse Polish Notation) expression:
\[
5 \quad 3 \quad 2 \quad * \quad +
\]

\subsubsection*{Step-by-Step Evaluation Trace}
\begin{enumerate}
    \item Push operand $5$: Stack $= [5]$.
    \item Push operand $3$: Stack $= [5, 3]$.
    \item Push operand $2$: Stack $= [5, 3, 2]$.
    \item Operator $*$ encountered:
    \begin{itemize}
        \item Pop $\text{op}_2 = 2$, Pop $\text{op}_1 = 3$.
        \item Compute: $\text{result} = 3 \times 2 = 6$.
        \item Push $6$: Stack $= [5, 6]$.
    \end{itemize}
    \item Operator $+$ encountered:
    \begin{itemize}
        \item Pop $\text{op}_2 = 6$, Pop $\text{op}_1 = 5$.
        \item Compute: $\text{result} = 5 + 6 = 11$.
        \item Push $11$: Stack $= [11]$.
    \end{itemize}
\end{enumerate}
Final result:
\[
\boxed{11}
\]

\begin{quote}
\textbf{Crucial Rule for Non-Commutative Operators:} In division and subtraction, order is vital:
\[
\text{op}_2 = \text{pop}(), \quad \text{op}_1 = \text{pop}() \implies \text{compute } (\text{op}_1 - \text{op}_2) \text{ or } (\text{op}_1 / \text{op}_2).
\]
\end{quote}

\subsection{Problem 6: Implement a Queue Using Two Stacks}

\textbf{Problem Statement:} Implement FIFO queue semantics using two LIFO stacks. Trace:
\[
\text{enqueue}(10), \; \text{enqueue}(20), \; \text{enqueue}(30), \; \text{dequeue}(), \; \text{enqueue}(40), \; \text{dequeue}()
\]

\begin{lstlisting}[style=cppstyle, caption={C++ Implementation of Queue using Two Stacks.}]
#include <iostream>
#include <stack>
using namespace std;

class QueueUsingStacks {
private:
    stack<int> inStack;
    stack<int> outStack;

    void transferIfNeeded() {
        if (outStack.empty()) {
            while (!inStack.empty()) {
                outStack.push(inStack.top());
                inStack.pop();
            }
        }
    }

public:
    void enqueue(int value) {
        inStack.push(value);
    }

    bool dequeue(int& value) {
        transferIfNeeded();
        if (outStack.empty()) return false;
        value = outStack.top();
        outStack.pop();
        return true;
    }
};

int main() {
    QueueUsingStacks q;
    int val;

    q.enqueue(10);
    q.enqueue(20);
    q.enqueue(30);

    q.dequeue(val);
    cout << "Dequeued: " << val << '\n'; // 10

    q.enqueue(40);

    q.dequeue(val);
    cout << "Dequeued: " << val << '\n'; // 20

    return 0;
}
\end{lstlisting}

\textbf{Complexity Analysis:}
\begin{itemize}
    \item \textbf{Enqueue:} $O(1)$ worst case.
    \item \textbf{Dequeue:} $O(1)$ amortized; $O(n)$ worst case (when \texttt{outStack} is empty).
    \item \textbf{Space:} $O(n)$ total storage across both stacks.
\end{itemize}

\newpage
\begin{center}
\vspace{0.5em}
{\color{primary}\rule{0.45\textwidth}{1.2pt}}\\[4pt]
{\Large\bfseries\color{primary} Section III: Linked-List Problems}\\[2pt]
{\color{primary}\rule{0.45\textwidth}{1.2pt}}
\vspace{0.5em}
\end{center}

\section{Linked-List Problems}

\subsection{Problem 7: Reverse a Singly Linked List}

\textbf{Problem Statement:} Reverse the list $10 \to 20 \to 30 \to 40 \to \text{NULL}$ in place.

\subsubsection*{Algorithmic Strategy}
Maintain three pointers: \texttt{previous = nullptr}, \texttt{current = head}, and \texttt{nextNode = nullptr}.
\begin{enumerate}
    \item Save next node: $\texttt{nextNode} = \texttt{current->next}$.
    \item Reverse current pointer: $\texttt{current->next} = \texttt{previous}$.
    \item Advance previous: $\texttt{previous} = \texttt{current}$.
    \item Advance current: $\texttt{current} = \texttt{nextNode}$.
\end{enumerate}

\begin{lstlisting}[style=cppstyle, caption={C++ Solution for In-Place Linked List Reversal.}]
#include <iostream>
using namespace std;

struct Node {
    int data;
    Node* next;
    explicit Node(int val) : data(val), next(nullptr) {}
};

Node* reverseList(Node* head) {
    Node* previous = nullptr;
    Node* current = head;

    while (current != nullptr) {
        Node* nextNode = current->next;
        current->next = previous;
        previous = current;
        current = nextNode;
    }
    return previous; // New head
}

void display(Node* head) {
    while (head) {
        cout << head->data << (head->next ? " -> " : " -> NULL\n");
        head = head->next;
    }
}

int main() {
    Node* head = new Node(10);
    head->next = new Node(20);
    head->next->next = new Node(30);
    head->next->next->next = new Node(40);

    cout << "Original list: "; display(head);
    head = reverseList(head);
    cout << "Reversed list: "; display(head);

    return 0;
}
\end{lstlisting}

\textbf{Complexity Analysis:} Time $O(n)$; Auxiliary space $O(1)$.

\subsection{Problem 8: Detect a Cycle in a Linked List (Floyd's Algorithm)}

\textbf{Problem Statement:} Detect if a singly linked list contains a cycle using $O(1)$ auxiliary memory.

\subsubsection*{Floyd's Tortoise and Hare Algorithm}
Use two pointers:
\begin{itemize}
    \item \texttt{slow}: advances $1$ node per step.
    \item \texttt{fast}: advances $2$ nodes per step.
\end{itemize}
If a cycle exists, the relative distance between them decreases by $1$ in each step inside the loop, guaranteeing that \texttt{fast} will catch up with \texttt{slow} in at most $L$ steps (where $L$ is loop length). If no cycle exists, \texttt{fast} reaches \texttt{NULL}.

\begin{lstlisting}[style=cppstyle, caption={C++ Cycle Detection using Floyd's Algorithm.}]
bool hasCycle(Node* head) {
    Node* slow = head;
    Node* fast = head;

    while (fast != nullptr && fast->next != nullptr) {
        slow = slow->next;
        fast = fast->next->next;
        if (slow == fast) {
            return true; // Cycle detected
        }
    }
    return false; // Reached end: no cycle
}
\end{lstlisting}

\textbf{Complexity Analysis:} Time $O(n)$; Auxiliary space $O(1)$.

\newpage
\begin{center}
\vspace{0.5em}
{\color{primary}\rule{0.45\textwidth}{1.2pt}}\\[4pt]
{\Large\bfseries\color{primary} Section IV: Heap and Priority-Queue Problems}\\[2pt]
{\color{primary}\rule{0.45\textwidth}{1.2pt}}
\vspace{0.5em}
\end{center}

\section{Heap and Priority-Queue Problems}

\subsection{Problem 9: Build a Max-Heap from an Unordered Array}

\textbf{Problem Statement:} Construct a max-heap from array $A = [4, 10, 3, 5, 1]$.

\subsubsection*{Worked Solution}
For array of size $n = 5$, the last non-leaf node index is:
\[
\left\lfloor \frac{n}{2} \right\rfloor - 1 = \left\lfloor \frac{5}{2} \right\rfloor - 1 = 1
\]
Heapify in reverse order from index $1$ down to $0$:
\begin{itemize}
    \item \textbf{At index 1 ($A[1] = 10$):} Children are $A[3] = 5, A[4] = 1$. Since $10 > 5$ and $10 > 1$, max-heap property is already satisfied.
    \item \textbf{At index 0 ($A[0] = 4$):} Children are $A[1] = 10, A[2] = 3$. Largest child is $10$. Swap $A[0]$ and $A[1]$:
    \[
    A = [10, 4, 3, 5, 1]
    \]
    Sift down $4$ at index $1$: Children are $A[3] = 5, A[4] = 1$. Largest is $5$. Swap $A[1]$ and $A[3]$:
    \[
    \boxed{A = [10, 5, 3, 4, 1]}
    \]
\end{itemize}

\begin{quote}
\textbf{Key Complexity Theorem:} Building a heap bottom-up takes $\boxed{O(n)}$ time, not $O(n \log n)$, because the sum of heights across all nodes converges to $O(n)$:
\[
\sum_{h=0}^{\lfloor \log n \rfloor} \frac{n}{2^{h+1}} O(h) = O(n).
\]
\end{quote}

\subsection{Problem 10: Insert and Delete Maximum from a Max-Heap}

\textbf{Problem Statement:} Given max-heap $A = [10, 5, 3, 4, 1]$:
\begin{enumerate}
    \item Insert element $15$.
    \item Delete the maximum element.
\end{enumerate}

\subsubsection*{1. Insertion Trace}
Append $15$ at index $5$: $A = [10, 5, 3, 4, 1, 15]$.
\begin{itemize}
    \item Parent of index $5$ is $\lfloor (5-1)/2 \rfloor = 2$, holding value $3$. Since $15 > 3$, swap:
    \[
    A = [10, 5, 15, 4, 1, 3]
    \]
    \item Parent of index $2$ is $\lfloor (2-1)/2 \rfloor = 0$, holding value $10$. Since $15 > 10$, swap:
    \[
    \boxed{A = [15, 5, 10, 4, 1, 3]}
    \]
\end{itemize}

\subsubsection*{2. Deletion Trace}
Extract root $15$. Move last element ($3$) to root and shrink heap size:
\[
A = [3, 5, 10, 4, 1]
\]
Sift down root $3$: Children are $5$ (idx 1) and $10$ (idx 2). Larger child is $10$. Swap $3$ and $10$:
\[
\boxed{A = [10, 5, 3, 4, 1]}
\]

\subsection{Problem 11: Priority Queue Extraction Order}

\textbf{Problem Statement:} A max-priority queue contains elements $\{12, 5, 20, 8, 15\}$. In what sequence are they extracted?
\begin{itemize}
    \item \textbf{Max-Priority Queue:} Extracts in descending order of priority:
    \[
    20 \longrightarrow 15 \longrightarrow 12 \longrightarrow 8 \longrightarrow 5
    \]
    \item \textbf{Min-Priority Queue:} Extracts in ascending order:
    \[
    5 \longrightarrow 8 \longrightarrow 12 \longrightarrow 15 \longrightarrow 20
    \]
\end{itemize}

\newpage
\begin{center}
\vspace{0.5em}
{\color{primary}\rule{0.45\textwidth}{1.2pt}}\\[4pt]
{\Large\bfseries\color{primary} Section V: Circular Queue Problems}\\[2pt]
{\color{primary}\rule{0.45\textwidth}{1.2pt}}
\vspace{0.5em}
\end{center}

\section{Circular Queue Problems}

\subsection{Problem 12: Trace Circular Queue Operations}

\textbf{Problem Statement:} A circular queue has buffer capacity $N = 5$. Initially, $\texttt{front} = \texttt{rear} = -1$. Trace the following operations and determine if the final queue is full:
\[
\text{enq}(10), \text{enq}(20), \text{enq}(30), \text{enq}(40), \text{deq}(), \text{deq}(), \text{enq}(50), \text{enq}(60), \text{enq}(70)
\]

\subsubsection*{Step-by-Step Execution Trace}
Index advance formula: $\text{next} = (\text{index} + 1) \bmod N$. Full condition: $(\text{rear} + 1) \bmod N == \text{front}$.

\begin{table}[htbp]
\centering
\small
\begin{tabularx}{\textwidth}{l c c X}
\toprule
\textbf{Operation} & \textbf{Front} & \textbf{Rear} & \textbf{Logical Queue Content} \\
\midrule
Initial State & -1 & -1 & Queue is Empty \\
$\text{enqueue}(10)$ & 0 & 0 & $[10]$ \\
$\text{enqueue}(20)$ & 0 & 1 & $[10, 20]$ \\
$\text{enqueue}(30)$ & 0 & 2 & $[10, 20, 30]$ \\
$\text{enqueue}(40)$ & 0 & 3 & $[10, 20, 30, 40]$ \\
$\text{dequeue}()$   & 1 & 3 & $[20, 30, 40]$ ($10$ removed) \\
$\text{dequeue}()$   & 2 & 3 & $[30, 40]$ ($20$ removed) \\
$\text{enqueue}(50)$ & 2 & 4 & $[30, 40, 50]$ \\
$\text{enqueue}(60)$ & 2 & 0 & $[30, 40, 50, 60]$ (Wrapped to idx 0) \\
$\text{enqueue}(70)$ & 2 & 1 & $[30, 40, 50, 60, 70]$ (Wrapped to idx 1) \\
\bottomrule
\end{tabularx}
\caption{Trace of circular queue operations.}
\end{table}

\textbf{Final Fullness Evaluation:}
\[
(\text{rear} + 1) \bmod 5 = (1 + 1) \bmod 5 = 2 == \text{front}
\]
The queue contains 5 elements and is \textbf{strictly full}. Any further insertion triggers queue overflow.

\newpage
\begin{center}
\vspace{0.5em}
{\color{primary}\rule{0.45\textwidth}{1.2pt}}\\[4pt]
{\Large\bfseries\color{primary} Section VI: Tree Problems}\\[2pt]
{\color{primary}\rule{0.45\textwidth}{1.2pt}}
\vspace{0.5em}
\end{center}

\section{Tree Problems}

\subsection{Problem 13: Tree Traversal Orders}

\textbf{Problem Statement:} Consider the Binary Search Tree:
\begin{itemize}
    \item Root: $8$. Left child: $3$. Right child: $10$.
    \item Children of $3$: left $1$, right $6$. Children of $6$: left $4$, right $7$.
    \item Children of $10$: right $14$. Child of $14$: left $13$.
\end{itemize}
Determine the four standard traversals.

\subsubsection*{Solutions}
\begin{itemize}
    \item \textbf{Preorder (Root, Left, Right):}
    \[
    8, \; 3, \; 1, \; 6, \; 4, \; 7, \; 10, \; 14, \; 13
    \]
    \item \textbf{Inorder (Left, Root, Right):}
    \[
    \boxed{1, \; 3, \; 4, \; 6, \; 7, \; 8, \; 10, \; 13, \; 14} \quad (\text{strictly sorted in ascending order})
    \]
    \item \textbf{Postorder (Left, Right, Root):}
    \[
    1, \; 4, \; 7, \; 6, \; 3, \; 13, \; 14, \; 10, \; 8
    \]
    \item \textbf{Level-Order (Breadth-First):}
    \[
    8, \; 3, \; 10, \; 1, \; 6, \; 14, \; 4, \; 7, \; 13
    \]
\end{itemize}

\subsection{Problem 14: Leaves, Nodes, and Edges in a Full Binary Tree}

\textbf{Problem Statement:} A full binary tree has $I = 20$ internal nodes. Calculate leaves $L$, total nodes $N$, and total edges $E$.

\subsubsection*{Worked Calculation}
For any full binary tree (where every non-leaf node has exactly 2 children):
\[
L = I + 1 = 20 + 1 = \boxed{21 \text{ leaves}}
\]
\[
N = I + L = 20 + 21 = \boxed{41 \text{ total nodes}}
\]
\[
E = N - 1 = 41 - 1 = \boxed{40 \text{ edges}} \quad (\text{also verified by } 2I = 2 \times 20 = 40)
\]

\subsection{Problem 15: Minimum Height of a Binary Tree}

\textbf{Problem Statement:} Find the minimum possible height (in edges) of a binary tree containing $N = 100$ nodes.

\subsubsection*{Worked Calculation}
A binary tree of height $h$ edges has maximum node capacity:
\[
N_{\max} = 2^{h+1} - 1
\]
For $N = 100$:
\[
100 \leq 2^{h+1} - 1 \implies 101 \leq 2^{h+1}
\]
Since $2^6 = 64$ and $2^7 = 128$:
\[
h + 1 \geq 7 \implies \boxed{h_{\min} = 6 \text{ edges}}
\]

\subsection{Problem 16: Maximum Data Keys in a B+ Tree}

\textbf{Problem Statement:} A B+ tree has order $m = 4$ and height $h = 3$ edges. Each internal node can have at most $4$ children, and each leaf can hold at most $3$ data keys. Calculate maximum data-key capacity.

\subsubsection*{Worked Calculation}
The maximum number of leaf nodes at height $h = 3$ is:
\[
L_{\max} = m^h = 4^3 = 64 \text{ leaf nodes}
\]
Each leaf holds up to $m - 1 = 3$ data records. Total maximum data records:
\[
\boxed{N_{\max} = 64 \times 3 = 192 \text{ records}}
\]

\newpage
\begin{center}
\vspace{0.5em}
{\color{primary}\rule{0.45\textwidth}{1.2pt}}\\[4pt]
{\Large\bfseries\color{primary} Section VII: Graph Problems}\\[2pt]
{\color{primary}\rule{0.45\textwidth}{1.2pt}}
\vspace{0.5em}
\end{center}

\section{Graph Problems}

\subsection{Problem 17: Edge Calculation using Handshaking Lemma}

\textbf{Problem Statement:} An undirected graph has five vertices with degrees: $2, 3, 2, 3, 2$. Find the number of edges.

\subsubsection*{Worked Calculation}
By the Handshaking Lemma:
\[
\sum_{v \in V} \deg(v) = 2E \implies 2 + 3 + 2 + 3 + 2 = 12
\]
\[
2E = 12 \implies \boxed{E = 6 \text{ edges}}
\]

\subsection{Problem 18: Maximum Edges in Simple Graphs}

\textbf{Problem Statement:} Calculate maximum edges for $V = 10$ vertices:
\begin{enumerate}
    \item Simple undirected graph:
    \[
    E_{\max} = \frac{V(V-1)}{2} = \frac{10 \times 9}{2} = \boxed{45 \text{ edges}}
    \]
    \item Simple directed graph without self-loops:
    \[
    E_{\max} = V(V-1) = 10 \times 9 = \boxed{90 \text{ arcs}}
    \]
\end{enumerate}

\subsection{Problem 19: Complete Bipartite Graph Edges}

\textbf{Problem Statement:} Calculate total edges in complete bipartite graph $K_{4,6}$.
\[
E = a \times b = 4 \times 6 = \boxed{24 \text{ edges}}
\]

\subsection{Problem 20: BFS Traversal and Shortest Path}

\textbf{Problem Statement:} An undirected graph has adjacency lists:
\begin{lstlisting}[style=pseudocode]
A: B, C
B: A, D, E
C: A, F
D: B
E: B, F
F: C, E
\end{lstlisting}
Perform BFS starting from $A$ and find the shortest path from $A$ to $F$.

\subsubsection*{Execution Trace}
\begin{itemize}
    \item Step 1: Enqueue $A$, mark visited: $\text{Queue} = [A]$.
    \item Step 2: Dequeue $A$; enqueue unvisited neighbors $B, C$: $\text{Queue} = [B, C]$.
    \item Step 3: Dequeue $B$; enqueue unvisited neighbors $D, E$: $\text{Queue} = [C, D, E]$.
    \item Step 4: Dequeue $C$; enqueue unvisited neighbor $F$: $\text{Queue} = [D, E, F]$.
    \item Step 5: Dequeue $D, E, F$.
\end{itemize}
\textbf{BFS Visit Order:} $\boxed{A, B, C, D, E, F}$.\\
\textbf{Shortest Path to $F$:} $\boxed{A \to C \to F}$ (length: 2 edges).

\newpage
\begin{center}
\vspace{0.5em}
{\color{primary}\rule{0.45\textwidth}{1.2pt}}\\[4pt]
{\Large\bfseries\color{primary} Section VIII: Hash Table Problems}\\[2pt]
{\color{primary}\rule{0.45\textwidth}{1.2pt}}
\vspace{0.5em}
\end{center}

\section{Hash Table Problems}

\subsection{Problem 21: Division Method Hash Index Calculation}

\textbf{Problem Statement:} Table size $m = 11$, $h(k) = k \bmod 11$. Calculate slot indices for keys: $32, 44, 55, 67, 78$.

\begin{table}[htbp]
\centering
\small
\begin{tabular}{c c c l}
\toprule
\textbf{Key ($k$)} & \textbf{Calculation ($k \bmod 11$)} & \textbf{Slot Index} & \textbf{Collision Status} \\
\midrule
32 & $32 \bmod 11$ & 10 & No collision \\
44 & $44 \bmod 11$ & 0  & First key at slot 0 \\
55 & $55 \bmod 11$ & 0  & \textbf{Collides with 44 at slot 0} \\
67 & $67 \bmod 11$ & 1  & First key at slot 1 \\
78 & $78 \bmod 11$ & 1  & \textbf{Collides with 67 at slot 1} \\
\bottomrule
\end{tabular}
\caption{Hash indices and collisions under division hashing.}
\end{table}

\subsection{Problem 22: Linear Probing Insertion}

\textbf{Problem Statement:} Table size $m = 10$, $h(k) = k \bmod 10$. Insert keys in order: $23, 43, 13, 27, 33$ using linear probing.

\subsubsection*{Step-by-Step Insertion}
\begin{itemize}
    \item Key 23: $23 \bmod 10 = 3$. Slot 3 empty $\implies$ store \textbf{23 at index 3}.
    \item Key 43: $43 \bmod 10 = 3$. Slot 3 occupied $\implies$ probe slot 4 $\implies$ store \textbf{43 at index 4}.
    \item Key 13: $13 \bmod 10 = 3$. Slots 3, 4 occupied $\implies$ probe slot 5 $\implies$ store \textbf{13 at index 5}.
    \item Key 27: $27 \bmod 10 = 7$. Slot 7 empty $\implies$ store \textbf{27 at index 7}.
    \item Key 33: $33 \bmod 10 = 3$. Slots 3, 4, 5 occupied $\implies$ probe slot 6 $\implies$ store \textbf{33 at index 6}.
\end{itemize}

\begin{table}[htbp]
\centering
\small
\begin{tabular}{c c c c c c c c c c c}
\toprule
\textbf{Index} & 0 & 1 & 2 & 3 & 4 & 5 & 6 & 7 & 8 & 9 \\
\midrule
\textbf{Content} & --- & --- & --- & \textbf{23} & \textbf{43} & \textbf{13} & \textbf{33} & \textbf{27} & --- & --- \\
\bottomrule
\end{tabular}
\caption{Final state of hash table after linear probing.}
\end{table}

\subsection{Problem 23: Load Factor Calculation}

\textbf{Problem Statement:} Table has $m = 50$ slots and contains $n = 35$ entries.
\[
\alpha = \frac{n}{m} = \frac{35}{50} = 0.70 \implies \boxed{70\% \text{ occupancy}}
\]

\subsection{Problem 24: Rehashing Threshold Calculation}

\textbf{Problem Statement:} Table capacity $m = 200$. Rehashing triggers when $\alpha > 0.75$. What is the maximum number of items stored before rehashing occurs?
\[
\frac{n}{200} \leq 0.75 \implies n \leq 200 \times 0.75 \implies \boxed{n_{\max} = 150 \text{ entries}}
\]
Inserting the $151^{\text{st}}$ item raises $\alpha = 151/200 = 0.755 > 0.75$, triggering table doubling and rehashing.

\newpage
\begin{center}
\vspace{0.5em}
{\color{primary}\rule{0.45\textwidth}{1.2pt}}\\[4pt]
{\Large\bfseries\color{primary} Section IX: Representation and Architecture Selection}\\[2pt]
{\color{primary}\rule{0.45\textwidth}{1.2pt}}
\vspace{0.5em}
\end{center}

\section{Graph Representation and Architecture Selection}

\subsection{Problem 25: Memory Comparison of Matrix vs List}

\textbf{Problem Statement:} An undirected graph has $V = 6$ vertices and $E = 7$ edges. Compare raw structural cell counts:
\begin{itemize}
    \item \textbf{Adjacency Matrix:} $V^2 = 6^2 = \boxed{36 \text{ entries}}$.
    \item \textbf{Adjacency List:} $V$ head pointers $+ 2E$ edge nodes $= 6 + 14 = \boxed{20 \text{ structural entries}}$.
\end{itemize}

\subsection{Problem 26: Master Data Structure Selection Guide}

\begin{table}[htbp]
\centering
\small
\begin{tabularx}{\textwidth}{p{7.5cm} l X}
\toprule
\textbf{Engineering Requirement} & \textbf{Optimal Structure} & \textbf{Primary Theoretical Justification} \\
\midrule
Undo operation in text editors & Stack & Strict LIFO retrieval \\
Print job scheduling / CPU task queue & Queue & Fair FIFO order processing \\
Fixed-size streaming packet buffer & Circular Queue & $O(1)$ slot recycling without shifting \\
Extract highest-urgency emergency task & Max-Priority Queue & Highest-key service first \\
Fast insertion/deletion without contiguous memory & Linked List & Dynamic pointer rearrangement \\
Retrieve maximum key repeatedly & Max-Heap & $O(1)$ root peek, $O(\log n)$ extract \\
Organizational hierarchy / File directory & Tree & Parent-child non-cyclic branching \\
Database indexing with range queries & B+ Tree & Massive disk fanout; linked leaves \\
Dictionary with expected $O(1)$ search & Hash Table & Direct key-to-slot mapping \\
Graph traversal on sparse networks & Adjacency List & $O(V+E)$ space; optimal neighbor loop \\
Repeated edge-existence test in dense graph & Adjacency Matrix & $O(1)$ direct array index lookup \\
Shortest path in unweighted graph & BFS (Queue) & Explores vertices by edge-hop levels \\
Cycle detection / Topological sorting & DFS (Stack) & Explores deepest paths before backtrack \\
Prefix matching / Auto-complete & Trie & Reuses shared string prefixes \\
\bottomrule
\end{tabularx}
\caption{Master data structure selection chart for technical examinations.}
\end{table}

\newpage
\begin{center}
\vspace{0.5em}
{\color{primary}\rule{0.45\textwidth}{1.2pt}}\\[4pt]
{\Large\bfseries\color{primary} Section X: Rapid-Fire Revision Questions}\\[2pt]
{\color{primary}\rule{0.45\textwidth}{1.2pt}}
\vspace{0.5em}
\end{center}

\section{Rapid-Fire Revision Questions}

\begin{longtable}{r p{9.8cm} p{6.8cm}}
\toprule
\textbf{\#} & \textbf{Recruitment Examination Question} & \textbf{Concise Technical Answer} \\
\midrule
\endfirsthead
\toprule
\textbf{\#} & \textbf{Recruitment Examination Question} & \textbf{Concise Technical Answer} \\
\midrule
\endhead
\bottomrule
\endfoot
\bottomrule
\endlastfoot
1  & Time complexity of finding minimum in unsorted array? & $O(n)$ \\
2  & Time complexity of iterative binary search? & $O(\log n)$ \\
3  & Which principle does a stack follow? & LIFO (Last In, First Out) \\
4  & Which principle does a queue follow? & FIFO (First In, First Out) \\
5  & Purpose of two stacks in queue implementation? & Order reversal to achieve FIFO semantics \\
6  & Time complexity of iterative singly linked list reversal? & $O(n)$ time, $O(1)$ auxiliary space \\
7  & Algorithm detecting linked list cycle with $O(1)$ space? & Floyd's Cycle-Finding (Tortoise and Hare) \\
8  & Time complexity of bottom-up heap construction? & $O(n)$ \\
9  & Location of maximum element in a max-heap? & At the root ($A[0]$) \\
10 & Worst-case insertion complexity in binary heap? & $O(\log n)$ \\
11 & What does a circular queue reuse? & Vacated front array slots via modulo \\
12 & What determines removal order in priority queue? & Priority key value \\
13 & Number of edges in a tree with $n$ vertices? & $n - 1$ \\
14 & Maximum nodes in binary tree of height $h$ edges? & $2^{h+1} - 1$ \\
15 & Relationship between leaves and internal nodes in full binary tree? & $L = I + 1$ \\
16 & Height of skewed binary tree with $n$ nodes? & $n - 1$ edges \\
17 & Worst-case search complexity of unbalanced BST? & $O(n)$ \\
18 & Which traversal outputs keys of BST in sorted order? & Inorder traversal \\
19 & AVL tree balance factor restriction? & $|BF| \leq 1$ (i.e., $-1, 0, +1$) \\
20 & Maximum children in order-$m$ B-tree? & $m$ \\
21 & Where are actual data records stored in B+ tree? & Exclusively in leaf nodes \\
22 & Maximum B+ tree leaf capacity at height $h$, order $m$? & $(m-1)m^h$ \\
23 & Maximum edges in simple undirected graph with $V$ vertices? & $\frac{V(V-1)}{2}$ \\
24 & Sum of degrees in an undirected graph? & $2E$ (Handshaking Lemma) \\
25 & Maximum edges in complete bipartite graph $K_{a,b}$? & $a \times b$ \\
26 & Graph traversal algorithm implemented using Queue? & BFS (Breadth-First Search) \\
27 & Graph traversal algorithm implemented using Stack/Recursion? & DFS (Depth-First Search) \\
28 & Expected search complexity in a well-distributed hash table? & $O(1)$ \\
29 & Worst-case search complexity in a hash table? & $O(n)$ (all keys collide) \\
30 & Mathematical definition of hash table load factor $\alpha$? & $\alpha = n / m$ \\
31 & Collision resolution method prone to primary clustering? & Linear probing \\
32 & Purpose of tombstones during open addressing deletion? & Preserving unbroken probe search chains \\
33 & Space complexity of graph adjacency matrix? & $O(V^2)$ \\
34 & Space complexity of graph adjacency list? & $O(V + E)$ \\
35 & Graph representation providing instantaneous $O(1)$ edge test? & Adjacency matrix \\
36 & Representation most memory-efficient for sparse graphs? & Adjacency list \\
\end{longtable}
